\input{preamble}

% OK, start here.
%
\begin{document}

\title{Schémas en groupoïdes}


\maketitle

\phantomsection
\label{section-phantom}

\tableofcontents

\section{Introduction}
\label{section-introduction}

\noindent
Ce chapitre est consacré aux généralités sur les schémas en groupoïdes.
Voir par exemple le bel article \cite{K-M} de Keel et Mori.





\section{Notations}
\label{section-notation}

\noindent
Soit $S$ un schéma. Si $U$ et $T$ sont des schémas sur $S$, nous notons
$U(T)$ l'ensemble des points à valeurs dans $T$ du schéma $U$ {\it au-dessus de} $S$. En formule :
$U(T) = \Mor_S(T, U)$. Nous nous efforçons de réserver la lettre $T$ pour désigner
un ``schéma test'' sur $S$, comme dans la discussion qui suit.
Soient $X$ et $Y$ des schémas sur
$S$, et soit $f : X \to Y$ un morphisme de schémas sur $S$.
Pour tout schéma $T$ sur $S$, on obtient une application induite entre ensembles
$$
f : X(T) \longrightarrow Y(T)
$$
que, comme indiqué, nous notons aussi $f$. En fait, cette construction
est fonctorielle en $T/S$. Le lemme de Yoneda, voir Catégories,
lemme \ref{categories-lemma-yoneda}, affirme que $f$ détermine et est
déterminé par cette transformation de foncteurs $f : h_X \to h_Y$.
Plus généralement, nous employons la même notation pour les applications entre produits
fibrés. Par exemple, si
$X$, $Y$ et $Z$ sont des schémas sur $S$, et si
$m : X \times_S Y \to Z \times_S Z$ est
un morphisme de schémas sur $S$, alors nous considérons $m$ comme correspondant
à une famille d'applications entre les ensembles de points à valeurs dans $T$
$$
X(T) \times Y(T) \longrightarrow Z(T) \times Z(T).
$$
Et ainsi de suite.

\medskip\noindent
Nous conservons notre convention consistant à indexer les morphismes de projection à partir de
l'indice $0$; on a donc $\text{pr}_0 : X \times_S Y \to X$ et
$\text{pr}_1 : X \times_S Y \to Y$.






\section{Relations d'équivalence}
\label{section-equivalence-relations}

\noindent
Rappelons qu'une {\it relation} $R$ sur un ensemble $A$ est simplement un sous-ensemble
$R \subset A \times A$. On écrit habituellement $a R b$ pour indiquer que
$(a, b) \in R$. On dit que la relation est {\it transitive} si
$a R b, b R c \Rightarrow a R c$. On dit que la relation est
{\it réflexive} si $a R a$ pour tout $a \in A$. On dit que la relation est
{\it symétrique} si $a R b \Rightarrow b R a$.
Une relation est appelée {\it relation d'équivalence} si
elle est transitive, réflexive et symétrique.

\medskip\noindent
Dans le cadre des schémas, nous allons assouplir quelque peu la notion de
relation et exiger seulement que $R \to A \times A$ soit
un morphisme. Voici la définition.

\begin{definition}
\label{definition-equivalence-relation}
Soit $S$ un schéma. Soit $U$ un schéma sur $S$.
\begin{enumerate}
\item Une {\it prérelation} sur $U$ relativement à $S$ est un morphisme quelconque
de schémas $j : R \to U \times_S U$. Dans ce cas, on pose
$t = \text{pr}_0 \circ j$ et $s = \text{pr}_1 \circ j$, de sorte
que $j = (t, s)$.
\item Une {\it relation} sur $U$ relativement à $S$ est un monomorphisme
de schémas $j : R \to U \times_S U$.
\item Une {\it prérelation d'équivalence} est une prérelation
$j : R \to U \times_S U$ telle que l'image de
$j : R(T) \to U(T) \times U(T)$ est une relation d'équivalence pour
tout $T/S$.
\item On dit qu'un morphisme de schémas $R \to U \times_S U$ est
une {\it relation d'équivalence sur $U$ relativement à $S$}
si et seulement si, pour tout schéma $T$ sur $S$, les points à valeurs dans $T$
du schéma $R$ définissent une relation d'équivalence
sur l'ensemble des points à valeurs dans $T$ du schéma $U$.
\end{enumerate}
\end{definition}

\noindent
Autrement dit, une relation d'équivalence est une prérelation d'équivalence
telle que $j$ soit une relation.

\begin{lemma}
\label{lemma-restrict-relation}
Soit $S$ un schéma.
Soit $U$ un schéma sur $S$.
Soit $j : R \to U \times_S U$ une prérelation.
Soit $g : U' \to U$ un morphisme de schémas.
Enfin, posons
$$
R' = (U' \times_S U')\times_{U \times_S U} R
\xrightarrow{j'}
U' \times_S U'
$$
Alors $j'$ est une prérelation sur $U'$ relativement à $S$.
Si $j$ est une relation, alors $j'$ est une relation.
Si $j$ est une prérelation d'équivalence, alors $j'$ est une prérelation d'équivalence.
Si $j$ est une relation d'équivalence, alors $j'$ est une relation d'équivalence.
\end{lemma}

\begin{proof}
Démonstration omise.
\end{proof}

\begin{definition}
\label{definition-restrict-relation}
Soit $S$ un schéma.
Soit $U$ un schéma sur $S$.
Soit $j : R \to U \times_S U$ une prérelation.
Soit $g : U' \to U$ un morphisme de schémas.
La prérelation $j' : R' \to U' \times_S U'$ est appelée
la {\it restriction}, ou l'{\it image réciproque}, de la prérelation $j$ sur $U'$.
Dans cette situation, nous écrivons parfois $R' = R|_{U'}$.
\end{definition}

\begin{lemma}
\label{lemma-pre-equivalence-equivalence-relation-points}
Soit $j : R \to U \times_S U$ une prérelation.
Considérons la relation sur les points du schéma $U$ définie par
la règle
$$
x \sim y
\Leftrightarrow
\exists\ r \in R :
t(r) = x,
s(r) = y.
$$
Si $j$ est une prérelation d'équivalence, alors la relation ainsi définie est une
relation d'équivalence.
\end{lemma}

\begin{proof}
Supposons que $x \sim y$ et que $y \sim z$.
Choisissons $r \in R$ tel que $t(r) = x$, $s(r) = y$ et
choisissons $r' \in R$ tel que $t(r') = y$, $s(r') = z$.
Choisissons un corps $K$ s'insérant dans le diagramme commutatif
suivant
$$
\xymatrix{
\kappa(r) \ar[r] & K \\
\kappa(y) \ar[u] \ar[r] & \kappa(r') \ar[u]
}
$$
Désignons par $x_K, y_K, z_K : \Spec(K) \to U$
les morphismes
$$
\begin{matrix}
\Spec(K) \to \Spec(\kappa(r))
\to
\Spec(\kappa(x)) \to U \\
\Spec(K) \to \Spec(\kappa(r))
\to
\Spec(\kappa(y)) \to U \\
\Spec(K) \to \Spec(\kappa(r'))
\to
\Spec(\kappa(z)) \to U
\end{matrix}
$$
Par construction, $(x_K, y_K) \in j(R(K))$ et
$(y_K, z_K) \in j(R(K))$. Comme $j$ est une prérelation d'équivalence,
on a aussi $(x_K, z_K) \in j(R(K))$.
Cela implique clairement que $x \sim z$.

\medskip\noindent
La démonstration du fait que $\sim$ est réflexive et symétrique est omise.
\end{proof}

\begin{lemma}
\label{lemma-etale-equivalence-relation}
Soit $j : R \to U \times_S U$ une prérelation. Supposons que
\begin{enumerate}
\item $s, t$ soient non ramifiés,
\item pour tout corps algébriquement clos $k$ sur $S$
l'application $R(k) \to U(k) \times U(k)$ définisse une relation d'équivalence,
\item il existe des morphismes $e : U \to R$, $i : R \to R$,
$c : R \times_{s, U, t} R \to R$ tels que
$$
\xymatrix@C-1pc{
U \ar[r]_e \ar[d]_\Delta &
R \ar[d]_j &
R \ar[d]^j \ar[r]_i &
R \ar[d]^j &
R \times_{s, U, t} R \ar[d]^{j \times j} \ar[r]_c &
R \ar[d]^j \\
U \times_S U \ar[r] &
U \times_S U &
U \times_S U \ar[r]^{flip} &
U \times_S U &
U \times_S U \times_S U \ar[r]^{\text{pr}_{02}} &
U \times_S U
}
$$
sont commutatifs.
\end{enumerate}
Alors $j$ est une relation d'équivalence.
\end{lemma}

\begin{proof}
D'après la condition (1) et le
lemme \ref{morphisms-lemma-unramified-permanence} de Morphismes,
on voit que $j$ est non ramifié. Ainsi,
$\Delta_j : R \to R \times_{U \times_S U} R$ est une immersion ouverte d'après
le lemme \ref{morphisms-lemma-diagonal-unramified-morphism} de Morphismes.
Or la condition (2) affirme que $\Delta_j$ est bijective sur les
points à valeurs dans $k$, donc $\Delta_j$ est un isomorphisme, donc $j$
est un monomorphisme. Il découle alors aisément des diagrammes
commutatifs que $R(T) \subset U(T) \times U(T)$ est une relation
d'équivalence pour tout schéma $T$ sur $S$.
\end{proof}













\section{Schémas en groupes}
\label{section-group-schemes}

\noindent
Rappelons qu'un {\it groupe} est un couple
$(G, m)$, où $G$ est un ensemble et $m : G \times G \to G$ est
une application d'ensembles possédant les propriétés suivantes :
\begin{enumerate}
\item (associativité) $m(g, m(g', g'')) = m(m(g, g'), g'')$
pour tous $g, g', g'' \in G$,
\item (élément neutre) il existe un unique élément $e \in G$
(appelé {\it élément neutre}, {\it unité}, ou $1$ de $G$) tel que
$m(g, e) = m(e, g) = g$ pour tout $g \in G$, et
\item (inverse) pour tout $g \in G$ il existe un élément $i(g) \in G$
tel que $m(g, i(g)) = m(i(g), g) = e$, où $e$ est
l'élément neutre.
\end{enumerate}
Nous obtenons ainsi une application $e : \{*\} \to G$ et une application
$i : G \to G$ de sorte que le quadruplet $(G, m, e, i)$
satisfait les axiomes énumérés ci-dessus.

\medskip\noindent
Un {\it homomorphisme de groupes} $\psi : (G, m) \to (G', m')$
est une application d'ensembles $\psi : G \to G'$ telle que
$m'(\psi(g), \psi(g')) = \psi(m(g, g'))$. Il s'ensuit automatiquement
que $\psi(e) = e'$ et $i'(\psi(g)) = \psi(i(g))$.
(Notations évidentes.) Nous nous en servirons ci-dessous.

\begin{definition}
\label{definition-group-scheme}
Soit $S$ un schéma.
\begin{enumerate}
\item Un {\it schéma en groupes sur $S$} est un couple $(G, m)$, où
$G$ est un schéma sur $S$ et $m : G \times_S G \to G$ est
un morphisme de schémas sur $S$ vérifiant la propriété suivante :
Pour tout schéma $T$ sur $S$, le couple $(G(T), m)$
est un groupe.
\item Un {\it morphisme $\psi : (G, m) \to (G', m')$ de schémas en groupes sur $S$}
est un morphisme $\psi : G \to G'$ de schémas sur $S$ tel que, pour
tout $T/S$, l'application induite $\psi : G(T) \to G'(T)$ soit un homomorphisme
de groupes.
\end{enumerate}
\end{definition}

\noindent
Soit $(G, m)$ un schéma en groupes sur le schéma $S$.
D'après la discussion ci-dessus (et celle de la section \ref{section-notation})
nous obtenons des morphismes de schémas sur $S$ :
(élément neutre) $e : S \to G$ et (inverse) $i : G \to G$ tels que
pour tout $T$, le quadruplet $(G(T), m, e, i)$ satisfait les
axiomes d'un groupe énumérés ci-dessus.

\medskip\noindent
Soient $(G, m)$ et $(G', m')$ des schémas en groupes sur $S$.
Soit $f : G \to G'$ un morphisme de schémas sur $S$.
Il résulte de la définition que $f$ est un morphisme
de schémas en groupes sur $S$ si et seulement si le diagramme suivant
est commutatif :
$$
\xymatrix{
G \times_S G \ar[r]_-{f \times f} \ar[d]_m &
G' \times_S G' \ar[d]^{m'} \\
G \ar[r]^f & G'
}
$$

\begin{lemma}
\label{lemma-base-change-group-scheme}
Soit $(G, m)$ un schéma en groupes sur $S$.
Soit $S' \to S$ un morphisme de schémas.
Le changement de base $(G_{S'}, m_{S'})$ est un schéma en groupes sur $S'$.
\end{lemma}

\begin{proof}
La démonstration est omise.
\end{proof}

\begin{definition}
\label{definition-closed-subgroup-scheme}
Soit $S$ un schéma. Soit $(G, m)$ un schéma en groupes sur $S$.
\begin{enumerate}
\item Un {\it sous-schéma en groupes fermé} de $G$ est un sous-schéma fermé
$H \subset G$ tel que $m|_{H \times_S H}$ se factorise par $H$ et induise
sur $H$ une structure de schéma en groupes sur $S$.
\item Un {\it sous-schéma en groupes ouvert} de $G$ est un sous-schéma ouvert
$G' \subset G$ tel que $m|_{G' \times_S G'}$ se factorise par $G'$
et induise sur $G'$ une structure de schéma en groupes sur $S$.
\end{enumerate}
\end{definition}

\noindent
On pourrait également dire que $H$ est un sous-schéma en groupes fermé de $G$
si c'est un schéma en groupes sur $S$ muni d'un morphisme de schémas en groupes
$i : H \to G$ sur $S$ qui identifie $H$ à un sous-schéma fermé de $G$.

\begin{lemma}
\label{lemma-closed-subgroup-scheme}
Soit $S$ un schéma. Soit $(G, m, e, i)$ un schéma en groupes sur $S$.
\begin{enumerate}
\item Un sous-schéma fermé $H \subset G$ est un sous-schéma en groupes fermé
si et seulement si $e : S \to G$, $m|_{H \times_S H} : H \times_S H \to G$,
et $i|_H : H \to G$ se factorisent par $H$.
\item Un sous-schéma ouvert $H \subset G$ est un sous-schéma en groupes ouvert
si et seulement si $e : S \to G$, $m|_{H \times_S H} : H \times_S H \to G$,
et $i|_H : H \to G$ se factorisent par $H$.
\end{enumerate}
\end{lemma}

\begin{proof}
En considérant les points à valeurs dans $T$, cela se traduit par les conditions bien connues
caractérisant les parties d'un groupe qui en sont des sous-groupes.
\end{proof}

\begin{definition}
\label{definition-smooth-group-scheme}
Soit $S$ un schéma. Soit $(G, m)$ un schéma en groupes sur $S$.
\begin{enumerate}
\item On dit que $G$ est un {\it schéma en groupes lisse} si son morphisme
structural $G \to S$ est lisse.
\item On dit que $G$ est un {\it schéma en groupes plat} si son morphisme
structural $G \to S$ est plat.
\item On dit que $G$ est un {\it schéma en groupes séparé} si son morphisme
structural $G \to S$ est séparé.
\end{enumerate}
On pourra compléter cette liste au besoin.
\end{definition}






\section{Exemples de schémas en groupes}
\label{section-examples-group-schemes}

\begin{example}[Schéma en groupes multiplicatif]
\label{example-multiplicative-group}
Considérons le foncteur qui associe
à tout schéma $T$ le groupe $\Gamma(T, \mathcal{O}_T^*)$
des unités de l'anneau des sections globales du faisceau structural.
Ce foncteur est représentable par le schéma
$$
\mathbf{G}_m = \Spec(\mathbf{Z}[x, x^{-1}])
$$
Le morphisme définissant la structure de groupe est le morphisme
\begin{eqnarray*}
\mathbf{G}_m \times \mathbf{G}_m & \to & \mathbf{G}_m \\
\Spec(\mathbf{Z}[x, x^{-1}] \otimes_{\mathbf{Z}} \mathbf{Z}[x, x^{-1}])
& \to &
\Spec(\mathbf{Z}[x, x^{-1}]) \\
\mathbf{Z}[x, x^{-1}] \otimes_{\mathbf{Z}} \mathbf{Z}[x, x^{-1}]
& \leftarrow &
\mathbf{Z}[x, x^{-1}] \\
x \otimes x & \leftarrow & x
\end{eqnarray*}
On voit donc que $\mathbf{G}_m$ est un schéma en groupes sur $\mathbf{Z}$.
Pour tout schéma $S$, le changement de base $\mathbf{G}_{m, S}$ est un
schéma en groupes sur $S$ dont le foncteur des points est
$$
T/S
\longmapsto
\mathbf{G}_{m, S}(T) = \mathbf{G}_m(T) = \Gamma(T, \mathcal{O}_T^*)
$$
comme précédemment.
\end{example}

\begin{example}[Racines de l'unité]
\label{example-roots-of-unity}
Soit $n \in \mathbf{N}$.
Considérons le foncteur qui associe
à tout schéma $T$ le sous-groupe de $\Gamma(T, \mathcal{O}_T^*)$
formé des racines $n$-ièmes de l'unité.
Ce foncteur est représentable par le schéma
$$
\mu_n = \Spec(\mathbf{Z}[x]/(x^n - 1)).
$$
Le morphisme définissant la structure de groupe est le morphisme
\begin{eqnarray*}
\mu_n \times \mu_n & \to & \mu_n \\
\Spec(
\mathbf{Z}[x]/(x^n - 1)
\otimes_{\mathbf{Z}}
\mathbf{Z}[x]/(x^n - 1))
& \to &
\Spec(\mathbf{Z}[x]/(x^n - 1)) \\
\mathbf{Z}[x]/(x^n - 1) \otimes_{\mathbf{Z}} \mathbf{Z}[x]/(x^n - 1)
& \leftarrow &
\mathbf{Z}[x]/(x^n - 1) \\
x \otimes x & \leftarrow & x
\end{eqnarray*}
On voit donc que $\mu_n$ est un schéma en groupes sur $\mathbf{Z}$.
Pour tout schéma $S$, le changement de base $\mu_{n, S}$ est un
schéma en groupes sur $S$ dont le foncteur des points est
$$
T/S
\longmapsto
\mu_{n, S}(T) = \mu_n(T) = \{f \in \Gamma(T, \mathcal{O}_T^*) \mid f^n = 1\}
$$
comme précédemment.
\end{example}


\begin{example}[Schéma en groupes additif]
\label{example-additive-group}
Considérons le foncteur qui associe
à tout schéma $T$ le groupe $\Gamma(T, \mathcal{O}_T)$
des sections globales du faisceau structural.
Ce foncteur est représentable par le schéma
$$
\mathbf{G}_a = \Spec(\mathbf{Z}[x])
$$
Le morphisme définissant la structure de groupe est le morphisme
\begin{eqnarray*}
\mathbf{G}_a \times \mathbf{G}_a & \to & \mathbf{G}_a \\
\Spec(\mathbf{Z}[x] \otimes_{\mathbf{Z}} \mathbf{Z}[x])
& \to &
\Spec(\mathbf{Z}[x]) \\
\mathbf{Z}[x] \otimes_{\mathbf{Z}} \mathbf{Z}[x]
& \leftarrow &
\mathbf{Z}[x] \\
x \otimes 1 + 1 \otimes x & \leftarrow & x
\end{eqnarray*}
On voit donc que $\mathbf{G}_a$ est un schéma en groupes sur $\mathbf{Z}$.
Pour tout schéma $S$, le changement de base $\mathbf{G}_{a, S}$ est un
schéma en groupes sur $S$ dont le foncteur des points est
$$
T/S
\longmapsto
\mathbf{G}_{a, S}(T) = \mathbf{G}_a(T) = \Gamma(T, \mathcal{O}_T)
$$
comme précédemment.
\end{example}

\begin{example}[Schéma en groupes linéaire général]
\label{example-general-linear-group}
Soit $n \geq 1$.
Considérons le foncteur qui associe
à tout schéma $T$ le groupe
$$
\text{GL}_n(\Gamma(T, \mathcal{O}_T))
$$
des matrices inversibles $n \times n$ à coefficients dans
l'anneau des sections globales du faisceau structural.
Ce foncteur est représentable par le schéma
$$
\text{GL}_n = \Spec(\mathbf{Z}[\{x_{ij}\}_{1 \leq i, j \leq n}][1/d])
$$
où $d = \det((x_{ij}))$ et où $(x_{ij})$ est la matrice $n \times n$
ayant pour coefficient $x_{ij}$ en position $(i, j)$.
Le morphisme définissant la structure de groupe est le morphisme
\begin{eqnarray*}
\text{GL}_n \times \text{GL}_n & \to & \text{GL}_n \\
\Spec(\mathbf{Z}[x_{ij}, 1/d] \otimes_{\mathbf{Z}}
\mathbf{Z}[x_{ij}, 1/d])
& \to &
\Spec(\mathbf{Z}[x_{ij}, 1/d]) \\
\mathbf{Z}[x_{ij}, 1/d] \otimes_{\mathbf{Z}} \mathbf{Z}[x_{ij}, 1/d]
& \leftarrow &
\mathbf{Z}[x_{ij}, 1/d] \\
\sum x_{ik} \otimes x_{kj} & \leftarrow & x_{ij}
\end{eqnarray*}
Ainsi, on voit que $\text{GL}_n$ est un schéma en groupes sur $\mathbf{Z}$.
Pour tout schéma $S$, le changement de base $\text{GL}_{n, S}$ est un
schéma en groupes sur $S$ dont le foncteur des points est donné par
$$
T/S
\longmapsto
\text{GL}_{n, S}(T) = \text{GL}_n(T) =\text{GL}_n(\Gamma(T, \mathcal{O}_T))
$$
comme précédemment.
\end{example}

\begin{example}
\label{example-determinant}
Le déterminant définit un morphisme de schémas en groupes
$$
\det : \text{GL}_n \longrightarrow \mathbf{G}_m
$$
sur $\mathbf{Z}$. Par changement de base, il donne un morphisme
de schémas en groupes $\text{GL}_{n, S} \to \mathbf{G}_{m, S}$
sur tout schéma de base $S$.
\end{example}

\begin{example}[Groupe constant]
\label{example-constant-group}
Soit $G$ un groupe abstrait. Considérons le foncteur
qui associe à tout schéma $T$ le groupe
des applications localement constantes $T \to G$ (où $T$ est muni de la topologie de Zariski
et $G$ de la topologie discrète). Ce foncteur est représentable par le schéma
$$
G_{\Spec(\mathbf{Z})} =
\coprod\nolimits_{g \in G} \Spec(\mathbf{Z}).
$$
Le morphisme définissant la structure de groupe est le morphisme
$$
G_{\Spec(\mathbf{Z})}
\times_{\Spec(\mathbf{Z})}
G_{\Spec(\mathbf{Z})}
\longrightarrow
G_{\Spec(\mathbf{Z})}
$$
qui envoie la composante correspondant au couple $(g, g')$ sur la
composante correspondant à $gg'$. Pour tout schéma $S$, le changement de base
$G_S$ est un schéma en groupes sur $S$ dont le foncteur des points est donné par
$$
T/S
\longmapsto
G_S(T) = \{f : T \to G \text{ localement constante}\}
$$
comme précédemment.
\end{example}





\section{Propriétés des schémas en groupes}
\label{section-properties-group-schemes}

\noindent
Dans cette section, nous rassemblons quelques propriétés élémentaires des schémas en groupes qui
sont valables sur toute base.

\begin{lemma}
\label{lemma-group-scheme-separated}
Soit $S$ un schéma.
Soit $G$ un schéma en groupes sur $S$.
Alors $G \to S$ est séparé (resp.\ quasi-séparé) si et seulement si
le morphisme unité $e : S \to G$ est une immersion fermée
(resp.\ quasi-compact).
\end{lemma}

\begin{proof}
Rappelons que, d'après le
lemme \ref{schemes-lemma-section-immersion} de Schémas,
$e$ est une immersion, laquelle est fermée
(resp.\ quasi-compacte) si $G \to S$ est séparé (resp.\ quasi-séparé).
Pour la réciproque, considérons le diagramme
$$
\xymatrix{
G \ar[r]_-{\Delta_{G/S}} \ar[d] &
G \times_S G \ar[d]^{(g, g') \mapsto m(i(g), g')} \\
S \ar[r]^e & G
}
$$
Montrer que ce diagramme est cartésien est un exercice d'application du point de vue
fonctoriel en géométrie algébrique. Autrement dit, on voit que
$\Delta_{G/S}$ s'obtient à partir de $e$ par changement de base. Ainsi, si $e$ est une
immersion fermée (resp.\ quasi-compact), il en est de même de $\Delta_{G/S}$, voir
Schémas, lemme \ref{schemes-lemma-base-change-immersion}
(resp.\ Schémas, Lemme
\ref{schemes-lemma-quasi-compact-preserved-base-change}).
\end{proof}

\begin{lemma}
\label{lemma-flat-action-on-group-scheme}
Soit $S$ un schéma.
Soit $G$ un schéma en groupes sur $S$.
Soit $T$ un schéma sur $S$ et soit $\psi : T \to G$ un morphisme sur $S$.
Si $T$ est plat sur $S$, alors le morphisme
$$
T \times_S G \longrightarrow G, \quad
(t, g) \longmapsto m(\psi(t), g)
$$
est plat. En particulier, si $G$ est plat sur $S$, alors
$m : G \times_S G \to G$ est plat.
\end{lemma}

\begin{proof}
Considérons le diagramme
$$
\xymatrix{
T \times_S G \ar[rrr]_{(t, g) \mapsto (t, m(\psi(t), g))} & & &
T \times_S G \ar[r]_{\text{pr}} \ar[d] &
G \ar[d] \\
& & &
T \ar[r] &
S
}
$$
La flèche horizontale supérieure gauche est un isomorphisme et le
carré est cartésien. Le lemme résulte donc de
Morphismes, Lemme \ref{morphisms-lemma-base-change-flat}.
\end{proof}

\begin{lemma}
\label{lemma-group-scheme-module-differentials}
\begin{reference}
\cite[Proposition 3.15]{BookAV}
\end{reference}
Soit $(G, m, e, i)$ un schéma en groupes sur le schéma $S$.
Notons $f : G \to S$ le morphisme structural.
Il existe alors des isomorphismes canoniques
$$
\Omega_{G/S} \cong f^*\mathcal{C}_{S/G} \cong f^*e^*\Omega_{G/S}
$$
où $\mathcal{C}_{S/G}$ désigne le faisceau conormal de
l'immersion $e$. En particulier, si $S$ est le spectre d'un corps, alors
$\Omega_{G/S}$ est un $\mathcal{O}_G$-module libre.
\end{lemma}

\begin{proof}
D'après Morphismes, Lemme \ref{morphisms-lemma-base-change-differentials}, nous avons
$$
\Omega_{G \times_S G/G} = \text{pr}_0^*\Omega_{G/S}
$$
où, dans le membre de gauche, nous considérons $G \times_S G$ comme un schéma sur $G$
au moyen de $\text{pr}_1$.
Soit $\tau : G \times_S G \to G \times_S G$ l'application dite ``de cisaillement''
donnée sur les points par $(g, h) \mapsto (m(g, h), h)$. Cette application est un automorphisme
de $G \times_S G$ considéré comme un schéma sur $G$ via la projection $\text{pr}_1$.
En combinant ces deux remarques, nous obtenons un isomorphisme
$$
\tau^*\text{pr}_0^*\Omega_{G/S} \to \text{pr}_0^*\Omega_{G/S}
$$
Puisque $\text{pr}_0 \circ \tau = m$, celui-ci peut se réécrire sous la forme
$$
m^*\Omega_{G/S} \to \text{pr}_0^*\Omega_{G/S}
$$
En tirant cet isomorphisme en arrière par
$(e \circ f, \text{id}_G) : G \to G \times_S G$
et en utilisant que $m \circ (e \circ f, \text{id}_G) = \text{id}_G$
et que $\text{pr}_0 \circ (e \circ f, \text{id}_G) = e \circ f$,
nous obtenons un isomorphisme
$$
\Omega_{G/S} \to f^*e^*\Omega_{G/S}
$$
comme souhaité. D'après
Morphismes, Lemme \ref{morphisms-lemma-differentials-relative-immersion-section},
nous avons $\mathcal{C}_{S/G} \cong e^*\Omega_{G/S}$.
Si $S$ est le spectre d'un corps, alors
tout $\mathcal{O}_S$-module sur $S$ est libre,
d'où la dernière assertion.
\end{proof}

\begin{lemma}
\label{lemma-group-scheme-addition-tangent-vectors}
Soit $S$ un schéma. Soit $G$ un schéma en groupes sur $S$.
Soit $s \in S$. Alors la composée
$$
T_{G/S, e(s)} \oplus T_{G/S, e(s)} = T_{G \times_S G/S, (e(s), e(s))}
\rightarrow T_{G/S, e(s)}
$$
est l'addition des vecteurs tangents. Ici, le signe $=$ provient de
Variétés, Lemme \ref{varieties-lemma-tangent-space-product},
et la flèche de droite est induite par $m : G \times_S G \to G$ via
Variétés, Lemme \ref{varieties-lemma-map-tangent-spaces}.
\end{lemma}

\begin{proof}
Nous utiliserons Variétés, Équation (\ref{varieties-equation-tangent-space-fibre})
et travaillerons avec les vecteurs tangents dans les fibres.
Un élément $\theta$ du premier facteur $T_{G_s/s, e(s)}$
est l'image de $\theta$ par l'application
$T_{G_s/s, e(s)} \to T_{G_s \times G_s/s, (e(s), e(s))}$
provenant de $(1, e) : G_s \to G_s \times G_s$.
Puisque $m \circ (1, e) = 1$, nous voyons que $\theta$ s'envoie sur $\theta$
par fonctorialité. Comme l'application est linéaire, nous voyons que
$(\theta_1, \theta_2)$ s'envoie sur $\theta_1 + \theta_2$.
\end{proof}





\section{Propriétés des schémas en groupes sur un corps}
\label{section-properties-group-schemes-field}

\noindent
Dans cette section, nous rassemblons quelques propriétés des schémas en groupes sur un
corps. Dans le cas des schémas en groupes qui sont (localement) algébriques
sur un corps, nous pouvons en dire beaucoup plus, voir
section \ref{section-algebraic-group-schemes}.

\begin{lemma}
\label{lemma-group-scheme-over-field-open-multiplication}
Si $(G, m)$ est un schéma en groupes sur un corps $k$, alors le
morphisme de multiplication $m : G \times_k G \to G$ est ouvert.
\end{lemma}

\begin{proof}
Le morphisme de multiplication est isomorphe au morphisme de projection
$\text{pr}_0 : G \times_k G \to G$
car le diagramme
$$
\xymatrix{
G \times_k G \ar[d]^m \ar[rrr]_{(g, g') \mapsto (m(g, g'), g')} & & &
G \times_k G \ar[d]^{(g, g') \mapsto g} \\
G \ar[rrr]^{\text{id}} & & & G
}
$$
est commutatif, les flèches horizontales étant des isomorphismes. La projection
est ouverte d'après
Morphismes, Lemme \ref{morphisms-lemma-scheme-over-field-universally-open}.
\end{proof}

\begin{lemma}
\label{lemma-group-scheme-over-field-translate-open}
Soit $(G, m)$ un schéma en groupes sur un corps $k$. Soit $U \subset G$
un ouvert et soit $T \to G$ un morphisme de schémas. Alors l'image de
la composée $T \times_k U \to G \times_k G \to G$ est ouverte.
\end{lemma}

\begin{proof}
Pour toute extension de corps $K/k$, le morphisme $G_K \to G$ est ouvert
(Morphismes, Lemme \ref{morphisms-lemma-scheme-over-field-universally-open}).
Tout point $\xi$ de $T \times_k U$ est l'image d'un morphisme
$(t, u) : \Spec(K) \to T \times_k U$ pour un certain $K$. Alors l'image de
$T_K \times_K U_K = (T \times_k U)_K \to G_K$ contient le translaté
$t \cdot U_K$, qui est ouvert. En combinant ces faits, nous voyons que
l'image de $T \times_k U \to G$ contient un voisinage ouvert de
l'image de $\xi$. Comme $\xi$ était arbitraire, le résultat s'ensuit.
\end{proof}

\begin{lemma}
\label{lemma-group-scheme-over-field-separated}
Soit $G$ un schéma en groupes sur un corps.
Alors $G$ est un schéma séparé.
\end{lemma}

\begin{proof}
Posons $S = \Spec(k)$, où $k$ est un corps, et soit $G$ un schéma en groupes
sur $S$. D'après
le Lemme \ref{lemma-group-scheme-separated},
nous devons montrer que $e : S \to G$ est une immersion fermée. D'après
Morphismes, Lemme
\ref{morphisms-lemma-algebraic-residue-field-extension-closed-point-fibre},
l'image de $e : S \to G$ est un point fermé de $G$.
Il est clair que $\mathcal{O}_G \to e_*\mathcal{O}_S$ est surjectif,
puisque $e_*\mathcal{O}_S$ est un faisceau gratte-ciel supporté en l'élément
neutre de $G$, de valeur $k$. Nous concluons que $e$ est une immersion fermée d'après
Schémas, Lemme \ref{schemes-lemma-characterize-closed-immersions}.
\end{proof}

\begin{lemma}
\label{lemma-group-scheme-field-geometrically-irreducible}
Soit $G$ un schéma en groupes sur un corps $k$.
Alors
\begin{enumerate}
\item tout anneau local $\mathcal{O}_{G, g}$ de $G$ possède un unique idéal
premier minimal,
\item il existe exactement une composante irréductible $Z$ de $G$
passant par $e$, et
\item $Z$ est géométriquement irréductible sur $k$.
\end{enumerate}
\end{lemma}

\begin{proof}
Pour tout point $g \in G$, il existe une extension de corps
$K/k$ et un point $K$-rationnel $g' \in G(K)$ dont l'image est
$g$. Si l'on considère $g'$ comme un point $K$-rationnel du
schéma en groupes $G_K$, nous voyons alors que
$\mathcal{O}_{G, g} \to \mathcal{O}_{G_K, g'}$ est un homomorphisme local
d'anneaux fidèlement plat (puisque $G_K \to G$ est plat et qu'un homomorphisme
local plat est fidèlement plat, voir
Algèbre, Lemme \ref{algebra-lemma-local-flat-ff}).
Le résultat pour $\mathcal{O}_{G_K, g'}$ entraîne le
résultat pour $\mathcal{O}_{G, g}$, voir
Algèbre, Lemme \ref{algebra-lemma-injective-minimal-primes-in-image}.
Ainsi, pour démontrer (1), il suffit de
le démontrer pour les points $k$-rationnels $g$ de $G$. Dans ce cas,
la translation par $g$ définit un automorphisme $G \to G$
qui envoie $e$ sur $g$. Ainsi $\mathcal{O}_{G, g} \cong \mathcal{O}_{G, e}$.
On voit de cette façon que (2) implique (1), puisque les composantes irréductibles
passant par $e$ correspondent bijectivement aux idéaux premiers minimaux
de $\mathcal{O}_{G, e}$.

\medskip\noindent
Pour démontrer (2) et (3), il suffit de démontrer (2) lorsque $k$
est algébriquement clos. Dans ce cas, soient $Z_1$, $Z_2$ deux
composantes irréductibles de $G$ passant par $e$.
Puisque $k$ est algébriquement clos, le sous-schéma fermé
$Z_1 \times_k Z_2 \subset G \times_k G$ est lui aussi irréductible, voir
Variétés, Lemme \ref{varieties-lemma-bijection-irreducible-components}.
Par conséquent, $m(Z_1 \times_k Z_2)$ est contenu dans une composante irréductible
de $G$. D'autre part, il contient
$Z_1$ et $Z_2$ puisque $m|_{e \times G} = \text{id}_G$ et
$m|_{G \times e} = \text{id}_G$. Nous concluons que $Z_1 = Z_2$, comme souhaité.
\end{proof}

\begin{remark}
\label{remark-warning-group-scheme-geometrically-irreducible}
Attention : le résultat du
Lemme \ref{lemma-group-scheme-field-geometrically-irreducible}
ne signifie pas que toute composante irréductible de $G/k$ soit
géométriquement irréductible. Par exemple, le schéma en groupes
$\mu_{3, \mathbf{Q}} = \Spec(\mathbf{Q}[x]/(x^3 - 1))$
sur $\mathbf{Q}$ possède deux composantes irréductibles correspondant
à la factorisation $x^3 - 1 = (x - 1)(x^2 + x + 1)$.
Le premier facteur correspond à la composante irréductible
passant par l'élément neutre, tandis que la seconde composante irréductible
n'est pas géométriquement irréductible sur $\Spec(\mathbf{Q})$.
\end{remark}

\begin{lemma}
\label{lemma-reduced-subgroup-scheme-perfect}
Soit $G$ un schéma en groupes sur un corps parfait $k$.
Alors le réduit $G_{red}$ de $G$ est un sous-schéma en groupes fermé de $G$.
\end{lemma}

\begin{proof}
Démonstration omise. Indication : utiliser le fait que $G_{red} \times_k G_{red}$ est réduit d'après
Variétés, Lemmes \ref{varieties-lemma-perfect-reduced} et
\ref{varieties-lemma-geometrically-reduced-any-base-change}.
\end{proof}

\begin{lemma}
\label{lemma-open-subgroup-closed-over-field}
Soit $k$ un corps. Soit $\psi : G' \to G$ un morphisme de schémas en groupes
sur $k$. Si $\psi(G')$ est ouvert dans $G$, alors $\psi(G')$ est fermé dans $G$.
\end{lemma}

\begin{proof}
Posons $U = \psi(G') \subset G$. Soit $Z = G \setminus \psi(G') = G \setminus U$,
muni de la structure réduite induite de sous-schéma fermé. D'après
le Lemme \ref{lemma-group-scheme-over-field-translate-open},
l'image de
$$
Z \times_k G' \longrightarrow
Z \times_k U \longrightarrow G
$$
est ouverte (la première flèche est surjective). D'autre part, puisque $\psi$
est un homomorphisme de schémas en groupes, l'image de $Z \times_k G' \to G$
est contenue dans $Z$ (car la translation par $\psi(g')$ préserve
$U$ pour tout point $g'$ de $G'$ ; un petit détail est omis).
Par conséquent, $Z \subset G$ est un ouvert (mais pas
nécessairement un sous-schéma ouvert). Ainsi $U = \psi(G')$ est fermé.
\end{proof}

\begin{lemma}
\label{lemma-immersion-group-schemes-closed-over-field}
Soit $i : G' \to G$ une immersion de schémas en groupes sur un corps $k$.
Alors $i$ est une immersion fermée, c'est-à-dire que $i(G')$ est un sous-schéma en groupes fermé
de $G$.
\end{lemma}

\begin{proof}
Pour montrer que $i$ est une immersion fermée, il suffit de montrer que
$i(G')$ est un sous-ensemble fermé de $G$. Soit $k \subset k'$ une inclusion dans un corps parfait
qui est une extension de $k$. Si $i(G'_{k'}) \subset G_{k'}$ est fermé, alors
$i(G') \subset G$ est fermé d'après
Morphismes, Lemme \ref{morphisms-lemma-fpqc-quotient-topology}
(puisque $G_{k'} \to G$ est plat, quasi-compact et surjectif).
Nous pouvons donc supposer, et supposons, que $k$ est parfait. Nous utiliserons sans autre
mention le fait que les produits de schémas réduits sur $k$ sont réduits.
Nous pouvons remplacer $G'$ et $G$ par leurs réduits, voir
le Lemme \ref{lemma-reduced-subgroup-scheme-perfect}.
Soit $\overline{G'} \subset G$ l'adhérence de $i(G')$, considérée
comme un sous-schéma fermé réduit. D'après
Variétés, Lemme \ref{varieties-lemma-closure-of-product},
nous concluons que $\overline{G'} \times_k \overline{G'}$
est l'adhérence de l'image de $G' \times_k G' \to G \times_k G$. Par conséquent
$$
m\Big(\overline{G'} \times_k \overline{G'}\Big)
\subset \overline{G'}
$$
puisque $m$ est continu. Il s'ensuit que $\overline{G'} \subset G$
est un sous-schéma en groupes fermé (réduit). D'après
le Lemme \ref{lemma-open-subgroup-closed-over-field},
nous voyons que $i(G') \subset \overline{G'}$ est également fermé,
ce qui implique que $i(G') = \overline{G'}$, comme souhaité.
\end{proof}

\begin{lemma}
\label{lemma-irreducible-group-scheme-over-field-qc}
Soit $G$ un schéma en groupes sur un corps $k$. Si $G$ est irréductible,
alors $G$ est quasi-compact.
\end{lemma}

\begin{proof}
Supposons que $K/k$ soit une extension de corps. Si $G_K$
est quasi-compact, alors $G$ l'est aussi puisque $G_K \to G$ est surjectif.
D'après le Lemme \ref{lemma-group-scheme-field-geometrically-irreducible},
nous voyons que $G_K$ est irréductible. Il suffit donc de démontrer le lemme
après avoir remplacé $k$ par une extension. Choisissons comme $K$ un corps algébriquement
clos, de très grand cardinal. Alors, d'après
Variétés, Lemme \ref{varieties-lemma-make-Jacobson},
nous voyons que $G_K$ est un schéma de Jacobson dont tous les points fermés ont
un corps résiduel égal à $K$. Autrement dit, nous pouvons supposer que $G$ est un
schéma de Jacobson dont tous les points fermés ont pour corps résiduel $k$.

\medskip\noindent
Soit $U \subset G$ un ouvert affine non vide. Soit $g \in G(k)$. Alors
$gU \cap U \not = \emptyset$. Nous voyons donc que $g$ appartient à l'image
du morphisme
$$
U \times_{\Spec(k)} U \longrightarrow G, \quad
(u_1, u_2) \longmapsto u_1u_2^{-1}
$$
Puisque l'image de ce morphisme est ouverte
(Lemme \ref{lemma-group-scheme-over-field-open-multiplication}),
nous voyons que cette image est $G$ tout entier (car $G$ est un schéma de Jacobson
dont tous les points fermés sont $k$-rationnels).
Puisque $U$ est affine, $U \times_{\Spec(k)} U$ l'est aussi. Ainsi $G$ est
l'image d'un schéma quasi-compact, et est donc quasi-compact.
\end{proof}

\begin{lemma}
\label{lemma-connected-group-scheme-over-field-irreducible}
Soit $G$ un schéma en groupes sur un corps $k$. Si $G$ est connexe,
alors $G$ est irréductible.
\end{lemma}

\begin{proof}
D'après Variétés, Lemme
\ref{varieties-lemma-geometrically-connected-if-connected-and-point},
nous voyons que $G$ est géométriquement connexe. Si nous montrons que $G_K$
est irréductible pour une extension de corps $K/k$, alors
le lemme en découle. Nous pouvons donc appliquer
Variétés, Lemme \ref{varieties-lemma-make-Jacobson},
pour nous ramener au cas où $k$ est algébriquement clos,
$G$ est un schéma de Jacobson et tous les points fermés sont $k$-rationnels.

\medskip\noindent
Soit $Z \subset G$ l'unique composante irréductible de $G$ qui contient
l'élément neutre, voir
le Lemme \ref{lemma-group-scheme-field-geometrically-irreducible}.
Munissons $Z$ de la structure réduite induite de sous-schéma fermé ;
nous voyons que $Z \times_k Z$ est réduit et irréductible
(Variétés, Lemmes
\ref{varieties-lemma-geometrically-reduced-any-base-change} et
\ref{varieties-lemma-bijection-irreducible-components}).
Nous en concluons que $m|_{Z \times_k Z} : Z \times_k Z \to G$ se factorise
par $Z$. Ainsi, $Z$ devient un sous-schéma en groupes fermé de $G$.

\medskip\noindent
Pour obtenir une contradiction, supposons qu'il existe une autre composante
irréductible $Z' \subset G$. Alors $Z \cap Z' = \emptyset$ d'après
le Lemme \ref{lemma-group-scheme-field-geometrically-irreducible}.
D'après le Lemme \ref{lemma-irreducible-group-scheme-over-field-qc},
nous voyons que $Z$ est quasi-compact. Nous pouvons donc choisir un ouvert quasi-compact
$U \subset G$ tel que $Z \subset U$ et $U \cap Z' = \emptyset$.
L'image $W$ de $Z \times_k U \to G$ est ouverte dans $G$ d'après
le Lemme \ref{lemma-group-scheme-over-field-translate-open}.
D'autre part, $W$ est quasi-compact en tant qu'image d'un
espace quasi-compact. Nous affirmons que $W$ est fermé.

\medskip\noindent
Preuve de l'affirmation. Puisque $W$ est quasi-compact et $G$ est séparé
(Lemme \ref{lemma-group-scheme-over-field-separated}), l'inclusion
$W \to G$ est quasi-compacte d'après Schémas, Lemme
\ref{schemes-lemma-quasi-compact-permanence}. Nous en concluons que
les points de l'adhérence de $W$ sont des spécialisations de points de $W$
(Morphismes, Lemme \ref{morphisms-lemma-reach-points-scheme-theoretic-image}).
Ainsi, nous devons montrer que toute composante irréductible
$Z'' \subset G$ de $G$ qui rencontre $W$ est contenue dans $W$.
Comme $G$ est un schéma de Jacobson et que ses points fermés sont rationnels,
$Z'' \cap W$ contient un point rationnel
$g \in Z''(k) \cap W(k)$, et donc $Z'' = Zg$. Mais $W = m(Z \times_k W)$
par construction, donc $Z'' \cap W \not = \emptyset$ implique
$Z'' \subset W$.

\medskip\noindent
D'après l'affirmation, $W \subset G$ est un sous-ensemble ouvert et fermé de $G$.
Or $W \cap Z' = \emptyset$, car sinon, d'après l'argument donné dans
le paragraphe précédent, nous aurions $Z' = Zg$ pour un certain $g \in W(k)$.
Alors, puisque $Z$ est un sous-groupe, nous pourrions même choisir $g \in U(k)$, ce qui
contredirait $Z' \cap U = \emptyset$. Ainsi, $W \subset G$ est un sous-ensemble propre, ouvert
et fermé, ce qui contredit l'hypothèse que $G$ est connexe.
\end{proof}

\begin{proposition}
\label{proposition-connected-component}
Soit $G$ un schéma en groupes sur un corps $k$. Il existe un sous-schéma
en groupes fermé canonique $G^0 \subset G$ ayant les propriétés suivantes
\begin{enumerate}
\item $G^0 \to G$ est une immersion fermée plate,
\item $G^0 \subset G$ est la composante connexe neutre,
\item $G^0$ est géométriquement irréductible, et
\item $G^0$ est quasi-compact.
\end{enumerate}
\end{proposition}

\begin{proof}
Soit $G^0$ la composante connexe neutre, munie de sa structure canonique
de schéma (Morphismes, Définition
\ref{morphisms-definition-scheme-structure-connected-component}).
Pour montrer que $G^0$ est un sous-schéma en groupes fermé, nous utiliserons le
critère du Lemme \ref{lemma-closed-subgroup-scheme}.
Le morphisme $e : \Spec(k) \to G$ se factorise par $G^0$ puisque nous avons choisi
$G^0$ comme la composante connexe de $G$ qui contient $e$.
Puisque $i : G \to G$ est un automorphisme fixant $e$, nous voyons que
$i$ envoie $G^0$ dans lui-même.
D'après Variétés, Lemme \ref{varieties-lemma-geometrically-connected-criterion},
le schéma $G^0$ est géométriquement connexe sur $k$.
Ainsi, $G^0 \times_k G^0$ est connexe
(Variétés, Lemme \ref{varieties-lemma-bijection-connected-components}).
Ainsi, $m(G^0 \times_k G^0) \subset G^0$ ensemblistement.
Ainsi, $m|_{G^0 \times_k G^0} : G^0 \times_k G^0 \to G$
se factorise par $G^0$ d'après
Morphismes, Lemme \ref{morphisms-lemma-characterize-flat-closed-immersions}.
Par conséquent, $G^0$ est un sous-schéma en groupes fermé de $G$.
D'après le Lemme \ref{lemma-connected-group-scheme-over-field-irreducible},
nous voyons que $G^0$ est irréductible. D'après
le Lemme \ref{lemma-group-scheme-field-geometrically-irreducible},
nous voyons que $G^0$ est géométriquement irréductible. D'après
le Lemme \ref{lemma-irreducible-group-scheme-over-field-qc},
nous voyons que $G^0$ est quasi-compact.
\end{proof}

\begin{lemma}
\label{lemma-profinite-product-over-field}
Soit $k$ un corps. Soit $T = \Spec(A)$, où $A$ est une colimite filtrante
d'algèbres qui sont des produits finis de copies de $k$. Pour tout schéma $X$
sur $k$, on a $|T \times_k X| = |T| \times |X|$ comme espaces topologiques.
\end{lemma}

\begin{proof}
En prenant un recouvrement par des ouverts affines, on se ramène au cas où $X$ est affine.
Écrivons $X = \Spec(B)$.
Écrivons $A = \colim A_i$ avec $A_i = \prod_{t \in T_i} k$ et $T_i$ fini.
Alors $T_i = |\Spec(A_i)|$ muni de la topologie discrète, et les morphismes
de transition $A_i \to A_{i'}$ sont donnés par des applications $T_{i'} \to T_i$. Ainsi,
$|T| = \lim T_i$ comme espace topologique, voir
Limites, Lemme \ref{limits-lemma-topology-limit}. De même, nous avons
\begin{align*}
|T \times_k X| & =
|\Spec(A \otimes_k B)| \\
& =
|\Spec(\colim A_i \otimes_k B)| \\
& =
\lim |\Spec(A_i \otimes_k B)| \\
& =
\lim |\Spec(\prod\nolimits_{t \in T_i} B)| \\
& =
\lim T_i \times |X| \\
& =
(\lim T_i) \times |X| \\
& =
|T| \times |X|
\end{align*}
d'après le lemme précédent et le fait que les limites commutent entre elles.
\end{proof}

\noindent
Le lemme suivant affirme qu'en fait on peut placer une
``famille profinie algébrique de points'' dans un ouvert affine.
Nous invitons le lecteur à lire d'abord le Lemme \ref{lemma-points-in-affine}.

\begin{lemma}
\label{lemma-compact-set-in-affine}
Soit $k$ un corps algébriquement clos. Soit $G$ un schéma en groupes sur $k$.
Supposons que $G$ soit un schéma de Jacobson et que tous ses points fermés soient $k$-rationnels.
Soit $T = \Spec(A)$, où $A$ est une colimite filtrante d'algèbres qui
sont des produits finis de copies de $k$. Pour tout morphisme $f : T \to G$,
il existe un ouvert affine $U \subset G$ contenant $f(T)$.
\end{lemma}

\begin{proof}
Soit $G^0 \subset G$ le sous-schéma en groupes fermé obtenu dans la
Proposition \ref{proposition-connected-component}. Les deux premiers paragraphes
servent à se ramener au cas $G = G^0$.

\medskip\noindent
Remarquons que $T$ est une limite projective filtrante d'espaces topologiques finis
(Limites, Lemme \ref{limits-lemma-topology-limit}), donc profini comme
espace topologique (Topologie, Définition \ref{topology-definition-profinite}).
Soit $W \subset G$ un ouvert quasi-compact contenant l'image de $T \to G$.
Après avoir remplacé $W$ par l'image de $G^0 \times W \to G \times G \to G$, on peut
supposer que $W$ est invariant sous l'action de $G^0$ par translations à gauche, voir
le Lemme \ref{lemma-group-scheme-over-field-translate-open}.
Considérons la composée
$$
\psi = \pi \circ f : T \xrightarrow{f} W \xrightarrow{\pi} \pi_0(W)
$$
L'espace $\pi_0(W)$ est profini
(Topologie, Lemme \ref{topology-lemma-spectral-pi0} et
Propriétés, Lemme
\ref{properties-lemma-quasi-compact-quasi-separated-spectral}).
Soit $F_\xi \subset T$ la fibre de $T \to \pi_0(W)$ au-dessus de $\xi \in \pi_0(W)$.
Supposons que, pour tout $\xi \in \pi_0(W)$, nous ayons trouvé un ouvert affine
$U_\xi \subset W$ tel que $f(F_\xi) \subset U_\xi$. Comme $\psi$ est fermée d'après
Topologie, Lemme \ref{topology-lemma-closed-map},
$\psi(T \setminus f^{-1}(U_\xi))$ est un sous-ensemble fermé de $\pi_0(W)$
qui ne contient pas $\xi$. Son complémentaire $V_\xi$ est donc un
voisinage ouvert de $\xi$, et $\psi^{-1}(V_\xi) \subseteq f^{-1}(U_\xi)$.
Le recouvrement ouvert de $\pi_0(W)$ par les $V_\xi$ peut être raffiné en
un recouvrement fini $\pi_0(W) = V_1 \amalg \ldots \amalg V_n$
par des sous-ensembles ouverts et fermés deux à deux disjoints (Topologie,
Lemme \ref{topology-lemma-profinite-refine-open-covering}).
Pour tout $i = 1, \ldots, n$, choisissons $\xi_i \in \pi_0(W)$
tel que $V_i \subset V_{\xi_i}$. Comme $V_i$ est fermé et
$U_{\xi_i}$ affine, $U_{\xi_i} \cap \pi^{-1}(V_i)$ est également affine.
Ainsi, en remplaçant $U_{\xi_i}$ par $U_{\xi_i} \cap \pi^{-1}(V_i)$, nous
avons $\psi^{-1}(V_i) = f^{-1}(U_{\xi_i})$. Comme les $V_i$ recouvrent
$\pi_0(W)$, les $f^{-1}(U_{\xi_i})$ recouvrent $T$, et donc
$f(T) \subseteq U = \bigcup_i U_{\xi_i}$. Alors $U$ est la réunion disjointe
d'un nombre fini d'ouverts affines et est donc affine, comme souhaité.

\medskip\noindent
Soit $Z$ une composante connexe de $G$ qui rencontre $f(T)$. Alors $Z$
possède un point $k$-rationnel $z$ (car tous les corps résiduels du schéma $T$
sont isomorphes à $k$). Ainsi, $Z = G^0 z$. Par notre choix de $W$, nous voyons
que $Z \subset W$. L'argument du paragraphe précédent nous ramène
au problème de trouver un voisinage ouvert affine de $f(T) \cap Z$ dans $W$.
Après translation par un point rationnel, nous pouvons supposer que $Z = G^0$
(détails omis). Remarquons que l'image réciproque schématique
$T' = f^{-1}(G^0) \subset T$ est un sous-schéma fermé, qui est du même type.
Après avoir remplacé $T$ par $T'$, nous pouvons supposer que $f(T) \subset G^0$.
Choisissons un voisinage ouvert affine $U \subset G$
de $e \in G$, de sorte qu'en particulier $U \cap G^0$ soit non vide. Nous allons montrer
qu'il existe $g \in G^0(k)$ tel que $f(T) \subset g^{-1}U$.
Cela achèvera la preuve, car $g^{-1}U \subset W$ en vertu de l'invariance de $W$
sous les translations à gauche par $G^0$.

\medskip\noindent
Les arguments des deux paragraphes précédents permettent de passer à $G^0$
et de ramener le problème au suivant :
Supposons que $G$ soit irréductible et que $U \subset G$ soit un
voisinage ouvert affine de $e$. Montrons que $f(T) \subset g^{-1}U$
pour un certain $g \in G(k)$. Considérons le morphisme
$$
U \times_k T \longrightarrow G \times_k T,\quad
(t, u) \longrightarrow (uf(t)^{-1}, t)
$$
qui est une immersion ouverte (car le prolongement de ce morphisme à
$G \times_k T \to G \times_k T$ est un isomorphisme).
D'après notre hypothèse sur $T$, nous avons $|U \times_k T| = |U| \times |T|$
et de même pour $G \times_k T$, voir
le Lemme \ref{lemma-profinite-product-over-field}.
L'image de l'immersion ouverte affichée est donc une réunion finie
de rectangles $\bigcup_{i = 1, \ldots, n} U_i \times V_i$, avec
$V_i \subset T$ et $U_i \subset G$ ouverts quasi-compacts. Cela signifie que
les ouverts possibles $Uf(t)^{-1}$, pour $t \in T$, sont en nombre fini, disons
$Uf(t_1)^{-1}, \ldots, Uf(t_r)^{-1}$. Puisque $G$ est irréductible,
l'intersection
$$
Uf(t_1)^{-1} \cap \ldots \cap Uf(t_r)^{-1}
$$
est non vide et, puisque $G$ est un schéma de Jacobson à points fermés $k$-rationnels,
nous pouvons choisir dans cette intersection un point $k$-rationnel $g \in G(k)$.
Nous voyons alors que $g \in Uf(t)^{-1}$ pour tout $t \in T$, ce qui
signifie que $f(t) \in g^{-1}U$, comme souhaité.
\end{proof}

\begin{remark}
\label{remark-easy}
Si $G$ est un schéma en groupes sur un corps, existe-t-il toujours un sous-schéma en groupes quasi-compact,
ouvert et fermé ? D'après la
Proposition \ref{proposition-connected-component},
cette question n'est intéressante que si $G$ possède une infinité de composantes connexes
(géométriquement).
\end{remark}

\begin{lemma}
\label{lemma-group-scheme-field-countable-affine}
Soit $G$ un schéma en groupes sur un corps.
Il existe un sous-schéma ouvert et fermé $G' \subset G$
qui est une réunion dénombrable d'ouverts affines.
\end{lemma}

\begin{proof}
Soit $e \in U(k)$ un voisinage ouvert quasi-compact de l'élément
neutre. En remplaçant $U$ par $U \cap i(U)$, nous pouvons supposer que
$U$ est invariant par l'application d'inversion. Comme $G$ est séparé, c'est
encore un ensemble quasi-compact. Posons
$$
G' = \bigcup\nolimits_{n \geq 1} m_n(U \times_k \ldots \times_k U)
$$
où $m_n : G \times_k \ldots \times_k G \to G$ est l'application de multiplication à $n$ facteurs
suivante
$(g_1, \ldots, g_n) \mapsto m(m(\ldots (m(g_1, g_2), g_3), \ldots ), g_n)$.
Chacune de ces applications est ouverte (voir
Lemme \ref{lemma-group-scheme-over-field-open-multiplication}) ;
ainsi, $G'$ est un sous-schéma en groupes ouvert. D'après le
Lemme \ref{lemma-open-subgroup-closed-over-field},
c'est aussi un sous-schéma en groupes fermé.
\end{proof}






\section{Propriétés des schémas en groupes algébriques}
\label{section-algebraic-group-schemes}

\noindent
Rappelons qu'un schéma sur un corps $k$ est (localement) algébrique s'il est
(localement) de type fini sur $\Spec(k)$, voir
Variétés, Définition \ref{varieties-definition-algebraic-scheme}.
C'est dans ce sens que nous employons le terme algébrique dans le titre de cette section.

\begin{lemma}
\label{lemma-group-scheme-finite-type-field}
Soit $k$ un corps. Soit $G$ un schéma en groupes localement algébrique sur $k$.
Alors $G$ est équidimensionnel et $\dim(G) = \dim_g(G)$ pour tout $g \in G$.
Pour tout point fermé $g \in G$, on a $\dim(G) = \dim(\mathcal{O}_{G, g})$.
\end{lemma}

\begin{proof}
Montrons d'abord que $\dim_g(G) = \dim_{g'}(G)$ pour tout
couple de points $g, g' \in G$. D'après
Morphismes, Lemme \ref{morphisms-lemma-dimension-fibre-after-base-change},
nous pouvons étendre à volonté le corps de base. Nous pouvons donc supposer que
$g$ et $g'$ sont tous deux définis sur $k$. Il existe alors un
automorphisme de $G$ envoyant $g$ sur $g'$, d'où l'égalité.
D'après
Morphismes, Lemme \ref{morphisms-lemma-dimension-fibre-at-a-point},
on a
$\dim_g(G) = \dim(\mathcal{O}_{G, g}) +
\text{trdeg}_k(\kappa(g))$.
D'autre part, la dimension de $G$ (ou de tout ouvert de $G$)
est la borne supérieure des dimensions des anneaux locaux de $G$, voir
Propriétés, Lemme \ref{properties-lemma-codimension-local-ring}.
Il est clair que cette borne supérieure est atteinte aux points fermés $g$, auquel cas
$\text{trdeg}_k(\kappa(g)) = 0$ (par le théorème des zéros de Hilbert, voir
Morphismes, section \ref{morphisms-section-points-finite-type}).
Le lemme en résulte.
\end{proof}

\noindent
Le résultat suivant est parfois appelé théorème de Cartier.

\begin{lemma}
\label{lemma-group-scheme-characteristic-zero-smooth}
Soit $k$ un corps de caractéristique $0$. Soit $G$ un
schéma en groupes localement algébrique sur $k$. Alors le morphisme
structural $G \to \Spec(k)$ est lisse, c'est-à-dire que $G$ est un
schéma en groupes lisse.
\end{lemma}

\begin{proof}
D'après le
Lemme \ref{lemma-group-scheme-module-differentials},
le module des différentielles de $G$ sur $k$ est libre.
La lissité résulte donc de
Variétés, Lemme \ref{varieties-lemma-char-zero-differentials-free-smooth}.
\end{proof}

\begin{remark}
\label{remark-when-reduced}
Tout schéma en groupes sur un corps de caractéristique $0$ est réduit, voir
\cite[I, théorème 1.1 et I, corollaire 3.9, et II, théorème 2.4]{Perrin-thesis}
ainsi que
\cite[Proposition 4.2.8]{Perrin}.
Cette question avait été posée dans
\cite[page 80]{Oort}.
Nous avons vu dans le
Lemme \ref{lemma-group-scheme-characteristic-zero-smooth}
que cela est vrai lorsque le schéma en groupes est localement de type fini.
\end{remark}

\begin{lemma}
\label{lemma-reduced-group-scheme-prefect-field-characteristic-p-smooth}
Soit $k$ un corps parfait de caractéristique $p > 0$ (voir le
Lemme \ref{lemma-group-scheme-characteristic-zero-smooth}
pour le cas de caractéristique nulle).
Soit $G$ un schéma en groupes localement algébrique sur $k$.
Si $G$ est réduit, alors le morphisme
structural $G \to \Spec(k)$ est lisse, c'est-à-dire que $G$ est un
schéma en groupes lisse.
\end{lemma}

\begin{proof}
D'après le
Lemme \ref{lemma-group-scheme-module-differentials},
le faisceau $\Omega_{G/k}$ est libre. Le lemme résulte donc de
Variétés, Lemme \ref{varieties-lemma-char-p-differentials-free-smooth}.
\end{proof}

\begin{remark}
\label{remark-reduced-smooth-not-true-general}
Soit $k$ un corps de caractéristique $p > 0$.
Soit $\alpha \in k$ un élément qui n'est pas une puissance $p$-ième.
Le sous-schéma en groupes fermé
$$
G = V(x^p + \alpha y^p) \subset \mathbf{G}_{a, k}^2
$$
est réduit et irréductible, mais non lisse (pas même normal).
\end{remark}

\noindent
Le lemme suivant est un cas particulier du
Lemme \ref{lemma-compact-set-in-affine},
avec une démonstration un peu plus simple.

\begin{lemma}
\label{lemma-points-in-affine}
Soit $k$ un corps algébriquement clos.
Soit $G$ un schéma en groupes localement algébrique sur $k$.
Soient $g_1, \ldots, g_n \in G(k)$ des points $k$-rationnels.
Alors il existe un ouvert affine $U \subset G$ contenant $g_1, \ldots, g_n$.
\end{lemma}

\begin{proof}
Nous raisonnons d'abord par récurrence sur $n$ pour nous ramener au cas où tous les $g_i$
appartiennent à la même composante connexe de $G$. En effet, sinon, nous pouvons
trouver une décomposition $G = W_1 \amalg W_2$, où les $W_i$ sont ouverts dans $G$, et
(après une éventuelle renumérotation) $g_1, \ldots, g_r \in W_1$ et
$g_{r + 1}, \ldots, g_n \in W_2$ pour un certain $0 < r < n$. Par
récurrence, nous pouvons trouver des ouverts affines $U_1$ et $U_2$ de $G$ tels que
$g_1, \ldots, g_r \in U_1$ et $g_{r + 1}, \ldots, g_n \in U_2$.
Alors
$$
g_1, \ldots, g_n \in (U_1 \cap W_1) \cup (U_2 \cap W_2)
$$
fournit une solution au problème. Nous pouvons donc supposer que $g_1, \ldots, g_n$
appartiennent tous à la même composante connexe de $G$. En translatant par $g_1^{-1}$,
nous pouvons supposer que $g_1, \ldots, g_n \in G^0$, où $G^0 \subset G$ est comme dans la
Proposition \ref{proposition-connected-component}. Choisissons un voisinage
ouvert affine $U$ de $e$ ; en particulier, $U \cap G^0$ est non vide.
Comme $G^0$ est irréductible, on voit que
$$
G^0 \cap (Ug_1^{-1} \cap \ldots \cap Ug_n^{-1})
$$
est non vide. Comme $G \to \Spec(k)$ est localement de type fini, il en va de même de
$G^0 \to \Spec(k)$ ; par conséquent, tout ouvert non vide
possède un point $k$-rationnel. Nous pouvons donc choisir $g \in G^0(k)$ tel que
$g \in Ug_i^{-1}$ pour tout $i$. Alors $g_i \in g^{-1}U$ pour tout $i$,
et $g^{-1}U$ est l'ouvert affine recherché.
\end{proof}

\begin{lemma}
\label{lemma-algebraic-quasi-projective}
Soit $k$ un corps. Soit $G$ un schéma en groupes algébrique sur $k$.
Alors $G$ est quasi-projectif sur $k$.
\end{lemma}

\begin{proof}
D'après Variétés, Lemme \ref{varieties-lemma-ample-after-field-extension},
nous pouvons supposer que $k$ est algébriquement clos. Soit $G^0 \subset G$
la composante connexe neutre de $G$ comme dans la
proposition \ref{proposition-connected-component}.
Alors toute autre composante connexe de $G$ possède un point $k$-rationnel
et est donc isomorphe à $G^0$ comme schéma.
Comme $G$ est quasi-compact et noethérien, il ne possède qu'un nombre fini de
composantes connexes. Nous nous ramenons ainsi au cas considéré au
paragraphe suivant.

\medskip\noindent
Soit $G$ un schéma en groupes algébrique connexe sur un corps algébriquement clos
$k$. Si $k$ est de caractéristique nulle, alors $G$ est lisse sur
$k$ d'après le Lemme \ref{lemma-group-scheme-characteristic-zero-smooth}.
Si $k$ est de caractéristique $p > 0$, soit $H = G_{red}$
la réduction de $G$. D'après
Diviseurs, proposition \ref{divisors-proposition-push-down-ample},
il suffit de montrer que $H$ possède un faisceau inversible ample.
(Pour un schéma algébrique sur $k$, posséder un faisceau inversible
ample équivaut à être quasi-projectif sur $k$ ; voir
par exemple le résultat très général
Plus sur les morphismes, Lemme \ref{more-morphisms-lemma-quasi-projective}.)
D'après le Lemme \ref{lemma-reduced-subgroup-scheme-perfect},
on voit que $H$ est un schéma en groupes sur $k$.
D'après le Lemme \ref{lemma-reduced-group-scheme-prefect-field-characteristic-p-smooth},
on voit que $H$ est lisse sur $k$.
Cela nous ramène à la situation considérée au
paragraphe suivant.

\medskip\noindent
Soit $G$ un schéma en groupes lisse, irréductible et quasi-compact sur un
corps algébriquement clos $k$. Observons que les anneaux locaux de $G$
sont réguliers et donc factoriels
(Variétés, Lemme \ref{varieties-lemma-smooth-regular} et
Plus sur l'algèbre, Lemme \ref{more-algebra-lemma-regular-local-UFD}).
Le complémentaire d'un ouvert affine non vide de $G$
est le support d'un diviseur de Cartier effectif $D$.
Cela résulte de Diviseurs, Lemme
\ref{divisors-lemma-complement-open-affine-effective-cartier-divisor}.
(Observons que $G$ est séparé d'après le
Lemme \ref{lemma-group-scheme-over-field-separated}.)
Nous en concluons qu'il existe un diviseur de Cartier effectif $D \subset G$
tel que $G \setminus D$ soit affine. Nous utiliserons ci-dessous que,
pour tout $n \geq 1$ et tous $g_1, \ldots, g_n \in G(k)$, le complémentaire
$G \setminus \bigcup D g_i$ est affine. En effet, c'est l'intersection
des ouverts affines $G \setminus Dg_i \cong G \setminus D$
dans le schéma séparé $G$.

\medskip\noindent
Nous pouvons choisir la ligne supérieure du diagramme
$$
\xymatrix{
G & U \ar[l]_j \ar[r]^\pi & \mathbf{A}^d_k \\
& W \ar[r]^{\pi'} \ar[u] & V \ar[u]
}
$$
de telle sorte que $U \not = \emptyset$, que $j : U \to G$ soit une immersion ouverte et que
$\pi$ soit \'etale, voir
Morphismes, Lemme \ref{morphisms-lemma-smooth-etale-over-affine-space}.
Il existe un ouvert affine non vide $V \subset \mathbf{A}^d_k$ tel que,
si $W = \pi^{-1}(V)$, le morphisme $\pi' = \pi|_W : W \to V$ soit fini \'etale.
En particulier, $\pi'$ est fini localement libre, disons de degré $n$.
Considérons le diviseur de Cartier effectif
$$
\mathcal{D} = \{(g, w) \mid m(g, j(w)) \in D\} \subset G \times W
$$
(Il s'agit de la restriction à $G \times W$ de l'image réciproque de $D \subset G$
par le morphisme plat $m : G \times G \to G$.)
Considérons le fermé\footnote{En utilisant les résultats
de Diviseurs, section \ref{divisors-section-norms},
nous pourrions prendre pour diviseur de Cartier effectif
$E$ la norme du diviseur de Cartier effectif $\mathcal{D}$
le long du morphisme fini localement libre $1 \times \pi'$, ce qui permettrait d'éviter
une partie des arguments.}
$T = (1 \times \pi')(\mathcal{D}) \subset G \times V$.
Comme $\pi'$ est fini localement libre, toute composante irréductible
de $T$ est de codimension $1$ dans $G \times V$. Comme $G \times V$
est lisse sur $k$, nous en concluons que ces composantes sont des diviseurs de Cartier
effectifs (Diviseurs, Lemme \ref{divisors-lemma-weil-divisor-is-cartier-UFD}
et les lemmes cités ci-dessus),
et donc que $T$ est le support d'un diviseur de Cartier effectif
$E$ dans $G \times V$. Si $v \in V(k)$, alors
$(\pi')^{-1}(v) = \{w_1, \ldots, w_n\} \subset W(k)$ et l'on voit que
$$
E_v = \bigcup\nolimits_{i = 1, \ldots, n} D j(w_i)^{-1}
$$
dans $G$, du point de vue ensembliste. En particulier, on voit que $G \setminus E_v$
est un ouvert affine (voir ci-dessus).
De plus, si $g \in G(k)$, il existe $v \in V$
tel que $g \not \in E_v$. En effet, l'ensemble $W'$ des $w \in W$ tels que
$g \not \in Dj(w)^{-1}$ est un ouvert non vide, et il suffit de choisir $v$
tel que la fibre de $W' \to V$ au-dessus de $v$ possède $n$ éléments.

\medskip\noindent
Considérons le faisceau inversible
$\mathcal{M} = \mathcal{O}_{G \times V}(E)$ sur $G \times V$.
D'après Variétés, Lemme \ref{varieties-lemma-rational-equivalence-for-Pic},
la classe d'isomorphisme $\mathcal{L}$ de la restriction
$\mathcal{M}_v = \mathcal{O}_G(E_v)$ est indépendante de $v \in V(k)$.
D'autre part, pour tout $g \in G(k)$, nous pouvons trouver $v$
tel que $g \not \in E_v$ et que $G \setminus E_v$
soit affine. Ainsi, la section canonique
(Diviseurs, Définition
\ref{divisors-definition-invertible-sheaf-effective-Cartier-divisor})
de $\mathcal{O}_G(E_v)$
correspond à une section $s_v$ de $\mathcal{L}$ qui ne
s'annule pas en $g$ et telle que $G_{s_v}$ soit affine.
Cela signifie que $\mathcal{L}$ est ample par définition
(Propriétés, Définition \ref{properties-definition-ample}).
\end{proof}

\begin{lemma}
\label{lemma-algebraic-center}
Soit $k$ un corps. Soit $G$ un schéma en groupes localement algébrique sur $k$.
Alors le centre de $G$ est un sous-schéma en groupes fermé de $G$.
\end{lemma}

\begin{proof}
Notons $\text{Aut}(G)$ le foncteur contravariant sur la catégorie des
schémas sur $k$ qui associe à $S/k$ l'ensemble des automorphismes
du schéma $G_S$ obtenu par changement de base, en tant que schéma en groupes sur $S$. Il existe une transformation
naturelle
$$
G \longrightarrow \text{Aut}(G),\quad
g \longmapsto \text{inn}_g
$$
qui, pour tout $S$, associe à un point $g$ de $G$ à valeurs dans ce schéma l'automorphisme intérieur de $G$
déterminé par $g$. Le centre $C$ de $G$ est, par définition, le noyau de
cette transformation, c'est-à-dire le foncteur qui associe à $S$ les
$g \in G(S)$ dont l'automorphisme intérieur associé est trivial. L'énoncé
du lemme affirme que ce foncteur est représentable par un sous-schéma en groupes
fermé de $G$.

\medskip\noindent
Choisissons un entier $n \geq 1$. Soit $G_n \subset G$ le $n$-ième voisinage
infinitésimal de l'élément neutre $e$ de $G$. Pour tout schéma $S/k$,
le changement de base $G_{n, S}$ est le $n$-ième voisinage infinitésimal de
$e_S : S \to G_S$. Nous voyons ainsi qu'il existe une transformation naturelle
$\text{Aut}(G) \to \text{Aut}(G_n)$, où le membre de droite est le
foncteur des automorphismes de $G_n$ comme schéma ($G_n$ n'est pas, en général,
un schéma en groupes). Observons que $G_n$ est le spectre d'un anneau
local artinien $A_n$, de corps résiduel $k$ et de dimension finie
comme espace vectoriel sur $k$
(Variétés, Lemme \ref{varieties-lemma-algebraic-scheme-dim-0}).
Comme tout automorphisme de $G_n$ induit en particulier une application
linéaire inversible $A_n \to A_n$, nous obtenons des transformations de foncteurs
$$
G \to \text{Aut}(G) \to \text{Aut}(G_n) \to \text{GL}(A_n)
$$
Le dernier foncteur à valeurs dans les groupes est représentable, voir
Exemple \ref{example-general-linear-group}, et la
dernière flèche est manifestement injective.
Ainsi, pour tout $n$, nous obtenons un sous-schéma en groupes fermé
$$
H_n = \Ker(G \to \text{Aut}(G_n)) = \Ker(G \to \text{GL}(A_n)).
$$
En première approximation, posons $H = \bigcap_{n \geq 1} H_n$
(intersection schématique). C'est un sous-schéma en groupes fermé
qui contient le centre $C$.

\medskip\noindent
Soit $h$ un $S$-point de $H$, où $S$ est localement noethérien.
Alors l'automorphisme $\text{inn}_h$ induit l'identité sur tous
les sous-schémas fermés $G_{n, S}$. Considérons le noyau
$K = \Ker(\text{inn}_h : G_S \to G_S)$.
C'est un sous-schéma en groupes fermé de $G_S$ sur $S$ contenant les
sous-schémas fermés $G_{n, S}$ pour $n \geq 1$.
Cela implique que $K$ contient un voisinage ouvert de
$e(S) \subset G_S$, voir
Algèbre, Remarque \ref{algebra-remark-intersection-powers-ideal}.
Soit $G^0 \subset G$ comme dans la proposition \ref{proposition-connected-component}.
Comme $G^0$ est géométriquement irréductible, nous en concluons que $K$ contient
$G^0_S$ (pour tout ouvert non vide $U \subset G^0_{k'}$ et toute extension de corps
$k'/k$, on a $U \cdot U^{-1} = G^0_{k'}$ ; voir la démonstration du
Lemme \ref{lemma-irreducible-group-scheme-over-field-qc}).
En appliquant ceci avec $S = H$, nous constatons que $G^0$ et $H$
sont des sous-schémas en groupes de $G$ dont les points commutent : pour tout schéma $S$
et tous les points à valeurs dans $S$, $g \in G^0(S)$ et $h \in H(S)$, on a
$gh = hg$ dans $G(S)$.

\medskip\noindent
Supposons que $k$ soit algébriquement clos. Nous pouvons alors choisir un point à valeurs dans $k$,
noté $g_i$, dans chaque composante irréductible $G_i$ de $G$. Observons que,
dans ce cas, les composantes connexes de $G$ sont les composantes irréductibles
de $G$, lesquelles sont les translatées de $G^0$ par nos $g_i$. Nous affirmons que
$$
C = H \cap \bigcap\nolimits_i \Ker(\text{inn}_{g_i} : G \to G)
\quad (\text{intersection schématique})
$$
En effet, $C$ est contenu dans le membre de droite. Réciproquement, tout
point à valeurs dans $S$, noté $h$, du membre de droite commute avec $G^0$
et avec les $g_i$, donc avec tout élément de $G = \bigcup G^0g_i$.

\medskip\noindent
Le cas d'un corps de base quelconque $k$ découle du résultat pour la
clôture algébrique $\overline{k}$ par descente. En effet, soit
$A \subset G_{\overline{k}}$ le sous-schéma en groupes fermé représentant
le centre de $G_{\overline{k}}$. Alors on a
$$
A \times_{\Spec(k)} \Spec(\overline{k}) =
\Spec(\overline{k}) \times_{\Spec(k)} A
$$
comme sous-schémas fermés de $G_{\overline{k} \otimes_k \overline{k}}$, par
fonctorialité du centre. On voit donc que $A$ descend
en un sous-schéma en groupes fermé $Z \subset G$ d'après
Descente, Lemme \ref{descent-lemma-closed-immersion}
(et Descente, Lemme \ref{descent-lemma-descending-property-closed-immersion}).
Alors $Z$ représente $C$ (nous omettons un petit argument), ce qui achève la démonstration.
\end{proof}







\section{Variétés abéliennes}
\label{section-abelian-varieties}

\noindent
Une excellente référence pour ce sujet est le livre de Mumford sur les
variétés abéliennes, voir \cite{AVar}. Nous encourageons le lecteur à
le consulter. Il existe de nombreuses définitions équivalentes ; en voici une.

\begin{definition}
\label{definition-abelian-variety}
Soit $k$ un corps. Une {\it variété abélienne} est un schéma en groupes sur
$k$ qui est aussi une variété propre et géométriquement intègre sur
$k$\footnote{Pour des définitions équivalentes, voir la Remarque \ref{remark-16-equivalent}.}.
\end{definition}

\noindent
Nous démontrons quelques lemmes sur cette notion, puis nous rassemblons
tous les résultats dans la
proposition \ref{proposition-review-abelian-varieties}.

\begin{lemma}
\label{lemma-abelian-variety-projective}
Soit $k$ un corps. Soit $A$ une variété abélienne sur $k$.
Alors $A$ est projective.
\end{lemma}

\begin{proof}
Cela résulte du
Lemme \ref{lemma-algebraic-quasi-projective} et de
Plus sur les morphismes, Lemme \ref{more-morphisms-lemma-projective}.
\end{proof}

\begin{lemma}
\label{lemma-abelian-variety-change-field}
Soit $k$ un corps. Soit $A$ une variété abélienne sur $k$.
Pour toute extension de corps $K/k$, le changement de base $A_K$ est une
variété abélienne sur $K$.
\end{lemma}

\begin{proof}
Démonstration omise. Notons que c'est la raison pour laquelle nous avons exigé que $A$ soit
géométriquement intègre ; sans cette condition, ce lemme
(et beaucoup d'autres ci-dessous) serait faux.
\end{proof}

\begin{lemma}
\label{lemma-abelian-variety-smooth}
Soit $k$ un corps. Soit $A$ une variété abélienne sur $k$.
Alors $A$ est lisse sur $k$.
\end{lemma}

\begin{proof}
Si $k$ est parfait, cela résulte du
Lemme \ref{lemma-group-scheme-characteristic-zero-smooth}
(caractéristique nulle) et du
Lemme \ref{lemma-reduced-group-scheme-prefect-field-characteristic-p-smooth}
(caractéristique positive).
Nous pouvons nous ramener à ce cas dans la situation générale par descente de la lissité
(Descente, Lemme \ref{descent-lemma-descending-property-smooth})
et en passant à la clôture parfaite grâce au
Lemme \ref{lemma-abelian-variety-change-field}.
\end{proof}

\begin{lemma}
\label{lemma-abelian-variety-abelian}
Une variété abélienne est un schéma en groupes commutatif, c'est-à-dire que la loi
de groupe est commutative.
\end{lemma}

\begin{proof}
Soit $k$ un corps. Soit $A$ une variété abélienne sur $k$.
D'après le Lemme \ref{lemma-abelian-variety-change-field}, nous pouvons remplacer
$k$ par sa clôture algébrique. Considérons le morphisme
$$
h : A \times_k A \longrightarrow A \times_k A,\quad
(x, y) \longmapsto (x, xyx^{-1}y^{-1})
$$
C'est un morphisme au-dessus de $A$ pour la première projection de chaque côté.
Soit $e \in A(k)$ l'élément neutre. Alors on voit que $h|_{e \times A}$ est
constant de valeur $(e, e)$. D'après Plus sur les morphismes, Lemme
\ref{more-morphisms-lemma-flat-proper-family-cannot-collapse-fibre},
il existe un voisinage ouvert $U \subset A$ de $e$
tel que $h|_{U \times A}$ se factorise par un certain $Z \subset U \times A$
fini sur $U$. Cela signifie que, pour $x \in U(k)$, le morphisme
$A \to A$, $y \mapsto xyx^{-1}y^{-1}$, ne prend qu'un nombre fini de valeurs.
Bien entendu, cela signifie qu'il est constant de valeur $e$. Ainsi,
$(x, y) \mapsto xyx^{-1}y^{-1}$ est
constant de valeur $e$ sur $U \times A$, ce qui implique
que la loi de groupe sur $A$ est commutative.
\end{proof}

\begin{lemma}
\label{lemma-apply-cube}
Soit $k$ un corps. Soit $A$ une variété abélienne sur $k$.
Soit $\mathcal{L}$ un $\mathcal{O}_A$-module inversible.
Il existe alors un isomorphisme
$$
m_{1, 2, 3}^*\mathcal{L} \otimes
m_1^*\mathcal{L} \otimes
m_2^*\mathcal{L} \otimes
m_3^*\mathcal{L} \cong
m_{1, 2}^*\mathcal{L} \otimes
m_{1, 3}^*\mathcal{L} \otimes
m_{2, 3}^*\mathcal{L}
$$
de modules inversibles sur $A \times_k A \times_k A$,
où $m_{i_1, \ldots, i_t} : A \times_k A \times_k A \to A$
est le morphisme $(x_1, x_2, x_3) \mapsto \sum x_{i_j}$.
\end{lemma}

\begin{proof}
Appliquons le théorème du cube
(Plus sur les morphismes, Théorème \ref{more-morphisms-theorem-of-the-cube})
à la différence
$$
\mathcal{M} =
m_{1, 2, 3}^*\mathcal{L} \otimes
m_1^*\mathcal{L} \otimes
m_2^*\mathcal{L} \otimes
m_3^*\mathcal{L} \otimes
m_{1, 2}^*\mathcal{L}^{\otimes -1} \otimes
m_{1, 3}^*\mathcal{L}^{\otimes -1} \otimes
m_{2, 3}^*\mathcal{L}^{\otimes -1}
$$
Cela s'applique car la restriction de $\mathcal{M}$
à $A \times A \times e = A \times A$ est égale à
$$
n_{1, 2}^*\mathcal{L} \otimes
n_1^*\mathcal{L} \otimes
n_2^*\mathcal{L} \otimes
n_{1, 2}^*\mathcal{L}^{\otimes -1} \otimes
n_1^*\mathcal{L}^{\otimes -1} \otimes
n_2^*\mathcal{L}^{\otimes -1} \cong \mathcal{O}_{A \times_k A}
$$
où $n_{i_1, \ldots, i_t} : A \times_k A \to A$
est le morphisme $(x_1, x_2) \mapsto \sum x_{i_j}$.
Il en va de même pour $A \times e \times A$ et $e \times A \times A$.
\end{proof}

\begin{lemma}
\label{lemma-pullbacks-by-n}
Soit $k$ un corps. Soit $A$ une variété abélienne sur $k$.
Soit $\mathcal{L}$ un $\mathcal{O}_A$-module inversible.
Alors
$$
[n]^*\mathcal{L} \cong
\mathcal{L}^{\otimes n(n + 1)/2} \otimes
([-1]^*\mathcal{L})^{\otimes n(n - 1)/2}
$$
où $[n] : A \to A$ envoie $x$ sur $x + x + \ldots + x$, somme de $n$ termes,
et où $[-1] : A \to A$ est l'application d'inversion de $A$.
\end{lemma}

\begin{proof}
Considérons le morphisme $A \to A \times_k A \times_k A$,
$x \mapsto (x, x, -x)$, où $-x = [-1](x)$. En prenant l'image réciproque
de la relation du Lemme \ref{lemma-apply-cube}, nous obtenons
$$
\mathcal{L} \otimes
\mathcal{L} \otimes
\mathcal{L} \otimes
[-1]^*\mathcal{L} \cong
[2]^*\mathcal{L}
$$
ce qui démontre le résultat pour $n = 2$. Supposons par récurrence que le résultat vaille
pour $1, 2, \ldots, n$. Considérons alors le morphisme
$A \to A \times_k A \times_k A$, $x \mapsto (x, x, [n - 1]x)$.
En prenant l'image réciproque
de la relation du Lemme \ref{lemma-apply-cube}, nous obtenons
$$
[n + 1]^*\mathcal{L} \otimes
\mathcal{L} \otimes
\mathcal{L} \otimes
[n - 1]^*\mathcal{L} \cong
[2]^*\mathcal{L} \otimes
[n]^*\mathcal{L} \otimes
[n]^*\mathcal{L}
$$
et le résultat découle d'un calcul élémentaire.
\end{proof}

\begin{lemma}
\label{lemma-degree-multiplication-by-d}
Soit $k$ un corps. Soit $A$ une variété abélienne sur $k$.
Soit $[d] : A \to A$ la multiplication par $d$.
Alors $[d]$ est fini localement libre de degré $d^{2\dim(A)}$.
\end{lemma}

\begin{proof}
{\emergencystretch=1em D'après le Lemme \ref{lemma-abelian-variety-projective}
(et Plus sur les morphismes, Lemme \ref{more-morphisms-lemma-projective}),
on voit que $A$ possède un module inversible ample $\mathcal{L}$.
Comme $[-1] : A \to A$ est un automorphisme, on voit que
$[-1]^*\mathcal{L}$ est lui aussi un $\mathcal{O}_A$-module inversible ample.
Ainsi, $\mathcal{N} = \mathcal{L} \otimes [-1]^*\mathcal{L}$
est ample, voir
Propriétés, Lemme \ref{properties-lemma-ample-tensor-globally-generated}.
Comme $\mathcal{N} \cong [-1]^*\mathcal{N}$, on voit que
$[d]^*\mathcal{N} \cong \mathcal{N}^{\otimes d^2}$ d'après le
Lemme \ref{lemma-pullbacks-by-n}.\par}

\medskip\noindent
Pour obtenir une contradiction, soit $C$ une courbe propre contenue dans une
fibre de $[d]$. Alors $\mathcal{N}^{\otimes d^2}|_C \cong \mathcal{O}_C$
est un $\mathcal{O}_C$-module inversible ample de degré $0$, ce qui
contredit par exemple Variétés, Lemme \ref{varieties-lemma-ample-curve}.
(On peut aussi utiliser Variétés, Lemme \ref{varieties-lemma-ample-positive}.)
Ainsi, toute fibre de $[d]$ est de dimension $0$, et $[d]$ est donc fini,
par exemple d'après Cohomologie des schémas, Lemme
\ref{coherent-lemma-characterize-finite}.
De plus, comme $A$ est lisse sur $k$ d'après le
Lemme \ref{lemma-abelian-variety-smooth},
on voit que $[d] : A \to A$ est plat d'après
Algèbre, Lemme \ref{algebra-lemma-CM-over-regular-flat}
(on utilise aussi que les schémas lisses sur un corps sont réguliers et que
les anneaux réguliers sont des anneaux de Cohen-Macaulay, voir
Variétés, lemme \ref{varieties-lemma-smooth-regular} et
Algèbre, lemme \ref{algebra-lemma-regular-ring-CM}).
Ainsi, $[d]$ est fini et plat, donc fini localement libre d'après
Morphismes, lemme \ref{morphisms-lemma-finite-flat}.

\medskip\noindent
Enfin, venons-en à la formule du degré. D'après
Variétés, lemme \ref{varieties-lemma-degree-finite-morphism-in-terms-degrees},
on voit que
$$
\deg_{\mathcal{N}^{\otimes d^2}}(A) = \deg([d]) \deg_\mathcal{N}(A)
$$
Puisque le degré de $A$ relativement à
$\mathcal{N}^{\otimes d^2}$, respectivement à $\mathcal{N}$,
est le coefficient de $n^{\dim(A)}$ dans le polynôme
$$
n \longmapsto \chi(A, \mathcal{N}^{\otimes nd^2}),\quad
\text{respectivement}\quad n \longmapsto \chi(A, \mathcal{N}^{\otimes n})
$$
on voit que $\deg([d]) = d^{2 \dim(A)}$.
\end{proof}

\begin{lemma}
\label{lemma-abelian-variety-multiplication-by-d-etale}
\begin{slogan}
La multiplication par un entier sur une variété abélienne est un morphisme étale
si et seulement si cet entier est inversible dans le corps de base.
\end{slogan}
Soit $k$ un corps. Soit $A$ une variété abélienne non nulle sur $k$.
Alors $[d] : A \to A$ est \'etale si et seulement si $d$ est inversible dans $k$.
\end{lemma}

\begin{proof}
Observons que $[d](x + y) = [d](x) + [d](y)$. Comme la translation par un
point est un automorphisme de $A$, on voit que le lieu où
$[d] : A \to A$ est \'etale est soit vide, soit égal à $A$ (certains détails
sont omis). Il suffit donc de déterminer si $[d]$ est \'etale en
l'élément neutre $e \in A(k)$. Puisque l'on sait que $[d]$ est fini localement libre
(lemme \ref{lemma-degree-multiplication-by-d}),
dire qu'il est \'etale en $e$ équivaut à
montrer que $\text{d}[d] : T_{A/k, e} \to T_{A/k, e}$ est injective. Voir
Variétés, lemme \ref{varieties-lemma-injective-tangent-spaces-unramified} et
Morphismes, lemme \ref{morphisms-lemma-flat-unramified-etale}.
Le lemme \ref{lemma-group-scheme-addition-tangent-vectors} montre que
$\text{d}[d]$ est donnée par la multiplication par $d$ sur $T_{A/k, e}$.
\end{proof}

\begin{lemma}
\label{lemma-abelian-variety-multiplication-by-p}
Soit $k$ un corps de caractéristique $p > 0$. Soit $A$ une variété abélienne
de dimension $g$
sur $k$. La fibre de $[p] : A \to A$ au-dessus de $0$ possède au plus
$p^g$ points distincts.
\end{lemma}

\begin{proof}
Pour le démontrer, on peut remplacer $k$ par sa clôture algébrique, ce que l'on fait.
D'après le lemme \ref{lemma-group-scheme-addition-tangent-vectors},
la différentielle de $[p]$ est la multiplication par $p$ en tant qu'application
$T_{A/k, e} \to T_{A/k, e}$ et est donc nulle (comparer
avec la démonstration du lemme \ref{lemma-abelian-variety-multiplication-by-d-etale}).
Comme $[p]$ commute aux translations, on en conclut que la différentielle de $[p]$
est nulle partout, c'est-à-dire que l'application induite
$[p]^*\Omega_{A/k} \to \Omega_{A/k}$ est nulle.
En passant aux points génériques, on voit que
l'homomorphisme correspondant $[p]^* : k(A) \to k(A)$
entre corps de fonctions induit l'application nulle sur $\Omega_{k(A)/k}$.
Soit $t_1, \ldots, t_g$ une p-base de $k(A)$ sur $k$
(Compléments d'algèbre, définition \ref{more-algebra-definition-p-basis} et
lemme \ref{more-algebra-lemma-p-basis}). Alors $[p]^*(t_i)$
possède une racine $p$-ième d'après
Algèbre, lemme \ref{algebra-lemma-derivative-zero-pth-power}.
On en conclut que
$k(A)[x_1, \ldots, x_g]/(x_1^p - t_1, \ldots, x_g^p - t_g)$ est une sous-extension
de celle définie par $[p]^* : k(A) \to k(A)$.
On peut donc trouver un ouvert affine $U \subset A$ tel que
$t_i \in \mathcal{O}_A(U)$ et $x_i \in \mathcal{O}_A([p]^{-1}(U))$.
On obtient une factorisation
$$
[p]^{-1}(U)
\xrightarrow{\pi_1}
\Spec(\mathcal{O}(U)[x_1, \ldots, x_g]/(x_1^p - t_1, \ldots, x_g^p - t_g))
\xrightarrow{\pi_2}
U
$$
de $[p]$ au-dessus de $U$. Quitte à rétrécir $U$, on peut supposer que $\pi_1$
est fini localement libre (par exemple par platitude générique -- en fait, il est
déjà fini localement libre dans notre cas).
D'après le lemme \ref{lemma-degree-multiplication-by-d}, on voit que
$[p]$ est de degré $p^{2g}$. Comme $\pi_2$
est de degré $p^g$, on voit que $\pi_1$ est également de degré $p^g$.
Le morphisme $\pi_2$ est un homéomorphisme universel ; ses fibres sont donc
réduites à un point. On en conclut que les fibres ensemblistes de $[p]^{-1}(U) \to U$
sont les fibres de $\pi_1$. Elles
ont donc au plus $p^g$ éléments. Comme $[p]$ est un homomorphisme de schémas
en groupes sur $k$, la fibre de $[p] : A(k) \to A(k)$ possède le
même cardinal pour tout $a \in A(k)$, ce qui achève la démonstration.
\end{proof}

\begin{proposition}
\label{proposition-review-abelian-varieties}
\begin{reference}
Exposé de manière remarquable dans \cite{AVar}.
\end{reference}
Soit $A$ une variété abélienne sur un corps $k$. Alors
\begin{enumerate}
\item $A$ est projective sur $k$,
\item $A$ est un schéma en groupes commutatif,
\item le morphisme $[n] : A \to A$ est surjectif pour tout $n \geq 1$,
\item si $k$ est algébriquement clos, alors $A(k)$ est un groupe abélien divisible,
\item $A[n] = \Ker([n] : A \to A)$ est un schéma en groupes fini de degré
$n^{2\dim A}$ sur $k$,
\item $A[n]$ est \'etale sur $k$ si et seulement si $n \in k^*$,
\item si $n \in k^*$ et si $k$ est algébriquement clos,
alors $A(k)[n] \cong (\mathbf{Z}/n\mathbf{Z})^{\oplus 2\dim(A)}$,
\item si $k$ est algébriquement clos de caractéristique $p > 0$, alors
il existe un entier $0 \leq f \leq \dim(A)$ tel que
$A(k)[p^m] \cong (\mathbf{Z}/p^m\mathbf{Z})^{\oplus f}$
pour tout $m \geq 1$.
\end{enumerate}
\end{proposition}

\begin{proof}
L'assertion (1) découle du lemme \ref{lemma-abelian-variety-projective}.
L'assertion (2) découle du lemme \ref{lemma-abelian-variety-abelian}.
L'assertion (3) découle du lemme \ref{lemma-degree-multiplication-by-d}.
Si $k$ est algébriquement clos, les morphismes surjectifs de variétés
sur $k$ induisent des applications surjectives sur les points $k$-rationnels, donc
(4) découle de (3).
L'assertion (5) découle du lemme \ref{lemma-degree-multiplication-by-d}
et du fait qu'un changement de base d'un morphisme fini localement libre
de degré $N$ est un morphisme fini localement libre de degré $N$.
L'assertion (6) découle du 
lemme \ref{lemma-abelian-variety-multiplication-by-d-etale}.
En effet, si $n$ est inversible dans $k$, alors $[n]$ est \'etale
et donc $A[n]$ est \'etale sur $k$.
D'autre part, si $n$ n'est pas inversible dans $k$, alors
$[n]$ n'est pas \'etale en $e$ et il en résulte que $A[n]$
n'est pas \'etale sur $k$ en $e$ (utiliser
Morphismes, lemmes \ref{morphisms-lemma-flat-unramified-etale} et
\ref{morphisms-lemma-set-points-where-fibres-unramified}).

\medskip\noindent
Supposons $k$ algébriquement clos. Posons $g = \dim(A)$. Démonstration de (7).
Soit $\ell$ un nombre premier inversible dans $k$. On voit alors que
$$
A[\ell](k) = A(k)[\ell]
$$
est un groupe abélien fini, annulé par $\ell$, d'ordre $\ell^{2g}$.
Il s'ensuit qu'il est isomorphe à $(\mathbf{Z}/\ell\mathbf{Z})^{2g}$
d'après le théorème de structure des groupes abéliens finis. Considérons ensuite
la suite exacte courte
$$
0 \to A(k)[\ell] \to A(k)[\ell^2] \xrightarrow{\ell} A(k)[\ell] \to 0
$$
En raisonnant de la même manière que ci-dessus, nous concluons que 
$A(k)[\ell^2] \cong (\mathbf{Z}/\ell^2\mathbf{Z})^{2g}$.
Par récurrence sur l'exposant, nous obtenons
$A(k)[\ell^m] \cong (\mathbf{Z}/\ell^m\mathbf{Z})^{2g}$.
Pour les entiers composés $n$ premiers à la caractéristique de $k$,
nous considérons les composantes primaires et obtenons la forme correcte de la
$n$-torsion de $A(k)$.
La démonstration de (8) procède exactement de la même manière, en utilisant que
le lemme \ref{lemma-abelian-variety-multiplication-by-p} donne
$A(k)[p] \cong (\mathbf{Z}/p\mathbf{Z})^{\oplus f}$ pour un certain $0 \leq f \leq g$.
\end{proof}

\begin{remark}
\label{remark-16-equivalent}
Soit $k$ un corps. Il existe $2 \times 4 \times 2 = 16$ définitions
équivalentes des variétés abéliennes. Soient
\begin{itemize}
\item projectif, propre,
\item {\emergencystretch=2em géométriquement irréductible, irréductible, géométriquement connexe,
connexe,\par}
\item lisse, géométriquement réduit
\end{itemize}
trois ensembles de propriétés ; choisissons-en une dans chacun d'eux, et soit $A$
un schéma en groupes sur $k$ possédant les propriétés choisies sur $k$.
Alors $A$ est une variété abélienne.
Si nous choisissons « propre, géométriquement irréductible, géométriquement
réduit », nous retrouvons la définition \ref{definition-abelian-variety}
(utiliser Variétés, lemme \ref{varieties-lemma-geometrically-integral}).
Les choix les plus faibles possibles seraient « propre, connexe
et géométriquement réduit », voir par exemple
Morphismes, lemme \ref{morphisms-lemma-locally-projective-proper} et
Variétés, lemme \ref{varieties-lemma-smooth-geometrically-normal}.
Supposons donc que $A$ soit un schéma en groupes propre, connexe
et géométriquement réduit sur $k$.
Alors $A$ est géométriquement irréductible d'après les
lemmes \ref{lemma-connected-group-scheme-over-field-irreducible} et
\ref{lemma-group-scheme-field-geometrically-irreducible},
et est donc une variété abélienne.
Enfin, si $A/k$ est une variété abélienne, alors elle est projective
et lisse sur $k$ (proposition
\ref{proposition-review-abelian-varieties}) ; elle satisfait donc
les choix les plus forts possibles « projectif, géométriquement irréductible, lisse ».
\end{remark}







\section{Actions de schémas en groupes}
\label{section-action-group-scheme}

\noindent
Soient $(G, m)$ un groupe et $V$ un ensemble.
Rappelons qu'une {\it action (à gauche)} de $G$ sur $V$ est donnée
par une application $a : G \times V \to V$ telle que
\begin{enumerate}
\item (associativité) $a(m(g, g'), v) = a(g, a(g', v))$ pour tous
$g, g' \in G$ et $v \in V$, et
\item (élément neutre) $a(e, v) = v$ pour tout $v \in V$.
\end{enumerate}
On dit également que $V$ est un {\it $G$-ensemble} (cela signifie en général que nous
omettons $a$ de la notation --- ce qui constitue un abus de notation).
Un {\it morphisme de $G$-ensembles} $\psi : V \to V'$ est une application
telle que $\psi(a(g, v)) = a(g, \psi(v))$ pour tout $v \in V$.

\begin{definition}
\label{definition-action-group-scheme}
Soit $S$ un schéma. Soit $(G, m)$ un schéma en groupes sur $S$.
\begin{enumerate}
\item Une {\it action de $G$ sur le schéma $X/S$} est
un morphisme $a : G \times_S X \to X$ sur $S$ tel que
pour tout $T/S$, l'application $a : G(T) \times X(T) \to X(T)$
définit la structure d'un $G(T)$-ensemble sur $X(T)$.
\item Supposons que $X$, $Y$ soient des schémas sur $S$, chacun muni
d'une action de $G$. Un morphisme {\it équivariant}, ou plus précisément
un morphisme {\it $G$-équivariant} $\psi : X \to Y$,
est un morphisme de schémas sur $S$
tel que pour tout $T/S$, l'application $\psi : X(T) \to Y(T)$ soit
un morphisme de $G(T)$-ensembles.
\end{enumerate}
\end{definition}

\noindent
Dans la situation (1), cela signifie que les diagrammes
\begin{equation}
\label{equation-action}
\vcenter{
\xymatrix{
G \times_S G \times_S X \ar[r]_-{1_G \times a} \ar[d]_{m \times 1_X} &
G \times_S X \ar[d]^a \\
G \times_S X \ar[r]^a & X
}
}
\quad\quad
\vcenter{
\xymatrix{
G \times_S X \ar[r]_-a & X \\
X\ar[u]^{e \times 1_X} \ar[ru]_{1_X}
}
}
\end{equation}
sont commutatifs. Dans la situation (2), cela signifie simplement que le diagramme
$$
\xymatrix{
G \times_S X \ar[r]_-{\text{id} \times \psi} \ar[d]_a &
G \times_S Y \ar[d]^a \\
X \ar[r]^\psi & Y
}
$$
est commutatif.

\begin{definition}
\label{definition-free-action}
Soient $S$, $G \to S$ et $X \to S$ comme dans la
définition \ref{definition-action-group-scheme}.
Soit $a : G \times_S X \to X$ une action de $G$ sur $X/S$.
Nous disons que l'action est {\it libre} si, pour tout schéma $T$ sur $S$,
l'action $a : G(T) \times X(T) \to X(T)$ est une action libre
du groupe $G(T)$ sur l'ensemble $X(T)$.
\end{definition}

\begin{lemma}
\label{lemma-free-action}
Dans la situation de la définition \ref{definition-free-action},
l'action $a$ est libre si et seulement si
$$
G \times_S X \to X \times_S X, \quad (g, x) \mapsto (a(g, x), x)
$$
est un monomorphisme.
\end{lemma}

\begin{proof}
Cela résulte immédiatement des définitions.
\end{proof}






\section{Torseurs et espaces homogènes principaux}
\label{section-principal-homogeneous}

\noindent
Dans
Cohomologie sur les sites, définition \ref{sites-cohomology-definition-torsor},
nous avons défini un torseur sous un faisceau de groupes sur un site.
Supposons que $\tau \in\allowbreak \{Zariski,\allowbreak \etale,\allowbreak smooth,\allowbreak syntomic,\allowbreak fppf\}$ soit une
topologie et que $(G, m)$ soit un schéma en groupes sur $S$. Puisque $\tau$ est plus fine que
la topologie canonique (voir
Descente, lemme \ref{descent-lemma-fpqc-universal-effective-epimorphisms}),
nous voyons que $\underline{G}$ (voir
Sites, définition \ref{sites-definition-representable-sheaf})
est un faisceau de groupes sur $(\Sch/S)_\tau$.
Nous savons donc déjà ce que signifie être un
torseur sous $\underline{G}$ sur $(\Sch/S)_\tau$. Une situation particulière
se présente lorsque ce faisceau est représentable. Dans les définitions suivantes,
nous définissons directement ce que signifie, pour le schéma qui le représente, être un
torseur sous $G$.

\begin{definition}
\label{definition-pseudo-torsor}
Soit $S$ un schéma.
Soit $(G, m)$ un schéma en groupes sur $S$.
Soit $X$ un schéma sur $S$, et soit
$a : G \times_S X \to X$ une action de $G$ sur $X$.
\begin{enumerate}
\item Nous disons que $X$ est un {\it pseudo-torseur sous $G$} ou que $X$ est
{\it formellement homogène principal sous $G$} si le
morphisme induit de schémas $G \times_S X \to X \times_S X$,
$(g, x) \mapsto (a(g, x), x)$ est un isomorphisme de schémas sur $S$.
\item {\emergencystretch=1em Un pseudo-torseur $X$ sous $G$ est dit {\it trivial} s'il existe
un isomorphisme $G$-équivariant $G \to X$ sur $S$, où $G$ agit sur
$G$ par multiplication à gauche.\par}
\end{enumerate}
\end{definition}

\noindent
Il est clair que, si $S' \to S$ est un morphisme de schémas, alors
le changement de base $X_{S'}$ d'un pseudo-torseur sous $G$ sur $S$ est un
pseudo-torseur sous $G_{S'}$ sur $S'$.

\begin{lemma}
\label{lemma-characterize-trivial-pseudo-torsors}
Dans la situation de la
définition \ref{definition-pseudo-torsor}.
\begin{enumerate}
\item Le schéma $X$ est un pseudo-torseur sous $G$ si et seulement si, pour tout schéma
$T$ sur $S$, l'ensemble $X(T)$ est vide ou l'action du groupe $G(T)$
sur $X(T)$ est simplement transitive.
\item Un pseudo-torseur $X$ sous $G$ est trivial si et seulement si le morphisme
$X \to S$ possède une section.
\end{enumerate}
\end{lemma}

\begin{proof}
Démonstration omise.
\end{proof}

\begin{definition}
\label{definition-principal-homogeneous-space}
Soit $S$ un schéma.
Soit $(G, m)$ un schéma en groupes sur $S$.
Soit $X$ un pseudo-torseur sous $G$ sur $S$.
\begin{enumerate}
\item Nous disons que $X$ est un {\it espace homogène principal}
ou un {\it torseur sous $G$} s'il existe un recouvrement fpqc\footnote{Cela signifie
que, par défaut, un torseur est un pseudo-torseur trivial sur un
recouvrement fpqc. C'est la définition de \cite[Expos\'e IV, 6.5]{SGA3}.
Elle est quelque peu malcommode pour nous, puisque nous travaillons le plus souvent pour la
topologie fppf.}
$\{S_i \to S\}_{i \in I}$ tel que chaque
$X_{S_i} \to S_i$ possède une section (c'est-à-dire qu'il est un pseudo-torseur trivial sous $G_{S_i}$).
\item Soit $\tau \in\allowbreak \{Zariski,\allowbreak \etale,\allowbreak smooth,\allowbreak syntomic,\allowbreak fppf\}$.
Nous disons que $X$ est un {\it torseur sous $G$ pour la topologie $\tau$}, ou un
{\it $\tau$-torseur sous $G$}, ou simplement un {\it $\tau$-torseur},
s'il existe un recouvrement $\tau$ $\{S_i \to S\}_{i \in I}$
tel que chaque $X_{S_i} \to S_i$ possède une section.
\item Si $X$ est un torseur sous $G$, nous disons qu'il est
{\it quasi-isotrivial} s'il est un torseur pour la topologie \'etale.
\item Si $X$ est un torseur sous $G$, nous disons qu'il est
{\it localement trivial} s'il est un torseur pour la topologie de Zariski.
\end{enumerate}
\end{definition}

\noindent
Nous disons parfois « soit $X$ un torseur sous $G$ sur $S$ » pour indiquer que
$X$ est un schéma sur $S$ muni d'une action de $G$ qui en fait
un espace homogène principal sur $S$.
Montrons maintenant que cela concorde avec la terminologie introduite plus haut
lorsque les deux s'appliquent.

\begin{lemma}
\label{lemma-torsor}
Soit $S$ un schéma.
Soit $(G, m)$ un schéma en groupes sur $S$.
Soit $X$ un schéma sur $S$, et soit
$a : G \times_S X \to X$ une action de $G$ sur $X$.
Soit $\tau \in\allowbreak \{Zariski,\allowbreak \etale,\allowbreak smooth,\allowbreak syntomic,\allowbreak fppf\}$.
Alors $X$ est un torseur sous $G$ pour la topologie $\tau$ si et seulement si
$\underline{X}$ est un torseur sous $\underline{G}$ sur $(\Sch/S)_\tau$.
\end{lemma}

\begin{proof}
Démonstration omise.
\end{proof}

\begin{remark}
\label{remark-fun-with-torsors}
Soit $(G, m)$ un schéma en groupes sur le schéma $S$.
Dans cette situation, les types naturels de questions suivants se posent :
\begin{enumerate}
\item Si $X \to S$ est un pseudo-torseur sous $G$ et si $X \to S$ est surjectif,
alors $X$ est-il nécessairement un torseur sous $G$ ?
\item Tout torseur sous $\underline{G}$ sur $(\Sch/S)_{fppf}$ est-il
représentable ? Autrement dit, tout torseur sous $\underline{G}$
provient-il d'un torseur fppf sous $G$ ?
\item Tout torseur sous $G$ est-il un
torseur fppf (resp.\ lisse, resp.\ \'etale, resp.\ de Zariski) ?
\end{enumerate}
En général, les réponses à ces questions sont négatives. Pour obtenir une réponse positive,
nous devons imposer des conditions supplémentaires à $G \to S$.
Par exemple :
si $S$ est le spectre d'un corps, alors la réponse à (1) est affirmative
car $\{X \to S\}$ est alors un recouvrement fpqc qui trivialise $X$.
Si $G \to S$ est affine, alors la réponse à (2) est affirmative
(cela résulte de Descente, lemme \ref{descent-lemma-affine}).
Si $G = \text{GL}_{n, S}$, alors la réponse à (3) est affirmative
et, en fait, tout torseur sous $\text{GL}_{n, S}$ est localement trivial
(cela résulte de
Descente, lemme \ref{descent-lemma-finite-locally-free-descends}).
\end{remark}



\section{Faisceaux quasi-cohérents équivariants}
\label{section-equivariant}

\noindent
Nous considérons les « fonctions » comme duales de l'« espace ». Ainsi, pour un morphisme d'espaces,
l'application induite sur les fonctions va dans le sens opposé. De plus, nous considérons les
sections d'un faisceau de modules comme des « fonctions ». Cela nous conduit naturellement
au sens des flèches choisi dans la définition suivante.

\begin{definition}
\label{definition-equivariant-module}
Soit $S$ un schéma, soit $(G, m)$ un schéma en groupes sur $S$, et
soit $a : G \times_S X \to X$ une action du schéma en groupes $G$
sur $X/S$. Un {\it module $G$-équivariant quasi-cohérent sur $\mathcal{O}_X$},
ou simplement un {\it module quasi-cohérent équivariant sur $\mathcal{O}_X$},
est une paire $(\mathcal{F}, \alpha)$, où $\mathcal{F}$ est un module quasi-cohérent
sur $\mathcal{O}_X$, et $\alpha$ est un morphisme de $\mathcal{O}_{G \times_S X}$-modules
donné par
$$
\alpha : a^*\mathcal{F} \longrightarrow \text{pr}_1^*\mathcal{F}
$$
où $\text{pr}_1 : G \times_S X \to X$ est la projection,
tel que
\begin{enumerate}
\item le diagramme
$$
\xymatrix{
(1_G \times a)^*\text{pr}_1^*\mathcal{F} \ar[r]_-{\text{pr}_{12}^*\alpha} &
\text{pr}_2^*\mathcal{F} \\
(1_G \times a)^*a^*\mathcal{F} \ar[u]^{(1_G \times a)^*\alpha} \ar@{=}[r] &
(m \times 1_X)^*a^*\mathcal{F} \ar[u]_{(m \times 1_X)^*\alpha}
}
$$
soit commutatif dans la catégorie des
$\mathcal{O}_{G \times_S G \times_S X}$-modules, et
\item le tiré en arrière
$$
(e \times 1_X)^*\alpha : \mathcal{F} \longrightarrow \mathcal{F}
$$
soit l'identité.
\end{enumerate}
Pour une explication, comparer avec les diagrammes correspondants de
l'équation (\ref{equation-action}).
\end{definition}

\noindent
Notons que la commutativité du premier diagramme garantit que
$(e \times 1_X)^*\alpha$ est un endomorphisme idempotent de $\mathcal{F}$,
de sorte que la condition (2) signifie simplement que c'est un isomorphisme.

\begin{lemma}
\label{lemma-pullback-equivariant}
Soit $S$ un schéma. Soit $G$ un schéma en groupes sur $S$.
Soit $f : Y \to X$ un morphisme $G$-équivariant entre des $S$-schémas
munis d'actions de $G$. La règle
$(\mathcal{F}, \alpha) \mapsto (f^*\mathcal{F}, (1_G \times f)^*\alpha)$
définit un foncteur de la catégorie des modules quasi-cohérents $G$-équivariants
sur $\mathcal{O}_X$ vers la catégorie des
modules quasi-cohérents $G$-équivariants sur $\mathcal{O}_Y$.
\end{lemma}

\begin{proof}
Démonstration omise.
\end{proof}

\noindent
Donnons un exemple.

\begin{example}
\label{example-Gm-on-affine}
Soit $A$ un anneau $\mathbf{Z}$-gradué, c'est-à-dire que $A$ est muni d'une décomposition en
somme directe $A = \bigoplus_{n \in \mathbf{Z}} A_n$ et que
$A_n \cdot A_m \subset A_{n + m}$.
Posons $X = \Spec(A)$. Nous obtenons alors une action de $\mathbf{G}_m$
$$
a : \mathbf{G}_m \times X \longrightarrow X
$$
à partir de l'homomorphisme d'anneaux $\mu : A \to A \otimes \mathbf{Z}[x, x^{-1}]$,
$f \mapsto f \otimes x^{\deg(f)}$.
En effet, pour le vérifier, nous devons vérifier que
$$
\xymatrix{
A \ar[r]_\mu \ar[d]_\mu &
A \otimes \mathbf{Z}[x, x^{-1}] \ar[d]^{\mu \otimes 1} \\
A \otimes \mathbf{Z}[x, x^{-1}] \ar[r]^-{1 \otimes m} &
A \otimes \mathbf{Z}[x, x^{-1}] \otimes \mathbf{Z}[x, x^{-1}]
}
$$
où $m(x) = x \otimes x$, voir l'exemple \ref{example-multiplicative-group}.
Cela est immédiatement clair en évaluant sur un élément homogène.
Supposons que $M$ soit un $A$-module gradué. Nous obtenons alors un
module quasi-cohérent $\mathbf{G}_m$-équivariant sur $\mathcal{O}_X$
$\mathcal{F} = \widetilde{M}$ en prenant $\alpha$ comme dans la
définition \ref{definition-equivariant-module}, correspondant au
morphisme de $A \otimes \mathbf{Z}[x, x^{-1}]$-modules
$$
M \otimes_{A, \mu} (A \otimes \mathbf{Z}[x, x^{-1}])
\longrightarrow
M \otimes_{A, \text{id}_A \otimes 1} (A \otimes \mathbf{Z}[x, x^{-1}])
$$
qui envoie $m \otimes 1 \otimes 1$ sur $m \otimes 1 \otimes x^{\deg(m)}$
pour $m \in M$ homogène.
\end{example}

\begin{lemma}
\label{lemma-complete-reducibility-Gm}
Soit $a : \mathbf{G}_m \times X \to X$ une action sur un schéma affine.
Alors $X$ est le spectre d'un anneau $\mathbf{Z}$-gradué
et l'action est celle de l'exemple \ref{example-Gm-on-affine}.
\end{lemma}

\begin{proof}
Soit $f \in A = \Gamma(X, \mathcal{O}_X)$. Alors nous pouvons écrire
$$
a^\sharp(f) = \sum\nolimits_{n \in \mathbf{Z}} f_n \otimes x^n
\quad\text{dans}\quad
A \otimes \mathbf{Z}[x, x^{-1}] =
\Gamma(\mathbf{G}_m \times X, \mathcal{O}_{\mathbf{G}_m \times X})
$$
sous la forme d'une somme finie, les $f_n$ de $A$ étant déterminés de façon unique.
Nous obtenons ainsi des applications $A \to A$, $f \mapsto f_n$.
Puisque $a$ est une action, en évaluant en $x = 1$,
nous obtenons $f = \sum f_n$. Puisque $a$ est une action,
nous trouvons que
$$
\sum (f_n)_m \otimes x^m \otimes x^n = \sum f_n x^n \otimes x^n
$$
(comparer avec le calcul de l'exemple \ref{example-Gm-on-affine}).
Ainsi $(f_n)_m = 0$ si $n \not = m$ et $(f_n)_n = f_n$.
Ainsi, si nous posons
$$
A_n = \{f \in A \mid f_n = f\}
$$
alors nous obtenons $A = \sum A_n$. En revanche, la somme est nécessairement
directe puisque $f = 0$ implique $f_n = 0$ dans la situation ci-dessus.
\end{proof}

\begin{lemma}
\label{lemma-Gm-equivariant-module}
Soit $A$ un anneau gradué. Soit $X = \Spec(A)$ muni de l'action
$a : \mathbf{G}_m \times X \to X$ de l'exemple \ref{example-Gm-on-affine}.
Soit $\mathcal{F}$ un module quasi-cohérent $\mathbf{G}_m$-équivariant
sur $\mathcal{O}_X$. Alors $M = \Gamma(X, \mathcal{F})$
possède une graduation canonique qui en fait un $A$-module gradué
et telle que l'isomorphisme $\widetilde{M} \to \mathcal{F}$
(Schémas, lemme \ref{schemes-lemma-quasi-coherent-affine})
soit un isomorphisme de modules $\mathbf{G}_m$-équivariants, où
la structure $\mathbf{G}_m$-équivariante sur $\widetilde{M}$
est celle de l'exemple \ref{example-Gm-on-affine}.
\end{lemma}

\begin{proof}
On peut soit démontrer ceci en répétant les arguments du
lemme \ref{lemma-complete-reducibility-Gm} pour le module $M$.
On peut aussi considérer le schéma
$(X', \mathcal{O}_{X'}) = (X, \mathcal{O}_X \oplus \mathcal{F})$,
où $\mathcal{F}$ est considéré comme un idéal de carré nul.
Il existe une action naturelle $a' : \mathbf{G}_m \times X' \to X'$
définie à l'aide de l'action sur $X$ et sur $\mathcal{F}$. Appliquons alors
le lemme \ref{lemma-complete-reducibility-Gm} à $X'$ et concluons.
(L'intérêt de cet argument est qu'il montre immédiatement
que les graduations de $A$ et de $M$ sont compatibles, c'est-à-dire que $M$
est un $A$-module gradué.)
Détails omis.
\end{proof}



\section{Groupoïdes}
\label{section-groupoids}

\noindent
Rappelons qu'un groupoïde est une catégorie dans laquelle tout morphisme
est un isomorphisme ; voir
Catégories, définition \ref{categories-definition-groupoid}.
Ainsi, un groupoïde possède un ensemble d'objets $\text{Ob}$,
un ensemble de flèches $\text{Arrows}$, des applications {\it source} et {\it but}
$s, t : \text{Arrows} \to \text{Ob}$, ainsi qu'une {\it loi de composition}
$c : \text{Arrows} \times_{s, \text{Ob}, t} \text{Arrows}
\to \text{Arrows}$.
Ces applications satisfont exactement aux axiomes suivants :
\begin{enumerate}
\item {\emergencystretch=2em (associativité) $c \circ (1, c) = c \circ (c, 1)$ comme applications
$\text{Arrows} \times_{s, \text{Ob}, t}
\text{Arrows} \times_{s, \text{Ob}, t}
\text{Arrows} \to \text{Arrows}$,\par}
\item (identité) il existe une application $e : \text{Ob} \to \text{Arrows}$
telle que
\begin{enumerate}
\item $s \circ e = t \circ e = \text{id}$ comme applications $\text{Ob} \to \text{Ob}$,
\item $c \circ (1, e \circ s) = c \circ (e \circ t, 1) = 1$
comme applications $\text{Arrows} \to \text{Arrows}$,
\end{enumerate}
\item (inverse) il existe une application $i : \text{Arrows} \to \text{Arrows}$
telle que
\begin{enumerate}
\item $s \circ i = t$, $t \circ i = s$ comme applications $\text{Arrows} \to \text{Ob}$,
et
\item $c \circ (1, i) = e \circ t$ et $c \circ (i, 1) = e \circ s$
comme applications $\text{Arrows} \to \text{Arrows}$.
\end{enumerate}
\end{enumerate}
Dans ce cas, les applications $e$ et $i$ sont déterminées de manière unique et
$i$ est une bijection. Notons que si $(\text{Ob}', \text{Arrows}', s', t', c')$
est une seconde catégorie groupoïde, alors un foncteur
$f : (\text{Ob}, \text{Arrows}, s, t, c) \to
(\text{Ob}', \text{Arrows}', s', t', c')$
est défini par un couple d'applications d'ensembles $f : \text{Ob} \to \text{Ob}'$ et
$f : \text{Arrows} \to \text{Arrows}'$ telles que
$s' \circ f = f \circ s$, $t' \circ f = f \circ t$ et
$c' \circ (f, f) = f \circ c$. La compatibilité avec l'identité et
l'inverse est automatique. Nous utiliserons ce fait ci-dessous.
(Attention : la compatibilité avec l'identité
doit être imposée dans le cas des catégories générales.)

\begin{definition}
\label{definition-groupoid}
Soit $S$ un schéma.
\begin{enumerate}
\item Un {\it schéma en groupoïdes sur $S$}, ou simplement un
{\it groupoïde sur $S$}, est un
quintuplet $(U, R, s, t, c)$, où
$U$ et $R$ sont des schémas sur $S$ et
$s, t : R \to U$ et $c : R \times_{s, U, t} R \to R$
sont des morphismes de schémas sur $S$ ayant la
propriété suivante : pour tout schéma
$T$ sur $S$, le quintuplet
$$
(U(T), R(T), s, t, c)
$$
est une catégorie groupoïde au sens décrit ci-dessus.
\item Un {\it morphisme
$f : (U, R, s, t, c) \to (U', R', s', t', c')$
de schémas en groupoïdes sur $S$} est défini par des morphismes
de schémas $f : U \to U'$ et $f : R \to R'$ ayant la
propriété suivante : pour tout schéma
$T$ sur $S$, les applications $f$ définissent un foncteur de la
catégorie groupoïde $(U(T), R(T), s, t, c)$ vers la
catégorie groupoïde $(U'(T), R'(T), s', t', c')$.
\end{enumerate}
\end{definition}

\noindent
Soit $(U, R, s, t, c)$ un groupoïde sur $S$.
Notons que, d'après les remarques précédant la définition et le lemme de Yoneda,
il existe d'uniques morphismes de schémas
$e : U \to R$ et
$i : R \to R$ sur $S$ tels que, pour tout schéma $T$ sur $S$,
l'application induite $e : U(T) \to R(T)$ soit l'application des identités et que
$i : R(T) \to R(T)$ soit l'application des inverses de
la catégorie groupoïde. Le septuplet $(U, R, s, t, c, e, i)$
satisfait aux diagrammes commutatifs correspondant à chacun des
axiomes (1), (2)(a), (2)(b), (3)(a) et (3)(b) ci-dessus et, réciproquement,
étant donné un septuplet ayant cette propriété, le quintuplet $(U, R, s, t, c)$
est un schéma en groupoïdes. Notons que $i$ est un isomorphisme
et que $e$ est une section à la fois de $s$ et de $t$.
En outre, étant donné un schéma en groupoïdes sur $S$, on note
$$
j = (t, s) : R \longrightarrow U \times_S U
$$
ce qui est compatible avec nos conventions de la section
\ref{section-equivalence-relations} ci-dessus.
On dit parfois « soit $(U, R, s, t, c, e, i)$ un
groupoïde sur $S$ » pour souligner l'existence des identités et des
inverses.

\begin{lemma}
\label{lemma-groupoid-pre-equivalence}
Étant donné un schéma en groupoïdes $(U, R, s, t, c)$ sur $S$,
le morphisme $j : R \to U \times_S U$ est une prérelation
d'équivalence.
\end{lemma}

\begin{proof}
La preuve est omise.
C'est un joli exercice sur les définitions.
\end{proof}

\begin{lemma}
\label{lemma-equivalence-groupoid}
Étant donné une relation d'équivalence $j : R \to U \times_S U$ sur $S$,
il existe une unique manière de l'étendre en un groupoïde
$(U, R, s, t, c)$ sur $S$.
\end{lemma}

\begin{proof}
La preuve est omise.
C'est un joli exercice sur les définitions.
\end{proof}

\begin{lemma}
\label{lemma-diagram}
Soit $S$ un schéma.
Soit $(U, R, s, t, c)$ un groupoïde sur $S$.
Dans le diagramme commutatif
$$
\xymatrix{
& U & \\
R \ar[d]_s \ar[ru]^t &
R \times_{s, U, t} R
\ar[l]^-{\text{pr}_0} \ar[d]^{\text{pr}_1} \ar[r]_-c &
R \ar[d]^s \ar[lu]_t \\
U & R \ar[l]_t \ar[r]^s & U
}
$$
les deux carrés inférieurs sont des carrés cartésiens.
De plus, le triangle supérieur (qui est en réalité un carré)
est lui aussi cartésien.
\end{lemma}

\begin{proof}
La preuve est omise.
C'est un exercice sur les définitions et le point de
vue fonctoriel en géométrie algébrique.
\end{proof}

\begin{lemma}
\label{lemma-diagram-pull}
Soit $S$ un schéma.
Soit $(U, R, s, t, c, e, i)$ un groupoïde sur $S$.
Le diagramme
\begin{equation}
\label{equation-pull}
\xymatrix{
R \times_{t, U, t} R
\ar@<1ex>[r]^-{\text{pr}_1} \ar@<-1ex>[r]_-{\text{pr}_0}
\ar[d]_{(\text{pr}_0, c \circ (i, 1))} &
R \ar[r]^t \ar[d]^{\text{id}_R} &
U \ar[d]^{\text{id}_U} \\
R \times_{s, U, t} R
\ar@<1ex>[r]^-c \ar@<-1ex>[r]_-{\text{pr}_0} \ar[d]_{\text{pr}_1} &
R \ar[r]^t \ar[d]^s &
U \\
R \ar@<1ex>[r]^s \ar@<-1ex>[r]_t &
U
}
\end{equation}
est commutatif. Les deux lignes supérieures sont isomorphes au moyen des applications verticales indiquées.
Les deux carrés inférieurs de gauche sont cartésiens.
\end{lemma}

\begin{proof}
La commutativité du diagramme résulte des axiomes d'un groupoïde.
Notons qu'en termes de groupoïdes, la flèche verticale supérieure gauche associe
à une paire de morphismes $(\alpha, \beta)$ de même but la paire
de morphismes $(\alpha, \alpha^{-1} \circ \beta)$. Dans tout groupoïde,
ceci définit une bijection entre
$\text{Arrows} \times_{t, \text{Ob}, t} \text{Arrows}$
et
$\text{Arrows} \times_{s, \text{Ob}, t} \text{Arrows}$. D'où la deuxième
assertion du lemme.
La dernière assertion résulte du lemme \ref{lemma-diagram}.
\end{proof}

\begin{lemma}
\label{lemma-base-change-groupoid}
Soit $(U, R, s, t, c)$ un groupoïde sur le schéma $S$.
Soit $S' \to S$ un morphisme. Alors les changements de base $U' = S' \times_S U$,
$R' = S' \times_S R$, munis des changements de base $s'$, $t'$, $c'$
des morphismes $s, t, c$, forment un schéma en groupoïdes
$(U', R', s', t', c')$ sur $S'$, et les projections
déterminent un morphisme
$(U', R', s', t', c') \to (U, R, s, t, c)$
de schémas en groupoïdes sur $S$.
\end{lemma}

\begin{proof}
La preuve est omise. Indication :
$R' \times_{s', U', t'} R' = S' \times_S (R \times_{s, U, t} R)$.
\end{proof}









\section{Faisceaux quasi-cohérents sur les groupoïdes}
\label{section-groupoids-quasi-coherent}

\noindent
Voir l'introduction de la section \ref{section-equivariant} pour nos
choix concernant le sens des flèches.

\begin{definition}
\label{definition-groupoid-module}
Soit $S$ un schéma et soit $(U, R, s, t, c)$ un schéma en groupoïdes sur $S$.
Un {\it module quasi-cohérent sur $(U, R, s, t, c)$}
est une paire $(\mathcal{F}, \alpha)$, où $\mathcal{F}$ est un
$\mathcal{O}_U$-module quasi-cohérent, et $\alpha$ est un homomorphisme de $\mathcal{O}_R$-modules
donné par
$$
\alpha : t^*\mathcal{F} \longrightarrow s^*\mathcal{F}
$$
tel que
\begin{enumerate}
\item le diagramme
$$
\xymatrix{
& \text{pr}_1^*t^*\mathcal{F} \ar[r]_-{\text{pr}_1^*\alpha} &
\text{pr}_1^*s^*\mathcal{F} \ar@{=}[rd] & \\
\text{pr}_0^*s^*\mathcal{F} \ar@{=}[ru] & & & c^*s^*\mathcal{F} \\
& \text{pr}_0^*t^*\mathcal{F} \ar[lu]^{\text{pr}_0^*\alpha} \ar@{=}[r] &
c^*t^*\mathcal{F} \ar[ru]_{c^*\alpha}
}
$$
soit commutatif dans la catégorie des
$\mathcal{O}_{R \times_{s, U, t} R}$-modules, et
\item l'image inverse
$$
e^*\alpha : \mathcal{F} \longrightarrow \mathcal{F}
$$
soit l'identité.
\end{enumerate}
Comparer avec les diagrammes commutatifs du lemme \ref{lemma-diagram}.
\end{definition}

\noindent
La commutativité du premier diagramme force l'opérateur $e^*\alpha$
à être idempotent. La seconde condition peut donc être reformulée en disant
que $e^*\alpha$ est un isomorphisme. En fait, cette condition implique que
$\alpha$ est un isomorphisme.

\begin{lemma}
\label{lemma-isomorphism}
Soit $S$ un schéma et soit $(U, R, s, t, c)$ un schéma en groupoïdes sur $S$.
Si $(\mathcal{F}, \alpha)$ est un module quasi-cohérent sur $(U, R, s, t, c)$,
alors $\alpha$ est un isomorphisme.
\end{lemma}

\begin{proof}
Prenons l'image inverse du diagramme commutatif de la
définition \ref{definition-groupoid-module}
par le morphisme $(i, 1) : R \to R \times_{s, U, t} R$.
On voit alors que $i^*\alpha \circ \alpha = s^*e^*\alpha$.
En prenant l'image inverse par le morphisme $(1, i)$, on obtient la relation
$\alpha \circ i^*\alpha = t^*e^*\alpha$. D'après la seconde hypothèse,
ces morphismes sont l'identité. Ainsi, $i^*\alpha$ est un inverse de
$\alpha$.
\end{proof}

\begin{lemma}
\label{lemma-pullback}
Soit $S$ un schéma. Considérons un morphisme
$f : (U, R, s, t, c) \to (U', R', s', t', c')$
de schémas en groupoïdes sur $S$. Alors l'image inverse $f^*$ donnée par
$$
(\mathcal{F}, \alpha) \mapsto (f^*\mathcal{F}, f^*\alpha)
$$
définit un foncteur de la catégorie des faisceaux quasi-cohérents sur
$(U', R', s', t', c')$ vers la catégorie des faisceaux quasi-cohérents sur
$(U, R, s, t, c)$.
\end{lemma}

\begin{proof}
La preuve est omise.
\end{proof}

\begin{lemma}
\label{lemma-pushforward}
Soit $S$ un schéma. Considérons un morphisme
$f : (U, R, s, t, c) \to (U', R', s', t', c')$
de schémas en groupoïdes sur $S$. Supposons que
\begin{enumerate}
\item $f : U \to U'$ soit quasi-compact et quasi-séparé,
\item le carré
$$
\xymatrix{
R \ar[d]_t \ar[r]_f & R' \ar[d]^{t'} \\
U \ar[r]^f & U'
}
$$
soit cartésien, et
\item $s'$ et $t'$ soient plats.
\end{enumerate}
Alors l'image directe $f_*$ donnée par
$$
(\mathcal{F}, \alpha) \mapsto (f_*\mathcal{F}, f_*\alpha)
$$
définit un foncteur de la catégorie des faisceaux quasi-cohérents sur
$(U, R, s, t, c)$ vers la catégorie des faisceaux quasi-cohérents sur
$(U', R', s', t', c')$, adjoint à droite de l'image inverse définie dans le
lemme \ref{lemma-pullback}.
\end{lemma}

\begin{proof}
Puisque $U \to U'$ est quasi-compact et quasi-séparé, on voit que
$f_*$ transforme les faisceaux quasi-cohérents en faisceaux quasi-cohérents
(Schémas, lemme \ref{schemes-lemma-push-forward-quasi-coherent}).
De plus, puisque les carrés
$$
\vcenter{
\xymatrix{
R \ar[d]_t \ar[r]_f & R' \ar[d]^{t'} \\
U \ar[r]^f & U'
}
}
\quad\text{et}\quad
\vcenter{
\xymatrix{
R \ar[d]_s \ar[r]_f & R' \ar[d]^{s'} \\
U \ar[r]^f & U'
}
}
$$
sont cartésiens, on obtient $(t')^*f_*\mathcal{F} = f_*t^*\mathcal{F}$
et $(s')^*f_*\mathcal{F} = f_*s^*\mathcal{F}$, voir
Cohomologie des schémas, lemme
\ref{coherent-lemma-flat-base-change-cohomology}.
Il est donc légitime de considérer $f_*\alpha$ comme un homomorphisme
$(t')^*f_*\mathcal{F} \to (s')^*f_*\mathcal{F}$. Un argument semblable
montre que $f_*\alpha$ satisfait à la condition de cocycle.
Ce foncteur est adjoint au foncteur image inverse, car les foncteurs image inverse
et image directe sur les modules d'espaces annelés sont adjoints.
Certains détails sont omis.
\end{proof}

\begin{lemma}
\label{lemma-colimits}
Soit $S$ un schéma. Soit $(U, R, s, t, c)$ un schéma en groupoïdes sur $S$.
La catégorie des modules quasi-cohérents sur $(U, R, s, t, c)$ admet des colimites.
\end{lemma}

\begin{proof}
Soit $i \mapsto (\mathcal{F}_i, \alpha_i)$ un diagramme sur la catégorie
d'indices $\mathcal{I}$. On peut former la colimite
$\mathcal{F} = \colim \mathcal{F}_i$,
qui est un faisceau quasi-cohérent sur $U$, voir
Schémas, section \ref{schemes-section-quasi-coherent}.
Puisque l'image inverse commute aux colimites, on a
$s^*\mathcal{F} = \colim s^*\mathcal{F}_i$ et, de même,
$t^*\mathcal{F} = \colim t^*\mathcal{F}_i$. On peut donc poser
$\alpha = \colim \alpha_i$. Nous omettons la preuve que $(\mathcal{F}, \alpha)$
est la colimite du diagramme dans la catégorie des modules quasi-cohérents
sur $(U, R, s, t, c)$.
\end{proof}

\begin{lemma}
\label{lemma-abelian}
Soit $S$ un schéma.
Soit $(U, R, s, t, c)$ un schéma en groupoïdes sur $S$.
Si $s$ et $t$ sont plats, alors la catégorie des modules quasi-cohérents sur
$(U, R, s, t, c)$ est abélienne.
\end{lemma}

\begin{proof}
Soit $\varphi : (\mathcal{F}, \alpha) \to (\mathcal{G}, \beta)$ un
homomorphisme de modules quasi-cohérents sur $(U, R, s, t, c)$. Puisque
$s$ est plat, on voit que
$$
0 \to s^*\Ker(\varphi)
\to s^*\mathcal{F} \to s^*\mathcal{G} \to s^*\Coker(\varphi) \to 0
$$
est exacte, de même que son analogue obtenu après image inverse par $t$. Ainsi, $\alpha$ et $\beta$
induisent des isomorphismes
$\kappa : t^*\Ker(\varphi) \to s^*\Ker(\varphi)$ et
$\lambda : t^*\Coker(\varphi) \to s^*\Coker(\varphi)$
qui satisfont à la condition de cocycle. Il est alors immédiat de
vérifier que $(\Ker(\varphi), \kappa)$ et
$(\Coker(\varphi), \lambda)$ sont respectivement un noyau et un conoyau dans la
catégorie des modules quasi-cohérents sur $(U, R, s, t, c)$. En outre,
la relation $\Coim(\varphi) = \Im(\varphi)$ résulte
du fait qu'elle est satisfaite sur $U$.
\end{proof}













\section{Colimites de modules quasi-cohérents}
\label{section-colimits}

\noindent
Dans cette section, nous démontrons quelques résultats techniques affirmant que, sous
des hypothèses convenables, tout module quasi-cohérent sur un groupoïde est
une colimite filtrante de « petits » modules quasi-cohérents.

\begin{lemma}
\label{lemma-construct-quasi-coherent}
Soit $(U, R, s, t, c)$ un schéma en groupoïdes sur $S$.
Supposons que $s, t$ soient plats, quasi-compacts et quasi-séparés.
Pour tout module quasi-cohérent $\mathcal{G}$ sur $U$, il existe
un isomorphisme canonique
$\alpha : t^*s_*t^*\mathcal{G} \to s^*s_*t^*\mathcal{G}$
qui fait de $(s_*t^*\mathcal{G}, \alpha)$ un module quasi-cohérent
sur $(U, R, s, t, c)$. Cette construction définit un foncteur
$$
\QCoh(\mathcal{O}_U) \longrightarrow \QCoh(U, R, s, t, c)
$$
qui est adjoint à droite du foncteur d'oubli
$(\mathcal{F}, \beta) \mapsto \mathcal{F}$.
\end{lemma}

\begin{proof}
L'image directe d'un module quasi-cohérent par un morphisme quasi-compact et
quasi-séparé est quasi-cohérente, voir Schémas, lemme
\ref{schemes-lemma-push-forward-quasi-coherent}. Ainsi, $s_*t^*\mathcal{G}$
est quasi-cohérent. Avec les notations du lemme \ref{lemma-diagram}, on a
$$
t^*s_*t^*\mathcal{G} =
\text{pr}_{1, *}\text{pr}_0^*t^*\mathcal{G} =
\text{pr}_{1, *}c^*t^*\mathcal{G} =
s^*s_*t^*\mathcal{G}
$$
L'égalité du milieu vient de ce que $t \circ c = t \circ \text{pr}_0$ comme
morphismes $R \times_{s, U, t} R \to U$, tandis que les première et dernière
égalités viennent de ce que l'on sait que le changement de base et l'image directe commutent
à ces étapes, d'après Cohomologie des schémas, lemme
\ref{coherent-lemma-flat-base-change-cohomology}.

\medskip\noindent
Pour vérifier la condition de cocycle de la définition \ref{definition-groupoid-module}
pour $\alpha$, ainsi que la propriété d'adjonction, décrivons la construction
$\mathcal{G} \mapsto (s_*t^*\mathcal{G}, \alpha)$ d'une autre manière.
Considérons le schéma en groupoïdes
$(R, R \times_{t, U, t} R, \text{pr}_0, \text{pr}_1, \text{pr}_{02})$
associé à la relation d'équivalence $R \times_{t, U, t} R$
sur $R$, voir le lemme \ref{lemma-equivalence-groupoid}.
Il existe un morphisme
$$
f :
(R, R \times_{t, U, t} R, \text{pr}_1, \text{pr}_0, \text{pr}_{02})
\longrightarrow
(U, R, s, t, c)
$$
de schémas en groupoïdes donné par $s : R \to U$ et par l'application $R \times_{t, U, t} R \to R$
qui est définie par $(r_0, r_1) \mapsto r_0^{-1} \circ r_1$ ; nous omettons la vérification
de la commutativité des diagrammes voulus. Puisque
$t, s : R \to U$ sont quasi-compacts, quasi-séparés et plats,
et puisque l'on a un carré cartésien
$$
\xymatrix{
R \times_{t, U, t} R \ar[d]_{\text{pr}_0}
\ar[rr]_-{(r_0, r_1) \mapsto r_0^{-1} \circ r_1} & & R \ar[d]^t \\
R \ar[rr]^s & & U
}
$$
Le lemme \ref{lemma-diagram-pull} montre que ce carré est cartésien ; par conséquent, le
lemme \ref{lemma-pushforward} s'applique à $f$.
Ainsi, les foncteurs image directe et image inverse des modules quasi-cohérents associés à
$f$ sont adjoints. Pour terminer la preuve, nous allons identifier ces foncteurs
aux foncteurs décrits ci-dessus. Pour cela, notons que
$$
t^* :
\QCoh(\mathcal{O}_U)
\longrightarrow
\QCoh(R, R \times_{t, U, t} R, \text{pr}_1, \text{pr}_0, \text{pr}_{02})
$$
est une équivalence d'après la théorie de la descente des faisceaux quasi-cohérents, car
$\{t : R \to U\}$ est un recouvrement fpqc, voir
Descente, proposition \ref{descent-proposition-fpqc-descent-quasi-coherent}.

\medskip\noindent
L'image directe par $f$, précomposée par l'équivalence $t^*$, envoie
$\mathcal{G}$ sur $(s_*t^*\mathcal{G}, \alpha)$ ;
nous omettons la vérification que l'isomorphisme $\alpha$
obtenu de cette manière est le même que celui construit ci-dessus.

\medskip\noindent
L'image inverse par $f$, postcomposée par l'inverse de l'équivalence $t^*$,
envoie $(\mathcal{F}, \beta)$ sur la descente relativement à $\{t : R \to U\}$
du module $s^*\mathcal{F}$ muni de la donnée de descente $\gamma$ sur
$R \times_{t, U, t} R$, cette donnée étant l'image inverse de $\beta$ par
$(r_0, r_1) \mapsto r_0^{-1} \circ r_1$.
Considérons l'isomorphisme $\beta : t^*\mathcal{F} \to s^*\mathcal{F}$.
La donnée de descente canonique (Descente, définition
\ref{descent-definition-descent-datum-effective-quasi-coherent})
sur $t^*\mathcal{F}$ relativement à $\{t : R \to U\}$
se traduit, via $\beta$, par l'application
$$
\text{pr}_0^*s^*\mathcal{F}
\xrightarrow{\text{pr}_0^*\beta^{-1}}
\text{pr}_0^*t^*\mathcal{F}
\xrightarrow{can}
\text{pr}_1^*t^*\mathcal{F}
\xrightarrow{\text{pr}_1^*\beta}
\text{pr}_1^*s^*\mathcal{F}
$$
Puisque $\beta$ satisfait à la condition de cocycle, cette application est égale à l'image inverse
de $\beta$ par $(r_0, r_1) \mapsto r_0^{-1} \circ r_1$. Pour le voir,
prenons la relation de cocycle proprement dite dans la définition \ref{definition-groupoid-module}
et prenons son image inverse par le morphisme
$(\text{pr}_0, c \circ (i, 1)) : R \times_{t, U, t} R \to R \times_{s, U, t} R$
qui intervient également dans le diagramme commutatif du
lemme \ref{lemma-diagram-pull}. Il en résulte que
$(s^*\mathcal{F}, \gamma)$ est isomorphe à $(t^*\mathcal{F}, can)$.
En définitive, nous concluons que l'image inverse par $f$, postcomposée par
l'inverse de l'équivalence $t^*$, est isomorphe au foncteur d'oubli
$(\mathcal{F}, \beta) \mapsto \mathcal{F}$.
\end{proof}

\begin{remark}
\label{remark-adjunction-map}
Dans la situation du lemme \ref{lemma-construct-quasi-coherent}, notons
$$
F : \QCoh(U, R, s, t, c) \to \QCoh(\mathcal{O}_U),\quad
(\mathcal{F}, \beta) \mapsto \mathcal{F}
$$
le foncteur d'oubli, et notons
$$
G : \QCoh(\mathcal{O}_U) \to \QCoh(U, R, s, t, c),\quad
\mathcal{G} \mapsto (s_*t^*\mathcal{G}, \alpha)
$$
l'adjoint à droite construit dans le lemme. Alors l'unité
$\eta : \text{id} \to G \circ F$ de l'adjonction, évaluée
en $(\mathcal{F}, \beta)$, est donnée par l'application
$$
\mathcal{F} \to s_*s^*\mathcal{F} \xrightarrow{\beta^{-1}} s_*t^*\mathcal{F}
$$
Nous omettons la vérification.
\end{remark}

\begin{lemma}
\label{lemma-push-pull}
Soit $f : Y \to X$ un morphisme de schémas. Soit $\mathcal{F}$
un $\mathcal{O}_X$-module quasi-cohérent, soit $\mathcal{G}$
un $\mathcal{O}_Y$-module quasi-cohérent et soit
$\varphi : \mathcal{G} \to f^*\mathcal{F}$ un homomorphisme de modules. Supposons que
\begin{enumerate}
\item $\varphi$ soit injectif,
\item $f$ soit quasi-compact, quasi-séparé, plat et surjectif,
\item $X$ et $Y$ soient localement noethériens, et
\item $\mathcal{G}$ soit un $\mathcal{O}_Y$-module cohérent.
\end{enumerate}
Alors $\mathcal{F} \cap f_*\mathcal{G}$, défini comme le produit fibré
$$
\xymatrix{
\mathcal{F} \ar[r] & f_*f^*\mathcal{F} \\
\mathcal{F} \cap f_*\mathcal{G} \ar[u] \ar[r] &
f_*\mathcal{G} \ar[u]
}
$$
est un $\mathcal{O}_X$-module cohérent.
\end{lemma}

\begin{proof}
Nous utiliserons librement la caractérisation des modules cohérents de
Cohomologie des schémas, lemme \ref{coherent-lemma-coherent-Noetherian},
ainsi que le fait que les modules cohérents forment une sous-catégorie de Serre
de $\QCoh(\mathcal{O}_X)$, voir
Cohomologie des schémas,
lemme \ref{coherent-lemma-coherent-Noetherian-quasi-coherent-sub-quotient}.
Si $f$ possède une section $\sigma$, on voit que
$\mathcal{F} \cap f_*\mathcal{G}$ est contenu dans l'image de
$\sigma^*\mathcal{G} \to \sigma^*f^*\mathcal{F} = \mathcal{F}$,
donc est cohérent. En général, pour montrer que $\mathcal{F} \cap f_*\mathcal{G}$
est cohérent, il suffit de montrer que
$f^*(\mathcal{F} \cap f_*\mathcal{G})$ est cohérent (voir
Descente, lemme \ref{descent-lemma-finite-type-descends}).
Puisque $f$ est plat, ce module est égal à $f^*\mathcal{F} \cap f^*f_*\mathcal{G}$.
Puisque $f$ est plat, quasi-compact et quasi-séparé, on a
$f^*f_*\mathcal{G} = p_*q^*\mathcal{G}$, où $p, q : Y \times_X Y \to Y$
sont les projections, voir
Cohomologie des schémas, lemme \ref{coherent-lemma-flat-base-change-cohomology}.
Puisque $p$ possède une section, on conclut.
\end{proof}

\noindent
Soit $S$ un schéma. Soit $(U, R, s, t, c)$ un schéma en groupoïdes sur $S$.
Supposons que $U$ soit localement noethérien. Dans le lemme ci-dessous, un
faisceau quasi-cohérent $(\mathcal{F}, \alpha)$ sur $(U, R, s, t, c)$ est dit
{\it cohérent} si $\mathcal{F}$ est un $\mathcal{O}_U$-module cohérent.

\begin{lemma}
\label{lemma-colimit-coherent}
Soit $(U, R, s, t, c)$ un schéma en groupoïdes sur $S$.
Supposons que
\begin{enumerate}
\item $U$, $R$ soient noethériens,
\item $s, t$ soient plats, quasi-compacts et quasi-séparés.
\end{enumerate}
Alors tout module quasi-cohérent $(\mathcal{F}, \beta)$ sur $(U, R, s, t, c)$
est une colimite filtrante de modules cohérents.
\end{lemma}

\begin{proof}
Nous utiliserons sans plus ample mention la caractérisation donnée dans
Cohomologie des schémas, lemme \ref{coherent-lemma-coherent-Noetherian}, des
modules cohérents sur un schéma localement noethérien. On peut écrire
$\mathcal{F} = \colim \mathcal{H}_i$ comme colimite filtrante
de sous-modules cohérents $\mathcal{H}_i \subset \mathcal{F}$, voir
Cohomologie des schémas, lemme \ref{coherent-lemma-directed-colimit-coherent}.
Étant donné un faisceau quasi-cohérent $\mathcal{H}$ sur $U$, on note
$(s_*t^*\mathcal{H}, \alpha)$ le faisceau quasi-cohérent sur $(U, R, s, t, c)$ du
lemme \ref{lemma-construct-quasi-coherent}. Considérons le morphisme d'adjonction
$(\mathcal{F}, \beta) \to (s_*t^*\mathcal{F}, \alpha)$ dans
$\QCoh(U, R, s, t, c)$, voir remarque \ref{remark-adjunction-map}.
Posons
$$
(\mathcal{F}_i, \beta_i) =
(\mathcal{F}, \beta)
\times_{(s_*t^*\mathcal{F}, \alpha)}
(s_*t^*\mathcal{H}_i, \alpha)
$$
{\emergencystretch=1em dans $\QCoh(U,\allowbreak R,\allowbreak s,\allowbreak t,\allowbreak c)$. Puisque la restriction à $U$ est un foncteur exact
sur $\QCoh(U,\allowbreak R,\allowbreak s,\allowbreak t,\allowbreak c)$ d'après la démonstration du lemme \ref{lemma-abelian},
on obtient un diagramme cartésien\par}
$$
\xymatrix{
\mathcal{F} \ar[r] & s_*t^*\mathcal{F} \\
\mathcal{F}_i \ar[r] \ar[u] & s_*t^*\mathcal{H}_i \ar[u]
}
$$
autrement dit $\mathcal{F}_i = \mathcal{F} \cap s_*t^*\mathcal{H}_i$.
D'après la description du morphisme d'adjonction dans la remarque \ref{remark-adjunction-map},
ce diagramme est isomorphe au diagramme
$$
\xymatrix{
\mathcal{F} \ar[r] & s_*s^*\mathcal{F} \\
\mathcal{F}_i \ar[r] \ar[u] & s_*t^*\mathcal{H}_i \ar[u]
}
$$
où la flèche verticale de droite s'obtient en appliquant $s_*$ au morphisme
$$
t^*\mathcal{H}_i \to t^*\mathcal{F} \xrightarrow{\beta} s^*\mathcal{F}
$$
Cette flèche est injective puisque $t$ est un morphisme plat. Il s'ensuit que
$\mathcal{F}_i$ est cohérent par le lemme \ref{lemma-push-pull}.
Enfin, puisque $s$ est quasi-compact et quasi-séparé, on voit que
$s_*$ commute aux colimites (voir
Cohomologie des schémas, lemme \ref{coherent-lemma-colimit-cohomology}).
Ainsi $s_*t^*\mathcal{F} = \colim s_*t^*\mathcal{H}_i$ et donc
$(\mathcal{F}, \beta) = \colim (\mathcal{F}_i, \beta_i)$, comme voulu.
\end{proof}

\noindent
Voici un lemme curieux, utile lorsqu'on travaille avec des groupoïdes
sur des corps. C'est en fait l'argument usuel pour démontrer que toute
représentation d'un groupe algébrique est une colimite de représentations
de dimension finie.

\begin{lemma}
\label{lemma-colimit-finite-type}
Soit $(U, R, s, t, c)$ un schéma en groupoïdes sur $S$.
Supposons que
\begin{enumerate}
\item $U$, $R$ soient affines,
\item il existe des $e_i \in \mathcal{O}_R(R)$ tels que
tout élément $g \in \mathcal{O}_R(R)$ s'écrive de manière unique sous la forme
$\sum s^*(f_i)e_i$ avec des $f_i \in \mathcal{O}_U(U)$.
\end{enumerate}
Alors tout module quasi-cohérent $(\mathcal{F}, \alpha)$ sur $(U, R, s, t, c)$
est une colimite filtrante de modules quasi-cohérents de type fini.
\end{lemma}

\begin{proof}
{\emergencystretch=1em L'hypothèse signifie que $\mathcal{O}_R(R)$ est un
$\mathcal{O}_U(U)$-module libre via $s$, de base $e_i$. Ainsi,
pour tout $\mathcal{O}_U$-module quasi-cohérent $\mathcal{G}$,
on a $s^*\mathcal{G}(R) = \bigoplus_i \mathcal{G}(U)e_i$.
On écrira $s(-)$ pour désigner l'image réciproque des sections par $s$, et
de même pour les autres morphismes.
Soit $(\mathcal{F}, \alpha)$ un module quasi-cohérent sur
$(U, R, s, t, c)$. Soit $\sigma \in \mathcal{F}(U)$. D'après ce qui précède,
on peut écrire\par}
$$
\alpha(t(\sigma)) = \sum s(\sigma_i) e_i
$$
pour des $\sigma_i \in \mathcal{F}(U)$ uniques (tous sauf un nombre fini
sont bien entendu nuls). On peut aussi écrire
$$
c(e_i) = \sum \text{pr}_1(f_{ij}) \text{pr}_0(e_j)
$$
en tant que fonctions sur $R \times_{s, U, t}R$. Alors la commutativité du diagramme
de la définition \ref{definition-groupoid-module} signifie que
$$
\sum \text{pr}_1(\alpha(t(\sigma_i))) \text{pr}_0(e_i)
=
\sum \text{pr}_1(s(\sigma_i)f_{ij}) \text{pr}_0(e_j)
$$
(calcul omis). En identifiant les coefficients de $\text{pr}_0(e_l)$,
on voit que $\alpha(t(\sigma_l)) = \sum s(\sigma_i)f_{il}$. Ainsi,
le sous-module $\mathcal{G} \subset \mathcal{F}$ engendré par les
éléments $\sigma_i$ définit un module quasi-cohérent de type fini
préservé par $\alpha$. C'est donc un sous-objet de $\mathcal{F}$ dans
$\QCoh(U, R, s, t, c)$. Ce sous-module contient $\sigma$
(comme on le voit en prenant l'image réciproque de la première relation par $e$). D'où le résultat.
\end{proof}

\noindent
Nous suggérons au lecteur de sauter le reste de cette section. Soit $S$ un schéma.
Soit $(U, R, s, t, c)$ un schéma en groupoïdes sur $S$. Soit $\kappa$ un
cardinal. Dans la suite, on dira qu'un faisceau quasi-cohérent
$(\mathcal{F}, \alpha)$ sur $(U, R, s, t, c)$ est $\kappa$-engendré si
$\mathcal{F}$ est $\kappa$-engendré comme $\mathcal{O}_U$-module, voir
Propriétés, définition \ref{properties-definition-kappa-generated}.

\begin{lemma}
\label{lemma-set-of-iso-classes}
Soit $(U, R, s, t, c)$ un schéma en groupoïdes sur $S$.
Soit $\kappa$ un cardinal.
Il existe un ensemble $T$ et une famille $(\mathcal{F}_t, \alpha_t)_{t \in T}$ de
modules quasi-cohérents $\kappa$-engendrés sur $(U, R, s, t, c)$
telle que tout module quasi-cohérent $\kappa$-engendré sur
$(U, R, s, t, c)$ soit isomorphe à l'un des $(\mathcal{F}_t, \alpha_t)$.
\end{lemma}

\begin{proof}
Pour tout module quasi-cohérent $\mathcal{F}$ sur $U$, il existe un
ensemble (éventuellement vide) de morphismes $\alpha : t^*\mathcal{F} \to s^*\mathcal{F}$
tels que $(\mathcal{F}, \alpha)$ soit un module quasi-cohérent sur
$(U, R, s, t, c)$. D'après
Propriétés, lemme \ref{properties-lemma-set-of-iso-classes},
il existe un ensemble de classes d'isomorphisme d'objets $\kappa$-engendrés qui sont des
$\mathcal{O}_U$-modules quasi-cohérents.
\end{proof}

\begin{lemma}
\label{lemma-colimit-kappa}
Soit $(U, R, s, t, c)$ un schéma en groupoïdes sur $S$.
Supposons que $s, t$ soient plats. Il existe un
cardinal $\kappa$ tel que tout module quasi-cohérent
$(\mathcal{F}, \alpha)$ sur $(U, R, s, t, c)$
soit la colimite filtrante de ses sous-modules $\kappa$-engendrés
quasi-cohérents.
\end{lemma}

\begin{proof}
Dans l'énoncé du lemme et dans cette démonstration,
un {\it sous-module} d'un module quasi-cohérent $(\mathcal{F}, \alpha)$
est un sous-module quasi-cohérent $\mathcal{G} \subset \mathcal{F}$
tel que $\alpha(t^*\mathcal{G}) = s^*\mathcal{G}$ comme sous-faisceaux de
$s^*\mathcal{F}$. Cela a un sens car, $s, t$ étant plats, les
images réciproques $s^*$ et $t^*$ sont exactes, c'est-à-dire préservent les sous-faisceaux.
La démonstration reprendra celle de
Propriétés, lemme \ref{properties-lemma-colimit-kappa}.
Nous invitons vivement le lecteur à lire d'abord cette démonstration.

\medskip\noindent
Choisissons un recouvrement ouvert affine $U = \bigcup_{i \in I} U_i$.
Pour chaque paire $i, j$, choisissons des recouvrements ouverts affines
$$
U_i \cap U_j = \bigcup\nolimits_{k \in I_{ij}} U_{ijk}
\quad\text{et}\quad
s^{-1}(U_i) \cap t^{-1}(U_j) = \bigcup\nolimits_{k \in J_{ij}} W_{ijk}.
$$
Écrivons $U_i = \Spec(A_i)$, $U_{ijk} = \Spec(A_{ijk})$,
$W_{ijk} = \Spec(B_{ijk})$.
Soit $\kappa$ un cardinal infini qui soit $\geq$ la cardinalité
de chacun des ensembles $I$, $I_{ij}$, $J_{ij}$.

\medskip\noindent
Soit $(\mathcal{F}, \alpha)$ un module quasi-cohérent sur $(U, R, s, t, c)$.
Posons $M_i = \mathcal{F}(U_i)$, $M_{ijk} = \mathcal{F}(U_{ijk})$.
Notons que
$$
M_i \otimes_{A_i} A_{ijk} = M_{ijk} = M_j \otimes_{A_j} A_{ijk}
$$
et que $\alpha$ donne des isomorphismes
$$
\alpha|_{W_{ijk}} :
M_i \otimes_{A_i, t} B_{ijk}
\longrightarrow
M_j \otimes_{A_j, s} B_{ijk}
$$
voir
Schémas, lemme \ref{schemes-lemma-widetilde-pullback}.
Au moyen de l'axiome du choix, choisissons une application
$$
(i, j, k, m) \mapsto S(i, j, k, m)
$$
qui associe à tous $i, j \in I$, $k \in I_{ij}$ ou $k \in J_{ij}$
et $m \in M_i$ une partie finie $S(i, j, k, m) \subset M_j$
telle que l'on ait
$$
m \otimes 1 = \sum\nolimits_{m' \in S(i, j, k, m)} m' \otimes a_{m'}
\quad\text{ou}\quad
\alpha(m \otimes 1) = \sum\nolimits_{m' \in S(i, j, k, m)} m' \otimes b_{m'}
$$
dans $M_{ijk}$ pour certains $a_{m'} \in A_{ijk}$ ou $b_{m'} \in B_{ijk}$.
Convenons en outre que $S(i, i, k, m) = \{m\}$ pour tous
$i, j = i, k, m$ lorsque $k \in I_{ij}$. Fixons une telle famille $S(i, j, k, m)$.

\medskip\noindent
Étant donnée une famille $\mathcal{S} = (S_i)_{i \in I}$ de parties
$S_i \subset M_i$ de cardinalité au plus $\kappa$, posons
$\mathcal{S}' = (S'_i)$, où
$$
S'_j = \bigcup\nolimits_{(i, j, k, m)\text{ tel que }m \in S_i}
S(i, j, k, m)
$$
Notons que $S_i \subset S'_i$. Notons que $S'_i$ est de cardinalité au plus
$\kappa$, car c'est une réunion, indexée par un ensemble de cardinalité au plus $\kappa$,
de parties finies. Posons $\mathcal{S}^{(0)} = \mathcal{S}$,
$\mathcal{S}^{(1)} = \mathcal{S}'$ et, par récurrence,
$\mathcal{S}^{(n + 1)} = (\mathcal{S}^{(n)})'$. Posons ensuite
$\mathcal{S}^{(\infty)} = \bigcup_{n \geq 0} \mathcal{S}^{(n)}$.
En écrivant $\mathcal{S}^{(\infty)} = (S^{(\infty)}_i)$, on voit que,
pour tout élément $m \in S^{(\infty)}_i$, l'image de $m$ dans
$M_{ijk}$ peut s'écrire comme une somme finie $\sum m' \otimes a_{m'}$
avec $m' \in S_j^{(\infty)}$. Ainsi, en posant
$$
N_i = A_i\text{-sous-module de }M_i\text{ engendré par }S^{(\infty)}_i
$$
on a
$$
N_i \otimes_{A_i} A_{ijk} = N_j \otimes_{A_j} A_{ijk}
\quad\text{et}\quad
\alpha(N_i \otimes_{A_i, t} B_{ijk}) = N_j \otimes_{A_j, s} B_{ijk}
$$
comme sous-modules de $M_{ijk}$ ou de $M_j \otimes_{A_j, s} B_{ijk}$.
Il existe donc un sous-module quasi-cohérent
$\mathcal{G} \subset \mathcal{F}$ avec $\mathcal{G}(U_i) = N_i$
tel que $\alpha(t^*\mathcal{G}) = s^*\mathcal{G}$ comme sous-modules
de $s^*\mathcal{F}$. Autrement dit,
$(\mathcal{G}, \alpha|_{t^*\mathcal{G}})$ est un sous-module de
$(\mathcal{F}, \alpha)$.
De plus, par construction, $\mathcal{G}$ est $\kappa$-engendré.

\medskip\noindent
{\emergencystretch=1em Soit $\{(\mathcal{G}_t, \alpha_t)\}_{t \in T}$ l'ensemble des
sous-modules quasi-cohérents $\kappa$-engendrés de $(\mathcal{F}, \alpha)$.
Si $t, t' \in T$, alors $\mathcal{G}_t + \mathcal{G}_{t'}$ est encore un
sous-module quasi-cohérent $\kappa$-engendré, car il est l'image du morphisme
$\mathcal{G}_t \oplus \mathcal{G}_{t'} \to \mathcal{F}$.
Le système, ordonné par inclusion, est donc filtrant.
Les arguments précédents montrent que toute section de $\mathcal{F}$ sur $U_i$
appartient à l'un des $\mathcal{G}_t$ (car on peut partir de $\mathcal{S}$
tel que la section donnée soit un élément de $S_i$). Par conséquent,
$\colim_t \mathcal{G}_t \to \mathcal{F}$ est à la fois injectif et surjectif,
comme voulu.\par}
\end{proof}





\section{Groupoïdes et schémas en groupes}
\label{section-groupoids-group-schemes}

\noindent
Il existe de nombreuses façons de construire un groupoïde à partir d'une action $a$
d'un groupe $G$ sur un ensemble $V$. Nous choisissons celle où l'on considère
un élément $g \in G$ comme une flèche de source $v$ et de but $a(g, v)$.
Cela conduit à la construction suivante pour les actions de schémas en
groupes.

\begin{lemma}
\label{lemma-groupoid-from-action}
Soit $S$ un schéma.
Soit $Y$ un schéma sur $S$.
Soit $(G, m)$ un schéma en groupes sur $Y$, de
section unité $e_G$ et de morphisme d'inversion $i_G$.
Soit $X/Y$ un schéma sur $Y$ et soit $a : G \times_Y X \to X$
une action de $G$ sur $X/Y$.
On obtient alors un schéma en groupoïdes $(U, R, s, t, c, e, i)$ sur $S$
de la manière suivante :
\begin{enumerate}
\item On pose $U = X$, et $R = G \times_Y X$.
\item On prend $s : R \to U$ égal à $(g, x) \mapsto x$.
\item On prend $t : R \to U$ égal à $(g, x) \mapsto a(g, x)$.
\item On prend $c : R \times_{s, U, t} R \to R$ égal à
$((g, x), (g', x')) \mapsto (m(g, g'), x')$.
\item On prend $e : U \to R$ égal à $x \mapsto (e_G(x), x)$.
\item On prend $i : R \to R$ égal à $(g, x) \mapsto (i_G(g), a(g, x))$.
\end{enumerate}
\end{lemma}

\begin{proof}
Omis. Indication : il suffit de vérifier que cela fonctionne au niveau
des ensembles. Pour cela, utiliser la description précédant le lemme, qui voit
$g$ comme une flèche de $v$ vers $a(g, v)$.
\end{proof}

\begin{lemma}
\label{lemma-action-groupoid-modules}
Soit $S$ un schéma.
Soit $Y$ un schéma sur $S$.
Soit $(G, m)$ un schéma en groupes sur $Y$.
Soit $X$ un schéma sur $Y$ et soit $a : G \times_Y X \to X$
une action de $G$ sur $X$ au-dessus de $Y$. Soit $(U, R, s, t, c)$
le schéma en groupoïdes construit dans le lemme \ref{lemma-groupoid-from-action}.
La règle
$(\mathcal{F}, \alpha) \mapsto (\mathcal{F}, \alpha)$ définit
une équivalence de catégories entre les objets $G$-équivariants qui sont des
$\mathcal{O}_X$-modules et la catégorie des modules quasi-cohérents
sur $(U, R, s, t, c)$.
\end{lemma}

\begin{proof}
L'assertion a un sens puisque $t = a$ et $s = \text{pr}_1$
comme morphismes $R = G \times_Y X \to X$, voir les
définitions \ref{definition-equivariant-module} et
\ref{definition-groupoid-module}.
En utilisant la traduction donnée dans le lemme \ref{lemma-groupoid-from-action},
les conditions de commutativité
des deux définitions coïncident exactement.
\end{proof}





\section{Le schéma en groupes stabilisateur}
\label{section-stabilizer}

\noindent
À partir d'un schéma en groupoïdes, on obtient un schéma en groupes comme suit.

\begin{lemma}
\label{lemma-groupoid-stabilizer}
Soit $S$ un schéma.
Soit $(U, R, s, t, c)$ un groupoïde sur $S$.
Le schéma $G$ défini par le carré cartésien
$$
\xymatrix{
G \ar[r] \ar[d] & R \ar[d]^{j = (t, s)} \\
U \ar[r]^-{\Delta} & U \times_S U
}
$$
est un schéma en groupes sur $U$, dont la loi de composition
$m$ est induite par la loi de composition $c$.
\end{lemma}

\begin{proof}
Cela résulte du fait que, dans une catégorie groupoïde, l'ensemble
des endomorphismes de tout objet forme un groupe.
\end{proof}

\noindent
Puisque $\Delta$ est une immersion, on voit que $G = j^{-1}(\Delta_{U/S})$
est un sous-schéma localement fermé de $R$. De ce point de vue,
le morphisme structural $j^{-1}(\Delta_{U/S}) \to U$ est induit par
$s$ ou par $t$ (c'est le même), et $m$ est induit par $c$.

\begin{definition}
\label{definition-stabilizer-groupoid}
Soit $S$ un schéma.
Soit $(U, R, s, t, c)$ un groupoïde sur $S$.
Le schéma en groupes $j^{-1}(\Delta_{U/S})\to U$
est appelé le {\it stabilisateur du schéma en groupoïdes
$(U, R, s, t, c)$}.
\end{definition}

\noindent
Dans la littérature, le schéma en groupes stabilisateur est souvent noté $S$
(sans doute parce que le mot stabilisateur commence par un ``s'');
nous ne pouvons le faire, puisque nous avons déjà employé $S$ pour le schéma de base.

\begin{lemma}
\label{lemma-groupoid-action-stabilizer}
Soit $S$ un schéma.
Soit $(U, R, s, t, c)$ un groupoïde sur $S$, et soit $G/U$ son stabilisateur.
Notons $R_t/U$ le schéma $R$ considéré comme schéma sur $U$ via le
morphisme $t : R \to U$.
Il existe une action à gauche canonique
$$
a : G \times_U R_t \longrightarrow R_t
$$
induite par la loi de composition $c$.
\end{lemma}

\begin{proof}
En termes de points au-dessus de $T/S$, on définit $a(g, r) = c(g, r)$.
\end{proof}

\begin{lemma}
\label{lemma-groupoid-action-stabilizer-pseudo-torsor}
Soit $S$ un schéma. Soit $(U, R, s, t, c)$ un schéma en groupoïdes
sur $S$. Soit $G$ le schéma en groupes stabilisateur de $R$.
Soit
$$
G_0 = G \times_{U, \text{pr}_0} (U \times_S U) = G \times_S U
$$
comme schéma en groupes sur $U \times_S U$. L'action de $G$ sur $R$ du
lemme \ref{lemma-groupoid-action-stabilizer}
induit une action de $G_0$ sur $R$ au-dessus de $U \times_S U$
qui fait de $R$ un pseudo-torseur sous $G_0$ sur $U \times_S U$.
\end{lemma}

\begin{proof}
Cela résulte du fait que, dans une catégorie groupoïde $\mathcal{C}$, l'ensemble
$\Mor_\mathcal{C}(x, y)$ est un espace homogène principal
sous le groupe $\Mor_\mathcal{C}(y, y)$.
\end{proof}

\begin{lemma}
\label{lemma-fibres-j}
Soit $S$ un schéma. Soit $(U, R, s, t, c)$ un schéma en groupoïdes sur $S$.
Soit $p \in U \times_S U$ un point. Notons
$R_p$ la fibre schématique de $j = (t, s) : R \to U \times_S U$.
Si $R_p \not = \emptyset$, alors l'action
$$
G_{0, \kappa(p)} \times_{\kappa(p)} R_p \longrightarrow R_p
$$
(voir
lemme \ref{lemma-groupoid-action-stabilizer-pseudo-torsor})
fait de $R_p$ un torseur sous $G_{\kappa(p)}$ sur $\kappa(p)$.
\end{lemma}

\begin{proof}
On a un pseudo-torseur d'après le lemme cité dans l'énoncé.
Et si $R_p$ n'est pas le schéma vide, alors $\{R_p \to p\}$
est un recouvrement fpqc qui trivialise le pseudo-torseur.
\end{proof}







\section{Restriction des groupoïdes}
\label{section-restrict-groupoid}

\noindent
Considérons un groupoïde (ordinaire)
$\mathcal{C} = (\text{Ob}, \text{Arrows}, s, t, c)$.
Supposons donnée une application d'ensembles $g : \text{Ob}' \to \text{Ob}$.
On peut alors construire un groupoïde
$\mathcal{C}' = (\text{Ob}', \text{Arrows}', s', t', c')$
en considérant un morphisme entre des éléments $x', y'$ de $\text{Ob}'$
comme un morphisme dans $\mathcal{C}$ entre $g(x'), g(y')$.
Autrement dit, on pose
$$
\text{Arrows}' =
\text{Ob}'
\times_{g, \text{Ob}, t}
\text{Arrows}
\times_{s, \text{Ob}, g}
\text{Ob}'.
$$
On fait les choix évidents pour $s'$, $t'$ et $c'$. Il existe un foncteur canonique
$\mathcal{C}' \to \mathcal{C}$ qui est pleinement fidèle,
mais pas nécessairement essentiellement surjectif. Ce groupoïde $\mathcal{C}'$,
muni du foncteur $\mathcal{C}' \to \mathcal{C}$,
est appelé la {\it restriction} du groupoïde
$\mathcal{C}$ à $\text{Ob}'$.

\begin{lemma}
\label{lemma-restrict-groupoid}
Soit $S$ un schéma.
Soit $(U, R, s, t, c)$ un schéma en groupoïdes sur $S$.
Soit $g : U' \to U$ un morphisme de schémas.
Considérons le diagramme suivant
$$
\xymatrix{
R' \ar[d] \ar[r] \ar@/_3pc/[dd]_{t'} \ar@/^1pc/[rr]^{s'}&
R \times_{s, U} U' \ar[r] \ar[d] &
U' \ar[d]^g \\
U' \times_{U, t} R \ar[d] \ar[r] &
R \ar[r]^s \ar[d]_t &
U \\
U' \ar[r]^g &
U
}
$$
où tous les carrés sont cartésiens. Il existe alors une
loi de composition canonique $c' : R' \times_{s', U', t'} R' \to R'$
telle que $(U', R', s', t', c')$ soit un schéma en groupoïdes sur
$S$ et que $U' \to U$, $R' \to R$ définissent un morphisme
$(U', R', s', t', c') \to (U, R, s, t, c)$ de groupoïdes en schémas sur $S$.
De plus, pour tout schéma $T$ sur $S$, le foncteur de groupoïdes
$$
(U'(T), R'(T), s', t', c') \to (U(T), R(T), s, t, c)
$$
est la restriction (voir ci-dessus) de $(U(T), R(T), s, t, c)$ par l'application
$U'(T) \to U(T)$.
\end{lemma}

\begin{proof}
Omis.
\end{proof}

\begin{definition}
\label{definition-restrict-groupoid}
Soit $S$ un schéma.
Soit $(U, R, s, t, c)$ un schéma en groupoïdes sur $S$.
Soit $g : U' \to U$ un morphisme de schémas.
Le morphisme de groupoïdes
$(U', R', s', t', c') \to (U, R, s, t, c)$
construit dans le lemme \ref{lemma-restrict-groupoid} est appelé
la {\it restriction de $(U, R, s, t, c)$ à $U'$}.
On utilise parfois la notation $R' = R|_{U'}$ dans ce cas.
\end{definition}

\begin{lemma}
\label{lemma-restrict-groupoid-relation}
Les notions de restriction des groupoïdes et des
relations de (pré-)équivalence définies dans les définitions
\ref{definition-restrict-groupoid} et \ref{definition-restrict-relation}
coïncident via les constructions des
lemmes \ref{lemma-groupoid-pre-equivalence} et
\ref{lemma-equivalence-groupoid}.
\end{lemma}

\begin{proof}
Ce que nous affirmons ici est que le $R'$ du
lemme \ref{lemma-restrict-groupoid} est également
égal à
$$
R' = (U' \times_S U')\times_{U \times_S U} R
\longrightarrow
U' \times_S U'
$$
En fait, c'eût peut-être été une manière plus claire d'énoncer ce lemme.
\end{proof}

\begin{lemma}
\label{lemma-restrict-stabilizer}
Soit $S$ un schéma.
Soit $(U, R, s, t, c)$ un schéma en groupoïdes sur $S$.
Soit $g : U' \to U$ un morphisme de schémas.
Soit $(U', R', s', t', c')$ la restriction de $(U, R, s, t, c)$ par $g$.
Soit $G$ le stabilisateur de $(U, R, s, t, c)$ et soit
$G'$ le stabilisateur de $(U', R', s', t', c')$.
Alors $G'$ est le changement de base de $G$ par $g$, c'est-à-dire
qu'il existe une identification canonique $G' = U' \times_{g, U} G$.
\end{lemma}

\begin{proof}
Omis.
\end{proof}






\section{Sous-schémas invariants}
\label{section-invariant}

\noindent
Dans cette section, nous examinons brièvement la notion de sous-schéma invariant.

\begin{definition}
\label{definition-invariant-open}
Soit $(U, R, s, t, c)$ un schéma en groupoïdes sur le schéma de base $S$.
\begin{enumerate}
\item Une partie $W \subset U$ est {\it $R$-invariante au sens ensembliste}
si $t(s^{-1}(W)) \subset W$.
\item Un ouvert $W \subset U$ est {\it $R$-invariant} si
$t(s^{-1}(W)) \subset W$.
\item Un sous-schéma fermé $Z \subset U$ est dit {\it $R$-invariant}
si $t^{-1}(Z) = s^{-1}(Z)$. On utilise ici l'image réciproque schématique, voir
Schémas, définition \ref{schemes-definition-inverse-image-closed-subscheme}.
\item Un monomorphisme de schémas $T \to U$ est {\it $R$-invariant} si
$T \times_{U, t} R = R \times_{s, U} T$ comme schémas sur $R$.
\end{enumerate}
\end{definition}

\noindent
Pour les parties et les sous-schémas ouverts $W \subset U$, l'$R$-invariance
équivaut aussi à demander que $s^{-1}(W) = t^{-1}(W)$
comme parties de $R$. Si $W \subset U$ est un sous-schéma ouvert $R$-équivariant,
alors la restriction de $R$ à $W$ est simplement $R_W = s^{-1}(W) = t^{-1}(W)$.
De même, si $Z \subset U$ est un sous-schéma fermé $R$-invariant, alors
la restriction de $R$ à $Z$ est simplement $R_Z = s^{-1}(Z) = t^{-1}(Z)$.

\begin{lemma}
\label{lemma-constructing-invariant-opens}
Soit $S$ un schéma.
Soit $(U, R, s, t, c)$ un schéma en groupoïdes sur $S$.
\begin{enumerate}
\item Pour toute partie $W \subset U$, la partie $t(s^{-1}(W))$
est $R$-invariante au sens ensembliste.
\item Si $s$ et $t$ sont ouverts, alors, pour tout ouvert $W \subset U$,
l'ouvert $t(s^{-1}(W))$ est un sous-schéma ouvert $R$-invariant.
\item Si $s$ et $t$ sont ouverts et quasi-compacts, alors $U$ possède un
recouvrement ouvert formé de sous-schémas ouverts quasi-compacts $R$-invariants.
\end{enumerate}
\end{lemma}

\begin{proof}
{\emergencystretch=2em La partie (1) résulte des
Lemmes \ref{lemma-pre-equivalence-equivalence-relation-points} et
\ref{lemma-groupoid-pre-equivalence}, à savoir que $t(s^{-1}(W))$
est l'ensemble des points de $U$ équivalents à un point de $W$.
Supposons ensuite $s$ et $t$ ouverts, et $W \subset U$ ouvert.
Comme $s$ est ouvert, l'ensemble $W' = t(s^{-1}(W))$ est un ouvert de $U$.
Supposons enfin que $s$, $t$ soient tous deux ouverts et quasi-compacts.
Alors, si $W \subset U$ est un ouvert quasi-compact,
$W' = t(s^{-1}(W))$ est lui aussi un ouvert quasi-compact et il est invariant d'après la
discussion précédente. En faisant parcourir à $W$ tous les ouverts affines de $U$,
on obtient (3).\par}
\end{proof}

\begin{lemma}
\label{lemma-first-observation}
Soit $S$ un schéma. Soit $(U, R, s, t, c)$ un schéma en groupoïdes sur $S$.
Supposons $s$ et $t$ quasi-compacts et plats, et $U$ quasi-séparé.
Soit $W \subset U$ un ouvert quasi-compact. Alors $t(s^{-1}(W))$
est l'intersection d'une famille non vide d'ouverts quasi-compacts de $U$.
\end{lemma}

\begin{proof}
Notons que $s^{-1}(W)$ est un ouvert quasi-compact de $R$.
En tant qu'application continue, $t$ envoie le sous-ensemble quasi-compact
$s^{-1}(W)$ sur un sous-ensemble quasi-compact $t(s^{-1}(W))$.
Comme $t$ est plat et que $s^{-1}(W)$ est stable par générisation,
il en va de même de $t(s^{-1}(W))$, voir
(Morphismes, Lemme \ref{morphisms-lemma-generalizations-lift-flat} et
Topologie, Lemme \ref{topology-lemma-lift-specializations-images}).
Choisissons un ouvert quasi-compact $W' \subset U$ contenant $t(s^{-1}(W))$. D'après
Propriétés, Lemme \ref{properties-lemma-quasi-compact-quasi-separated-spectral},
on voit que $W'$ est un espace spectral (c'est ici que l'on utilise que $U$ est quasi-séparé).
Le lemme résulte alors de
Topologie, Lemme \ref{topology-lemma-make-spectral-space},
appliqué à $t(s^{-1}(W)) \subset W'$.
\end{proof}

\begin{lemma}
\label{lemma-second-observation}
Hypothèses et notations comme dans le Lemme \ref{lemma-first-observation}.
Il existe un ouvert $R$-invariant $V \subset U$ et un ouvert quasi-compact
$W'$ tels que $W \subset V \subset W' \subset U$.
\end{lemma}

\begin{proof}
Posons $E = t(s^{-1}(W))$. Rappelons que $E$ est $R$-invariant au sens ensembliste
(Lemme \ref{lemma-constructing-invariant-opens}).
D'après le Lemme \ref{lemma-first-observation}, il existe un ouvert quasi-compact
$W'$ contenant $E$. Posons $Z = U \setminus W'$ et considérons
$T = t(s^{-1}(Z))$. Observons que $Z \subset T$ et que
$E \cap T = \emptyset$, car $s^{-1}(E) = t^{-1}(E)$ est disjoint
de $s^{-1}(Z)$. Puisque $T$ est l'image du sous-ensemble fermé
$s^{-1}(Z) \subset R$ par le morphisme quasi-compact $t : R \to U$,
on voit que tout point $\xi$ de l'adhérence $\overline{T}$
est une spécialisation d'un point de $T$, voir
Morphismes, Lemme \ref{morphisms-lemma-reach-points-scheme-theoretic-image} (et
Morphismes, Lemme \ref{morphisms-lemma-quasi-compact-scheme-theoretic-image}
pour voir que l'image schématique est l'adhérence de l'image).
Écrivons $\xi' \leadsto \xi$ avec $\xi' \in T$. Supposons que $r \in R$ et
$s(r) = \xi$. Comme $s$ est plat, on peut trouver une spécialisation $r' \leadsto r$
dans $R$ telle que $s(r') = \xi'$
(Morphismes, Lemme \ref{morphisms-lemma-generalizations-lift-flat}).
Alors $t(r') \leadsto t(r)$. On conclut que $t(r') \in T$, puisque $T$
est invariant au sens ensembliste d'après
le Lemme \ref{lemma-constructing-invariant-opens}.
Ainsi, $\overline{T}$ est un sous-ensemble fermé $R$-invariant au sens ensembliste
et $V = U \setminus \overline{T}$ est l'ouvert que nous
cherchions. Il est contenu dans $W'$, ce qui achève la démonstration.
\end{proof}




\section{Faisceaux quotients}
\label{section-quotient-sheaves}

\noindent
Soit $\tau \in\allowbreak \{Zariski,\allowbreak \etale,\allowbreak fppf,\allowbreak smooth,\allowbreak syntomic\}$.
Soit $S$ un schéma.
Soit $j : R \to U \times_S U$ une prérelation sur $S$.
Supposons que $U, R, S$ soient des objets d'un $\tau$-site $\Sch_\tau$
(voir Topologies, section \ref{topologies-section-procedure}).
On peut alors considérer les foncteurs
$$
h_U, h_R :
(\Sch/S)_\tau^{opp}
\longrightarrow
\textit{Ensembles}.
$$
Ce sont des faisceaux, voir
Descente, Lemme \ref{descent-lemma-fpqc-universal-effective-epimorphisms}.
Le morphisme $j$ induit une application $j : h_R \to h_U \times h_U$.
Pour tout objet $T \in \Ob((\Sch/S)_\tau)$,
on peut prendre la relation d'équivalence $\sim_T$ engendrée par
$j(T) : R(T) \to U(T) \times U(T)$ et considérer le quotient.
On obtient ainsi un préfaisceau
\begin{equation}
\label{equation-quotient-presheaf}
(\Sch/S)_\tau^{opp}
\longrightarrow
\textit{Ensembles}, \quad
T \longmapsto U(T)/\sim_T
\end{equation}

\begin{definition}
\label{definition-quotient-sheaf}
Soient $\tau$, $S$ et la prérelation $j : R \to U \times_S U$ comme ci-dessus.
Dans cette situation, le {\it faisceau quotient $U/R$} associé
à $j$ est le faisceau associé au préfaisceau
(\ref{equation-quotient-presheaf}) pour la topologie $\tau$.
Si $j : R \to U \times_S U$ provient de l'action d'un schéma en groupes
$G/S$ sur $U$ comme dans le Lemme \ref{lemma-groupoid-from-action}, on
note parfois le faisceau quotient $U/G$.
\end{definition}

\noindent
Cela signifie exactement que le diagramme
$$
\xymatrix{
h_R \ar@<1ex>[r] \ar@<-1ex>[r] &
h_U \ar[r] &
U/R
}
$$
est un diagramme coégalisateur dans la catégorie des faisceaux d'ensembles
sur $(\Sch/S)_\tau$. Au moyen du plongement de Yoneda, on
peut considérer $(\Sch/S)_\tau$ comme une sous-catégorie pleine des
faisceaux sur $(\Sch/S)_\tau$ et ainsi identifier les schémas
aux foncteurs représentables. Par cet abus de notation,
nous représenterons souvent le diagramme ci-dessus simplement par
$$
\xymatrix{
R \ar@<1ex>[r]^s \ar@<-1ex>[r]_t &
U \ar[r] &
U/R
}
$$
Nous travaillerons le plus souvent avec la topologie fppf lorsque nous considérerons
les faisceaux quotients de groupoïdes ou de relations d'équivalence.

\begin{definition}
\label{definition-representable-quotient}
Dans la situation de la Définition \ref{definition-quotient-sheaf}.
On dit que la prérelation $j$ admet un
{\it quotient représentable} si le faisceau $U/R$ est représentable.
On dira qu'un groupoïde $(U, R, s, t, c)$ admet un
{\it quotient représentable}
si le quotient $U/R$ correspondant à $j = (t, s)$ est représentable.
\end{definition}

\noindent
Le lemme suivant caractérise les schémas $M$ qui représentent le quotient.
Il s'applique par exemple si $\tau = fppf$, si $U \to M$ est plat,
de présentation finie et surjectif, et si $R \cong U \times_M U$.

\begin{lemma}
\label{lemma-criterion-quotient-representable}
Dans la situation de la Définition \ref{definition-quotient-sheaf}.
Supposons qu'il existe un schéma $M$ et un morphisme $U \to M$ tels que
\begin{enumerate}
\item le morphisme $U \to M$ égalise $s, t$,
\item le morphisme $U \to M$ induit une surjection de faisceaux
$h_U \to h_M$ pour la topologie $\tau$, et
\item l'application induite $(t, s) : R \to U \times_M U$ induit une
surjection de faisceaux $h_R \to h_{U \times_M U}$ pour la topologie $\tau$.
\end{enumerate}
Alors $M$ représente le faisceau quotient $U/R$.
\end{lemma}

\begin{proof}
La condition (1) dit que $h_U \to h_M$ se factorise par $U/R$.
La condition (2) dit que $U/R \to h_M$ est surjectif comme morphisme de faisceaux.
La condition (3) dit que $U/R \to h_M$ est injectif comme morphisme de faisceaux.
Le lemme en résulte.
\end{proof}

\noindent
Le lemme suivant est faux si l'on n'exige pas que $j$ soit une
prérelation d'équivalence (mais seulement, disons, une prérelation).

\begin{lemma}\emergencystretch=1em
\label{lemma-quotient-pre-equivalence}
Soit $\tau \in\allowbreak \{Zariski,\allowbreak \etale,\allowbreak fppf,\allowbreak smooth,\allowbreak syntomic\}$.
Soit $S$ un schéma.
Soit $j : R \to U \times_S U$ une prérelation d'équivalence sur $S$.
Supposons que $U, R, S$ soient des objets d'un $\tau$-site $\Sch_\tau$.
Pour $T \in \Ob((\Sch/S)_\tau)$ et
$a, b \in U(T)$, les conditions suivantes sont équivalentes :
\begin{enumerate}
\item $a$ et $b$ ont la même image dans $(U/R)(T)$, et
\item il existe un $\tau$-recouvrement $\{f_i : T_i \to T\}$ de $T$
et des morphismes $r_i : T_i \to R$ tels que
$a \circ f_i = s \circ r_i$ et $b \circ f_i = t \circ r_i$.
\end{enumerate}
Autrement dit, dans ce cas, le morphisme de $\tau$-faisceaux
$$
h_R \longrightarrow h_U \times_{U/R} h_U
$$
est surjectif.
\end{lemma}

\begin{proof}
Omis. Indication : cela fonctionne parce que, dans ce cas, le préfaisceau
(\ref{equation-quotient-presheaf}) est en réalité donné
par $T \mapsto U(T)/j(R(T))$, puisque $j(R(T)) \subset U(T) \times U(T)$
est une relation d'équivalence, voir
Définition \ref{definition-equivalence-relation}.
\end{proof}

\begin{lemma}\emergencystretch=1em
\label{lemma-quotient-pre-equivalence-relation-restrict}
Soit $\tau \in\allowbreak \{Zariski,\allowbreak \etale,\allowbreak fppf,\allowbreak smooth,\allowbreak syntomic\}$.
Soit $S$ un schéma.
Soit $j : R \to U \times_S U$ une prérelation d'équivalence sur $S$
et $g : U' \to U$ un morphisme de schémas sur $S$.
Soit $j' : R' \to U' \times_S U'$ la restriction de $j$ à $U'$.
Supposons que $U, U', R, S$ soient des objets d'un $\tau$-site $\Sch_\tau$.
Le morphisme de faisceaux quotients
$$
U'/R' \longrightarrow U/R
$$
est injectif. Si $g$ définit une surjection $h_{U'} \to h_U$ de faisceaux
pour la topologie $\tau$ (par exemple, si $\{g : U' \to U\}$ est un
$\tau$-recouvrement), alors $U'/R' \to U/R$ est un isomorphisme.
\end{lemma}

\begin{proof}
Supposons que $\xi, \xi' \in (U'/R')(T)$ soient des sections dont les
images dans $U/R$ sont égales.
On peut alors trouver un $\tau$-recouvrement $\mathcal{T} = \{T_i \to T\}$ de $T$
tel que $\xi|_{T_i}, \xi'|_{T_i}$ soient données par $a_i, a_i' \in U'(T_i)$. D'après le
Lemme \ref{lemma-quotient-pre-equivalence}
et les axiomes d'un site, on peut, après avoir raffiné
$\mathcal{T}$, supposer qu'il existe des morphismes $r_i : T_i \to R$
tels que $g \circ a_i = s \circ r_i$, $g \circ a_i' = t \circ r_i$.
Puisque, par construction,
$R' = R \times_{U \times_S U} (U' \times_S U')$,
on voit que $(r_i, (a_i, a_i')) \in R'(T_i)$, ce qui
montre que $a_i$ et $a_i'$ définissent la même section
de $U'/R'$ au-dessus de $T_i$. La condition de faisceau implique alors
$\xi = \xi'$.

\medskip\noindent
Si $h_{U'} \to h_U$ est une surjection
de faisceaux, alors, bien sûr, $U'/R' \to U/R$ est également surjectif.
Si $\{g : U' \to U\}$ est un $\tau$-recouvrement, alors
le morphisme de faisceaux $h_{U'} \to h_U$ est surjectif, voir
Sites, Lemme \ref{sites-lemma-covering-surjective-after-sheafification}.
Ainsi, $U'/R' \to U/R$ est également surjectif dans ce cas.
\end{proof}

\begin{lemma}\emergencystretch=1em
\label{lemma-quotient-groupoid-restrict}
Soit $\tau \in\allowbreak \{Zariski,\allowbreak \etale,\allowbreak fppf,\allowbreak smooth,\allowbreak syntomic\}$.
Soit $S$ un schéma.
Soit $(U, R, s, t, c)$ un schéma en groupoïdes sur $S$.
Soit $g : U' \to U$ un morphisme de schémas sur $S$.
Soit $(U', R', s', t', c')$ la restriction de $(U, R, s, t, c)$ à $U'$.
Supposons que $U, U', R, S$ soient des objets d'un $\tau$-site $\Sch_\tau$.
Le morphisme de faisceaux quotients
$$
U'/R' \longrightarrow U/R
$$
est injectif. Si la composée
$$
\xymatrix{
U' \times_{g, U, t} R \ar[r]_-{\text{pr}_1} \ar@/^3ex/[rr]^h
& R \ar[r]_s & U
}
$$
définit une surjection de faisceaux pour la topologie $\tau$, alors
le morphisme est bijectif. Cela vaut par exemple si
$\{h : U' \times_{g, U, t} R \to U\}$ est un $\tau$-recouvrement, ou
si $U' \to U$ définit une surjection de faisceaux pour la topologie $\tau$, ou si
$\{g : U' \to U\}$ est un recouvrement pour la topologie $\tau$.
\end{lemma}

\begin{proof}
L'injectivité résulte de la combinaison des
Lemmes \ref{lemma-groupoid-pre-equivalence} et
\ref{lemma-quotient-pre-equivalence-relation-restrict}.
Pour établir la surjectivité (voir
Sites, section \ref{sites-section-sheaves-injective},
pour une caractérisation des morphismes surjectifs de faisceaux), raisonnons comme suit.
Supposons que $T$ soit un schéma et que $\sigma \in U/R(T)$.
Il existe un recouvrement $\{T_i \to T\}$ tel que $\sigma|_{T_i}$
soit l'image d'un élément $f_i \in U(T_i)$. On peut donc
supposer que $\sigma$ est l'image de $f \in U(T)$.
Comme, par hypothèse, $h$ est une surjection de faisceaux, on
peut trouver un $\tau$-recouvrement $\{\varphi_i : T_i \to T\}$ et des morphismes
$f_i : T_i \to U' \times_{g, U, t} R$ tels que
$f \circ \varphi_i = h \circ f_i$. Notons
$f'_i = \text{pr}_0 \circ f_i : T_i \to U'$. Alors on voit que
$f'_i \in U'(T_i)$ s'envoie sur $g \circ f'_i \in U(T_i)$ et
que $g \circ f'_i \sim_{T_i} h \circ f_i = f \circ \varphi_i$,
avec les notations de (\ref{equation-quotient-presheaf}). Plus précisément, un
élément de $R(T_i)$ donnant la relation est $\text{pr}_1 \circ f_i$.
Cela signifie que la restriction
de $\sigma$ à $T_i$ appartient à l'image de $U'/R'(T_i) \to U/R(T_i)$,
comme voulu.

\medskip\noindent
Si $\{h\}$ est un $\tau$-recouvrement, il induit une surjection de faisceaux, voir
Sites, Lemme \ref{sites-lemma-covering-surjective-after-sheafification}.
Si $U' \to U$ est surjectif, alors $h$ l'est aussi, car $s$ possède une section
(à savoir l'élément neutre $e$ du schéma en groupoïdes).
\end{proof}

\begin{lemma}\emergencystretch=1em
\label{lemma-criterion-fibre-product}
Soit $S$ un schéma. Soit $f : (U, R, j) \to (U', R', j')$ un morphisme
entre des relations d'équivalence sur $S$. Supposons que
$$
\xymatrix{
R \ar[d]_s \ar[r]_f & R' \ar[d]^{s'} \\
U \ar[r]^f & U'
}
$$
soit cartésien. Pour tout
$\tau \in\allowbreak \{Zariski,\allowbreak \etale,\allowbreak fppf,\allowbreak smooth,\allowbreak syntomic\}$,
le diagramme
$$
\xymatrix{
U \ar[d] \ar[r] & U/R \ar[d]^f \\
U' \ar[r] & U'/R'
}
$$
est un carré cartésien de $\tau$-faisceaux.
\end{lemma}

\begin{proof}
D'après le Lemme \ref{lemma-quotient-pre-equivalence}, les faisceaux quotients
admettent une description simple, que nous utiliserons ci-dessous sans autre mention.
Montrons d'abord que
$$
U \longrightarrow U' \times_{U'/R'} U/R
$$
est injectif. En effet, supposons que $a, b \in U(T)$ aient la même image
dans le membre de droite. Alors $f(a) = f(b)$. Après avoir remplacé $T$ par les
membres d'un $\tau$-recouvrement, on peut supposer qu'il existe
$r \in R(T)$ tel que $a = s(r)$ et $b = t(r)$. Alors $r' = f(r)$
est un $T$-point de $R'$ tel que $s'(r') = t'(r')$. Ainsi,
$r' = e'(f(a))$ (où $e'$ est l'identité du groupoïde
associé à $j'$, voir le Lemme \ref{lemma-equivalence-groupoid}).
Comme le premier diagramme du lemme est cartésien, cela implique
nécessairement que $r$ est égal à $e(a)$. Ainsi, $a = b$.

\medskip\noindent
Enfin, montrons que la flèche affichée est surjective. Soit
$T$ un schéma sur $S$ et soit $(a', \overline{b})$ une section
du faisceau $U' \times_{U'/R'} U/R$ au-dessus de $T$. Après avoir remplacé $T$
par les membres d'un $\tau$-recouvrement, on peut supposer que $\overline{b}$
est la classe d'un élément $b \in U(T)$. Après avoir remplacé $T$
par les membres d'un $\tau$-recouvrement, on peut supposer qu'il existe
$r' \in R'(T)$ tel que $a' = t(r')$ et $s'(r') = f(b)$.
Comme le premier diagramme du lemme est cartésien, on peut trouver
$r \in R(T)$ tel que $s(r) = b$ et $f(r) = r'$. Il est alors clair
que $a = t(r) \in U(T)$ est une section qui s'envoie sur
$(a', \overline{b})$.
\end{proof}








\section{Descente en termes de groupoïdes}
\label{section-groupoids-descent}

\noindent
Les morphismes cartésiens sont définis comme suit.

\begin{definition}
\label{definition-cartesian-morphism}
Soit $S$ un schéma. Soit $f : (U', R', s', t', c') \to (U, R, s, t, c)$
un morphisme de schémas en groupoïdes sur $S$. On dit que $f$ est {\it cartésien}, ou
que {\it $(U', R', s', t', c')$ est cartésien au-dessus de $(U, R, s, t, c)$},
si le diagramme
$$
\xymatrix{
R' \ar[r]_f \ar[d]_{s'} & R \ar[d]^s \\
U' \ar[r]^f & U
}
$$
est un carré cartésien dans la catégorie des schémas. Un {\it morphisme de schémas en groupoïdes
cartésien au-dessus de $(U, R, s, t, c)$} est un morphisme de schémas en groupoïdes
compatible avec les morphismes structuraux vers $(U, R, s, t, c)$.
\end{definition}

\noindent
Les morphismes cartésiens sont reliés aux données de descente. Nous démontrons d'abord un
lemme général décrivant la catégorie des schémas en groupoïdes cartésiens au-dessus d'un
schéma en groupoïdes fixé.

\begin{lemma}
\label{lemma-characterize-cartesian-schemes}
Soit $S$ un schéma. Soit $(U, R, s, t, c)$ un schéma en groupoïdes sur $S$.
La catégorie des schémas en groupoïdes cartésiens au-dessus de $(U, R, s, t, c)$
est équivalente à la catégorie des couples $(V, \varphi)$, où $V$ est un
schéma sur $U$ et où
$$
\varphi :
V \times_{U, t} R
\longrightarrow
R \times_{s, U} V
$$
est un isomorphisme au-dessus de $R$ tel que $e^*\varphi = \text{id}_V$ et tel que
$$
c^*\varphi = \text{pr}_1^*\varphi \circ \text{pr}_0^*\varphi
$$
comme morphismes de schémas au-dessus de $R \times_{s, U, t} R$.
\end{lemma}

\begin{proof}
Dans le lemme, la notation de tiré en arrière désigne le changement de base. La formule
affichée a un sens parce que
$$
(R \times_{s, U, t} R) \times_{\text{pr}_1, R, \text{pr}_1} (V \times_{U, t} R)
=
(R \times_{s, U, t} R) \times_{\text{pr}_0, R, \text{pr}_0} (R \times_{s, U} V)
$$
comme schémas au-dessus de $R \times_{s, U, t} R$.

\medskip\noindent
Étant donné $(V, \varphi)$, posons $U' = V$ et $R' = V \times_{U, t} R$.
On prend pour $t' : R' \to U'$ la projection $V \times_{U, t} R \to V$.
On prend pour $s'$ le composé de $\varphi$ et de la projection
$R \times_{s, U} V \to V$. On prend pour $c'$ le composé
\begin{align*}
R' \times_{s', U', t'} R'
& \xrightarrow{\varphi, 1}
(R \times_{s, U} V) \times_V (V \times_{U, t} R) \\
& \xrightarrow{}
R \times_{s, U} V \times_{U, t} R \\
& \xrightarrow{\varphi^{-1}, 1}
V \times_{U, t} (R \times_{s, U, t} R) \\
& \xrightarrow{1, c}
V \times_{U, t} R = R'
\end{align*}
Un calcul, que nous omettons, montre que l'on obtient un schéma en groupoïdes
au-dessus de $(U, R, s, t, c)$. Il est clair que ce schéma en groupoïdes est
cartésien au-dessus de $(U, R, s, t, c)$.

\medskip\noindent
Réciproquement, étant donné $f : (U', R', s', t', c') \to (U, R, s, t, c)$
cartésien, les morphismes
$$
U' \times_{U, t} R \xleftarrow{t', f} R' \xrightarrow{f, s'} R \times_{s, U} U'
$$
sont des isomorphismes; on peut donc poser $V = U'$ et prendre pour $\varphi$ le
composé $(f, s') \circ (t', f)^{-1}$. Nous omettons la démonstration du fait que
$\varphi$ satisfait aux conditions du lemme. Nous omettons la démonstration du fait que
ces constructions sont inverses l'une de l'autre.
\end{proof}

\noindent
Soit $S$ un schéma. Soit $f : X \to Y$ un morphisme de schémas sur $S$. On
obtient alors un schéma en groupoïdes $(X, X \times_Y X, \text{pr}_1, \text{pr}_0, c)$
sur $S$. En effet, $j : X \times_Y X \to X \times_S X$ est une relation
d'équivalence et l'on peut prendre le groupoïde associé, voir le
Lemme \ref{lemma-equivalence-groupoid}.

\begin{lemma}
\label{lemma-cartesian-equivalent-descent-datum}
Soit $S$ un schéma. Soit $f : X \to Y$ un morphisme de schémas sur $S$.
La construction du Lemme \ref{lemma-characterize-cartesian-schemes}
détermine une équivalence
$$
\begin{matrix}
\text{catégorie des schémas en groupoïdes} \\
\text{cartésiens au-dessus de } (X, X \times_Y X, \ldots)
\end{matrix}
\longrightarrow
\begin{matrix}
\text{ catégorie des données de descente} \\
\text{ relatives à } X/Y
\end{matrix}
$$
\end{lemma}

\begin{proof}
Cela résulte immédiatement du
Lemme \ref{lemma-characterize-cartesian-schemes}
et de la définition des données de descente sur les schémas dans
Descente, Définition \ref{descent-definition-descent-datum}.
\end{proof}







\section{Conditions de séparation}
\label{section-separation}

\noindent
Il s'agit en réalité de conditions sur le morphisme $j : R \to U \times_S U$
lorsque l'on se donne un groupoïde $(U, R, s, t, c)$ sur $S$. Comme dans la section
précédente, nous formulons d'abord le diagramme correspondant.

\begin{lemma}
\label{lemma-diagram-diagonal}
Soit $S$ un schéma.
Soit $(U, R, s, t, c)$ un groupoïde sur $S$.
Soit $G \to U$ le schéma en groupes stabilisateur.
Le diagramme commutatif
$$
\xymatrix{
R \ar[d]^{\Delta_{R/U \times_S U}} \ar[rrr]_{f \mapsto (f, s(f))} & & &
R \times_{s, U} U \ar[d] \ar[r] & U \ar[d] \\
R \times_{(U \times_S U)} R \ar[rrr]^{(f, g) \mapsto (f, f^{-1} \circ g)} & & &
R \times_{s, U} G \ar[r] & G
}
$$
possède deux flèches horizontales de gauche qui sont des isomorphismes,
et son carré de droite est cartésien.
\end{lemma}

\begin{proof}
Omis.
Exercice sur les définitions et le point de vue fonctoriel en
géométrie algébrique.
\end{proof}

\begin{lemma}
\label{lemma-diagonal}
Soit $S$ un schéma.
Soit $(U, R, s, t, c)$ un groupoïde sur $S$.
Soit $G \to U$ le schéma en groupes stabilisateur.
\begin{enumerate}
\item Les conditions suivantes sont équivalentes :
\begin{enumerate}
\item $j : R \to U \times_S U$ est séparé,
\item $G \to U$ est séparé, et
\item $e : U \to G$ est une immersion fermée.
\end{enumerate}
\item Les conditions suivantes sont équivalentes :
\begin{enumerate}
\item $j : R \to U \times_S U$ est quasi-séparé,
\item $G \to U$ est quasi-séparé, et
\item $e : U \to G$ est quasi-compact.
\end{enumerate}
\end{enumerate}
\end{lemma}

\begin{proof}
Le schéma en groupes $G \to U$ est le changement de base de $R \to U \times_S U$
par le morphisme diagonal $U \to U \times_S U$, voir le
Lemme \ref{lemma-groupoid-stabilizer}. Par conséquent, si
$j$ est séparé (resp.\ quasi-séparé),
alors $G \to U$ est séparé (resp.\ quasi-séparé).
(Voir Schémas, Lemme
\ref{schemes-lemma-separated-permanence}).
Ainsi, (a) $\Rightarrow$ (b) dans (1) comme dans (2).

\medskip\noindent
Si $G \to U$ est séparé (resp.\ quasi-séparé), alors le morphisme
$U \to G$, en tant que section du morphisme structural $G \to U$, est une immersion
fermée (resp.\ quasi-compact), voir
Schémas, Lemme \ref{schemes-lemma-section-immersion}.
Ainsi, (b) $\Rightarrow$ (a) dans (1) comme dans (2).

\medskip\noindent
D'après le résultat du
Lemme \ref{lemma-diagram-diagonal}
(et Schémas, Lemmes \ref{schemes-lemma-base-change-immersion}
et \ref{schemes-lemma-quasi-compact-preserved-base-change}),
on voit que, si $e$ est une immersion fermée (resp.\ quasi-compact),
$\Delta_{R/U \times_S U}$ est une immersion
fermée (resp.\ quasi-compact).
Ainsi, (c) $\Rightarrow$ (a) dans (1) comme dans (2).
\end{proof}









\section{Groupoïdes finis plats, cas affine}
\label{section-finite-flat}

\noindent
Soit $S$ un schéma.
Soit $(U, R, s, t, c)$ un schéma en groupoïdes sur $S$.
Supposons que $U = \Spec(A)$ et $R = \Spec(B)$ soient affines.
Dans ce cas, on obtient deux homomorphismes d'anneaux
$s^\sharp, t^\sharp : A \longrightarrow B$.
Soit $C$ l'égalisateur de $s^\sharp$ et $t^\sharp$. En formule,
\begin{equation}
\label{equation-invariants}
C = \{a \in A \mid t^\sharp(a) = s^\sharp(a) \}.
\end{equation}
Nous l'appellerons parfois {\it anneau des fonctions $R$-invariantes sur $U$}.
Quelles propriétés $M = \Spec(C)$ possède-t-il ? La première observation est
que le diagramme
$$
\xymatrix{
R \ar[r]_s \ar[d]_t & U \ar[d] \\
U \ar[r] & M
}
$$
est commutatif, c'est-à-dire que le morphisme $U \to M$ égalise $s, t$.
De plus, si $T$ est un schéma affine quelconque, et si $U \to T$ est
un morphisme qui égalise $s, t$, alors $U \to T$ se factorise par $U \to M$.
Autrement dit, $U \to M$ est un coégalisateur dans la catégorie des schémas affines.

\medskip\noindent
Nous aimerions trouver des conditions garantissant que le morphisme $U \to M$ est
véritablement un ``quotient'' dans la catégorie des schémas. Nous en discuterons
longuement ailleurs (insérer ici une référence ultérieure) ; ici, nous examinons seulement
quelques cas particuliers. Plus précisément, nous nous concentrerons sur le cas où $s, t$ sont
finis localement libres.

\begin{example}
\label{example-quotient-projective-line}
Soit $k$ un corps. Soit $U = \text{GL}_{2, k}$. Soit $B \subset \text{GL}_2$
le sous-schéma en groupes fermé des matrices triangulaires supérieures.
Alors le faisceau quotient $\text{GL}_{2, k}/B$ (pour la topologie de Zariski, \'etale ou
fppf, voir la Définition \ref{definition-quotient-sheaf}) est
représentable par la droite projective : $\mathbf{P}^1 = \text{GL}_{2, k}/B$.
(Détails omis.)
En revanche, l'anneau des fonctions invariantes dans ce cas est simplement $k$.
Notons que, dans ce cas, les morphismes
$s, t : R = \text{GL}_{2, k} \times_k B \to \text{GL}_{2, k} = U$ sont lisses
de dimension relative $3$.
\end{example}

\noindent
Rappelons que, dans Exercices, Exercices \ref{exercises-exercise-trace-det} et
\ref{exercises-exercise-trace-det-rings}, nous avons défini le déterminant
et la norme pour les modules finis localement libres et les extensions d'anneaux
finies localement libres. Si $\varphi : A \to B$ est un homomorphisme d'anneaux fini localement libre, alors
nous noterons $\text{Norm}_\varphi(b) \in A$ la norme de $b \in B$. Dans le
cas d'un morphisme de schémas fini localement libre, la norme a été construite dans
Diviseurs, Lemme \ref{divisors-lemma-finite-locally-free-has-norm}.

\begin{lemma}
\label{lemma-determinant-trick}
Soit $S$ un schéma. Soit $(U, R, s, t, c)$ un schéma en groupoïdes sur $S$.
Supposons que $U = \Spec(A)$ et $R = \Spec(B)$ soient affines et que
$s, t : R \to U$ soient finis localement libres.
Soit $C$ comme dans (\ref{equation-invariants}).
Soit $f \in A$. Alors $\text{Norm}_{s^\sharp}(t^\sharp(f)) \in C$.
\end{lemma}

\begin{proof}
Considérons le diagramme commutatif
$$
\xymatrix{
& U & \\
R \ar[d]_s \ar[ru]^t &
R \times_{s, U, t} R
\ar[l]^-{\text{pr}_0} \ar[d]^{\text{pr}_1} \ar[r]_-c &
R \ar[d]^s \ar[lu]_t \\
U & R \ar[l]_t \ar[r]^s & U
}
$$
{\emergencystretch=2em Il s'agit du diagramme du lemme \ref{lemma-diagram}.
Considérons $f \in \Gamma(U, \mathcal{O}_U)$. La commutativité de la
partie supérieure du diagramme montre que
$\text{pr}_0^\sharp(t^\sharp(f)) = c^\sharp(t^\sharp(f))$ comme éléments de
$\Gamma(R \times_{S, U, t} R, \mathcal{O})$.
En examinant le carré cartésien inférieur droit,
la compatibilité de la construction de la norme avec le changement de base montre que
$s^\sharp(\text{Norm}_{s^\sharp}(t^\sharp(f))) =
\text{Norm}_{\text{pr}_1^\sharp}(c^\sharp(t^\sharp(f)))$.
De même, nous obtenons
$t^\sharp(\text{Norm}_{s^\sharp}(t^\sharp(f))) =
\text{Norm}_{\text{pr}_1^\sharp}(\text{pr}_0^\sharp(t^\sharp(f)))$.
Par conséquent, la première égalité de cette preuve donne
$s^\sharp(\text{Norm}_{s^\sharp}(t^\sharp(f))) =
t^\sharp(\text{Norm}_{s^\sharp}(t^\sharp(f)))$, comme souhaité.\par}
\end{proof}

\begin{lemma}
\label{lemma-finite-locally-free-disjoint-free}
Soit $S$ un schéma. Soit $(U, R, s, t, c)$ un schéma en groupoïdes sur $S$.
Supposons que $s, t : R \to U$ soient finis localement libres.
Alors
$$
U = \coprod\nolimits_{r \geq 1} U_r
$$
est une réunion disjointe d'ouverts invariants par $R$ telle que la restriction $R_r$ de
$R$ à $U_r$ soit telle que $s, t : R_r \to U_r$ soient finis localement
libres de rang $r$.
\end{lemma}

\begin{proof}
D'après
Morphismes, Lemme \ref{morphisms-lemma-finite-locally-free},
il existe une décomposition
$U = \coprod\nolimits_{r \geq 0} U_r$
telle que $s : s^{-1}(U_r) \to U_r$ soit fini localement libre de rang $r$.
Comme $s$ est surjectif, on voit que $U_0 = \emptyset$.
Notons que $u \in U_r \Leftrightarrow$ la fibre schématique
$s^{-1}(u)$ est de degré $r$ sur $\kappa(u)$. Maintenant, si $z \in R$ avec $s(z) = u$
et $t(z) = u'$, alors, avec les notations du Lemme \ref{lemma-diagram},
$$
\text{pr}_1^{-1}(z) \to \Spec(\kappa(z))
$$
est à la fois le changement de base de
$s^{-1}(u) \to \Spec(\kappa(u))$ et de $s^{-1}(u') \to \Spec(\kappa(u'))$
d'après le lemme cité. Ainsi $u \in U_r \Leftrightarrow u' \in U_r$ ;
autrement dit, les ouverts $U_r$ sont invariants par $R$.
En particulier, la restriction de $R$ à $U_r$ est simplement
$s^{-1}(U_r)$ et $s : R_r \to U_r$ est fini localement libre de rang $r$.
Comme $t : R_r \to U_r$ est isomorphe à $s$ par l'inversion de $R_r$,
on voit qu'il est également de rang $r$.
\end{proof}

\begin{lemma}
\label{lemma-integral-over-invariants}
Soit $S$ un schéma. Soit $(U, R, s, t, c)$ un schéma en groupoïdes sur $S$.
Supposons que $U = \Spec(A)$ et $R = \Spec(B)$ soient affines et que
$s, t : R \to U$ soient finis localement libres.
Soit $C \subset A$ comme dans (\ref{equation-invariants}).
Alors $A$ est entier sur $C$.
\end{lemma}

\begin{proof}
Tout d'abord, le Lemme \ref{lemma-finite-locally-free-disjoint-free}
nous apprend que $(U, R, s, t, c)$ est une réunion disjointe
de schémas en groupoïdes $(U_r, R_r, s, t, c)$ tels que chaque $s, t : R_r \to U_r$
soit de rang constant $r$. Comme $U$ est quasi-compact, on a $U_r = \emptyset$ pour
presque tout $r$. Il suffit de démontrer le lemme pour chaque $(U_r, R_r, s, t, c)$ ;
nous pouvons donc supposer que $s, t$ sont finis localement libres de rang $r$.

\medskip\noindent
Supposons que $s, t$ soient finis localement libres de rang $r$.
Soit $f \in A$. Considérons l'élément $x - f \in A[x]$, où nous regardons
$x$ comme la coordonnée sur $\mathbf{A}^1$.
Puisque
$$
(U \times \mathbf{A}^1, R \times \mathbf{A}^1,
s \times \text{id}_{\mathbf{A}^1},
t \times \text{id}_{\mathbf{A}^1},
c \times \text{id}_{\mathbf{A}^1})
$$
est lui aussi un schéma en groupoïdes dont la source et le but sont finis, nous pouvons lui appliquer
le Lemme \ref{lemma-determinant-trick} et nous voyons que
$P(x) = \text{Norm}_{s^\sharp}(t^\sharp(x - f))$
est un élément de $C[x]$. Puisque $s^\sharp : A \to B$ est fini localement
libre de rang $r$, on voit que $P$ est unitaire de degré $r$.
De plus, $P(f) = 0$ d'après Cayley-Hamilton
(Algèbre, Lemme \ref{algebra-lemma-charpoly}).
Plus précisément, Cayley-Hamilton affirme que le polynôme
$s^\sharp(P) \in B[x]$ annule l'élément $t^\sharp(f)$.
Puisque les coefficients de $P$ appartiennent à $C$, nous obtenons
$t^\sharp(P(f)) = 0$ dans $B$. Comme $t$ est fidèlement plat,
nous obtenons $P(f) = 0$.
\end{proof}

\begin{lemma}
\label{lemma-invariants-base-change}
Soit $S$ un schéma. Soit $(U, R, s, t, c)$ un schéma en groupoïdes sur $S$.
Supposons que $U = \Spec(A)$ et $R = \Spec(B)$ soient affines et que
$s, t : R \to U$ soient finis localement libres. Soit $C \subset A$ comme dans
(\ref{equation-invariants}). Soit $C \to C'$ un homomorphisme d'anneaux, et posons
$U' = \Spec(A \otimes_C C')$,
$R' = \Spec(B \otimes_C C')$.
Alors
\begin{enumerate}
\item Les applications $s, t, c$ donnent des applications $s', t', c'$.
Celles-ci font de $(U', R', s', t', c')$ un schéma en groupoïdes. Soit $C^1 \subset A'$
l'anneau des fonctions $R'$-invariantes sur $U'$.
\item L'application canonique $\varphi : C' \to C^1$ satisfait aux propriétés suivantes :
\begin{enumerate}
\item pour tout $f \in C^1$, il existe un $n > 0$ et un
polynôme $P \in C'[x]$ dont l'image dans $C^1[x]$ est
$(x - f)^n$, et
\item pour tout $f \in \Ker(\varphi)$, il existe
un $n > 0$ tel que $f^n = 0$.
\end{enumerate}
\item Si $C \to C'$ est plat, alors $\varphi$ est un isomorphisme.
\end{enumerate}
\end{lemma}

\begin{proof}
La preuve de la partie (1) est omise. Notons $A' = A \otimes_C C'$ et
$B' = B \otimes_C C'$. Alors nous avons
$$
C^1
= \{a \in A' \mid (t')^\sharp(a) = (s')^\sharp(a) \}
= \{a \in A \otimes_C C' \mid t^\sharp \otimes 1(a) = s^\sharp \otimes 1(a) \}.
$$
Autrement dit, $C^1$ est le noyau de l'application différence
$(t^\sharp - s^\sharp) \otimes 1$, qui est précisément le changement de base
de l'application $C$-linéaire $t^\sharp - s^\sharp : A \to B$ par $C \to C'$.
Ainsi, (3) en résulte.

\medskip\noindent
Preuve de la partie (2)(b). Puisque $C \to A$ est entier
(Lemme \ref{lemma-integral-over-invariants}) et injectif, on voit que
$\Spec(A) \to \Spec(C)$ est surjectif, voir
Algèbre, Lemme \ref{algebra-lemma-integral-overring-surjective}.
Ainsi, $\Spec(A') \to \Spec(C')$ est également surjectif
comme changement de base d'un morphisme surjectif
(Morphismes, Lemme \ref{morphisms-lemma-base-change-surjective}).
Par conséquent, $\Spec(C^1) \to \Spec(C')$ est lui aussi surjectif.
Cela implique (2)(b), par exemple d'après
Algèbre, Lemme \ref{algebra-lemma-image-dense-generic-points}.

\medskip\noindent
Preuve de la partie (2)(a). D'après le Lemme \ref{lemma-finite-locally-free-disjoint-free},
notre schéma en groupoïdes $(U, R, s, t, c)$ se décompose en une réunion disjointe finie
de schémas en groupoïdes $(U_r, R_r, s, t, c)$ tels que $s, t : R_r \to U_r$
soient finis localement libres de rang $r$. En tirant en arrière par $U' = \Spec(C') \to U$,
nous obtenons une décomposition analogue de $U'$ et de $U^1 = \Spec(C^1)$.
Nous montrerons dans le paragraphe suivant que (2)(a) vaut pour le système correspondant
d'anneaux $A_r, B_r, C_r, C'_r, C^1_r$, avec $n = r$.
Étant donné $f \in C^1$, soit alors $P_r \in C_r[x]$ le polynôme
dont l'image dans $C^1_r[x]$ est l'image de $(x - f)^r$.
En choisissant un entier $n$ suffisamment divisible, on voit qu'il
existe un polynôme $P \in C'[x]$ dont l'image dans $C^1[x]$ est
$(x - f)^n$ ; à savoir, nous prenons pour $P$ l'unique élément de
$C'[x]$ dont l'image dans $C'_r[x]$ est $P_r^{n/r}$.

\medskip\noindent
Dans ce paragraphe, nous démontrons (2)(a) lorsque les homomorphismes d'anneaux
$s^\sharp, t^\sharp : A \to B$ sont finis localement libres d'un rang fixé $r$.
Soit $f \in C^1 \subset A' = A \otimes_C C'$. Choisissons une
$C$-algèbre plate $D$ et une surjection $D \to C'$. Choisissons un relèvement
$g \in A \otimes_C D$ de $f$.
Considérons le polynôme
$$
P = \text{Norm}_{s^\sharp \otimes 1}(t^\sharp \otimes 1(x - g))
$$
dans $(A \otimes_C D)[x]$. D'après le Lemme \ref{lemma-determinant-trick}
et la partie (3) du présent lemme, les coefficients de $P$ appartiennent à $D$
(comparer avec la preuve du Lemme \ref{lemma-integral-over-invariants}).
D'autre part, l'image de $P$ dans $(A \otimes_C C')[x]$ est
$(x - f)^r$, car $t^\sharp \otimes 1(x - f) = s^\sharp(x - f)$
et $s^\sharp$ est fini localement libre de rang $r$.
Cela démontre l'assertion avec $P$ comme dans l'énoncé (2)(a),
donné par l'image de notre $P$ par l'application $D[x] \to C'[x]$.
\end{proof}

\begin{lemma}
\label{lemma-points}
Soit $S$ un schéma. Soit $(U, R, s, t, c)$ un schéma en groupoïdes sur $S$.
Supposons que $U = \Spec(A)$ et $R = \Spec(B)$ soient affines et que
$s, t : R \to U$ soient finis localement libres. Soit $C \subset A$ comme dans
(\ref{equation-invariants}). Alors $U \to M = \Spec(C)$ possède
les propriétés suivantes :
\begin{enumerate}
\item l'application sur les points $|U| \to |M|$ est surjective et
$u_0, u_1 \in |U|$ ont la même image si et seulement s'il
existe un $r \in |R|$ tel que $t(r) = u_0$ et $s(r) = u_1$, ce qui donne
la formule
$$
|M| = |U|/|R|
$$
\item pour tout corps algébriquement clos $k$, on a
$$
M(k) = U(k)/R(k)
$$
\end{enumerate}
\end{lemma}

\begin{proof}
Puisque $C \to A$ est entier (Lemme \ref{lemma-integral-over-invariants})
et injectif, on voit que $\Spec(A) \to \Spec(C)$ est surjectif, voir
Algèbre, Lemme \ref{algebra-lemma-integral-overring-surjective}.
Ainsi, $|U| \to |M|$ est surjectif.

\medskip\noindent
Soit $k$ un corps algébriquement clos et soit $C \to k$ un homomorphisme d'anneaux.
Puisque les morphismes surjectifs sont préservés par changement de base
(Morphismes, Lemme \ref{morphisms-lemma-base-change-surjective}),
on voit que $A \otimes_C k$ n'est pas nul. Or $k \subset A \otimes_C k$ est une
extension entière non nulle. Ainsi, tout corps résiduel de $A \otimes_C k$
est une extension algébrique de $k$, donc est égal à $k$. Nous voyons donc que
$U(k) \to M(k)$ est surjectif.

\medskip\noindent
Soient $a_0, a_1 : A \to k$ deux homomorphismes d'anneaux. S'il existe un homomorphisme d'anneaux
$b : B \to k$ telle que $a_0 = b \circ t^\sharp$ et $a_1 = b \circ s^\sharp$,
alors on voit par définition que $a_0|_C = a_1|_C$. Ainsi, l'application
$U(k) \to M(k)$ égalise les deux applications $R(k) \to U(k)$.
Réciproquement, supposons que $a_0|_C = a_1|_C$. Notons cette application
d'algèbres $c : C \to k$. Considérons le diagramme
$$
\xymatrix{
& &
B \ar@{-->}[lld] \\
k & &
A
\ar@<0.5ex>[ll]^{a_0}
\ar@<-0.5ex>[ll]_{a_1}
\ar@<1ex>[u]
\ar@<-1ex>[u] \\
& &
C \ar[u] \ar[llu]^c
}
$$
Si nous pouvons construire une flèche pointillée rendant le diagramme commutatif, alors
la preuve de la partie (2) du lemme est achevée. Puisque $s : A \to B$ est fini,
il existe un nombre fini d'homomorphismes d'anneaux
$b_1, \ldots, b_n : B \to k$ telles que $b_i \circ s^\sharp = a_1$.
Si la flèche pointillée n'existe pas, nous voyons qu'aucune des
$a'_i = b_i \circ t^\sharp$, $i = 1, \ldots, n$, n'est égale à $a_0$.
Par conséquent, les idéaux maximaux
$$
\mathfrak m'_i = \Ker(a_i' \otimes 1 : A \otimes_C k \to k)
$$
de $A \otimes_C k$ sont distincts de
$\mathfrak m = \Ker(a_0 \otimes 1 : A \otimes_C k \to k)$.
D'après Algèbre, Lemme \ref{algebra-lemma-silly}, nous obtiendrions un élément
$f \in A \otimes_C k$ tel que $f \in \mathfrak m$, mais que
$f \not \in \mathfrak m_i'$ pour $i = 1, \ldots, n$.
Considérons la norme
$$
g = \text{Norm}_{s^\sharp \otimes 1}(t^\sharp \otimes 1(f))
\in
A \otimes_C k
$$
D'après le Lemme \ref{lemma-determinant-trick}, celle-ci appartient à l'anneau des invariants
$C^1 \subset A \otimes_C k$ du schéma en groupoïdes obtenu par
changement de base (au moyen de l'application $c : C \to k$). D'une part,
$a_1(g) \in k^*$, puisque
la valeur de $t^\sharp(f)$ en tous les points (qui correspondent à
$b_1, \ldots, b_n$) situés au-dessus de $a_1$ est
inversible (insérer ici une référence ultérieure à cette propriété du déterminant).
D'autre part, puisque $f \in \mathfrak m$, nous voyons que
$f$ n'est pas une unité, et donc que $t^\sharp(f)$ n'est pas une unité
(car $t^\sharp \otimes 1$ est fidèlement plat),
et donc que sa norme n'est pas une unité (insérer ici une référence ultérieure
à cette propriété du déterminant). Nous concluons que $C^1$ contient
un élément qui n'est ni nilpotent
ni une unité. Nous allons maintenant montrer que cela conduit à une contradiction.
En effet, appliquons le Lemme \ref{lemma-invariants-base-change}
à l'application $c : C \to C' = k$ ; alors
nous voyons que l'application de $k$ dans l'anneau des invariants $C^1$ est injective
et que, de plus, pour tout élément $x \in C^1$, il existe un entier
$n > 0$ tel que $x^n \in k$. Ainsi, tout élément de $C^1$ est
soit une unité, soit nilpotent.

\medskip\noindent
Il nous reste à achever la preuve de (1). Nous savons déjà que
$|U| \to |M|$ est surjectif. Il est clair que $|U| \to |M|$ est
invariante par $|R|$. Enfin, supposons que $u_0, u_1 \in U$ aient la même
image $m \in M$. Alors les extensions de corps induites $\kappa(u_0)/\kappa(m)$
et $\kappa(u_1)/\kappa(m)$ sont algébriques (car $A$ est entier sur $C$,
comme utilisé ci-dessus). Par conséquent, si $k$ est une clôture algébrique de $\kappa(m)$,
nous pouvons trouver des $\kappa(m)$-plongements $\overline{u}_0 : \kappa(u_0) \to k$
et $\overline{u}_1 : \kappa(u_1) \to k$. Ceux-ci déterminent des points à valeurs dans $k$,
$\overline{u}_0, \overline{u}_1 \in U(k)$, ayant la même
image dans $M(k)$. D'après la partie (2), il existe un point
$\overline{r} \in R(k)$ tel que $s(\overline{r}) = \overline{u}_0$ et
$t(\overline{r}) = \overline{u}_1$. L'image $r \in R$ de $\overline{r}$
est un point tel que $s(r) = u_0$ et $t(r) = u_1$, comme souhaité.
\end{proof}

\begin{lemma}
\label{lemma-etale}
Soit $S$ un schéma. Soit $f : (U', R', s', t') \to (U, R, s, t, c)$ un
morphisme de schémas en groupoïdes sur $S$. Supposons que
\begin{enumerate}
\item $U$, $R$, $U'$, $R'$ soient affines,
\item $s, t, s', t'$ soient finis localement libres,
\item les diagrammes
$$
\xymatrix{
R' \ar[d]_{s'} \ar[r]_f & R \ar[d]^s \\
U' \ar[r]^f & U
}
\quad
\quad
\xymatrix{
R' \ar[d]_{t'} \ar[r]_f & R \ar[d]^t \\
U' \ar[r]^f & U
}
\quad
\quad
\xymatrix{
G' \ar[d] \ar[r]_f & G \ar[d] \\
U' \ar[r]^f & U
}
$$
soient cartésiens, où $G$ et $G'$ sont les schémas en groupes stabilisateurs, et
\item $f : U' \to U$ soit \'etale.
\end{enumerate}
Alors l'application $C \to C'$ de l'anneau des fonctions $R$-invariantes sur $U$
vers l'anneau des fonctions $R'$-invariantes sur $U'$ est \'etale et
$U' = \Spec(C') \times_{\Spec(C)} U$.
\end{lemma}

\begin{proof}
Posons $M = \Spec(C)$ et $M' = \Spec(C')$.
Écrivons $U = \Spec(A)$, $U' = \Spec(A')$, $R = \Spec(B)$ et
$R' = \Spec(B')$. Nous utiliserons les résultats des
Lemmes \ref{lemma-integral-over-invariants},
\ref{lemma-invariants-base-change} et
\ref{lemma-points}
sans autre mention.

\medskip\noindent
Supposons que $C$ soit un anneau local strictement hensélien. Soit $p \in M$
le point fermé et soit $p' \in M'$ un point au-dessus de $p$.
Affirmation : dans ce cas, il existe une décomposition en réunion disjointe
$(U', R', s', t', c') = (U, R, s, t, c) \amalg (U'', R'', s'', t'', c'')$
au-dessus de $(U, R, s, t, c)$ telle que, pour la décomposition en
réunion disjointe correspondante $M' = M \amalg M''$ au-dessus de $M$,
le point $p'$ corresponde à $p \in M$.

\medskip\noindent
L'affirmation implique le lemme. Supposons que $M_1 \to M$ soit un morphisme plat
de schémas affines. Nous pouvons alors tout changer de base vers $M_1$
sans modifier les hypothèses (1) -- (4).
Le Lemme \ref{lemma-invariants-base-change} montre que
$M_1$, resp.\ $M_1'$, est le spectre de l'anneau des
fonctions $R_1$-invariantes sur $U_1$,
resp.\ des fonctions $R'_1$-invariantes sur $U'_1$.
Supposons que $p' \in M'$ ait pour image $p \in M$.
Soit $M_1$ le spectre de l'hensélisé strict de
$\mathcal{O}_{M, p}$, de point fermé $p_1 \in M_1$.
Choisissons un point $p'_1 \in M'_1$ au-dessus de $p_1$ et de $p'$.
L'affirmation donne
$$
(U'_1, R'_1, s'_1, t'_1, c'_1) =
(U_1, R_1, s_1, t_1, c_1) \amalg
(U''_1, R''_1, s''_1, t''_1, c''_1)
$$
et, de manière correspondante, $M'_1 = M_1 \amalg M''_1$ comme schéma au-dessus de $M_1$.
Écrivons $M_1 = \Spec(C_1)$ et $C_1 = \colim C_i$ comme une limite inductive filtrante
de $C$-algèbres \'etales. Posons $M_i = \Spec(C_i)$.
Alors $M_1 = \lim M_i$, et de même pour les autres schémas.
D'après Limites, Lemmes \ref{limits-lemma-descend-opens} et
\ref{limits-lemma-descend-isomorphism},
nous pouvons trouver un $i$ tel que
$$
(U'_i, R'_i, s'_i, t'_i, c'_i) =
(U_i, R_i, s_i, t_i, c_i) \amalg
(U''_i, R''_i, s''_i, t''_i, c''_i)
$$
Nous concluons que $M'_i = M_i \amalg M''_i$. En particulier,
$M' \to M$ devient \'etale en un point au-dessus de $p'$ après un
changement de base \'etale. Cela implique que $M' \to M$ est \'etale en $p'$
(par exemple d'après Morphismes, Lemme
\ref{morphisms-lemma-set-points-where-fibres-etale}).
Nous démontrerons $U' \cong M' \times_M U$ après avoir prouvé l'affirmation.

\medskip\noindent
Preuve de l'affirmation. Observons que $U_p$ et $U'_{p'}$ ont un nombre fini de points.
Pour $u \in U_p$, l'extension $\kappa(u)/\kappa(p)$ est algébrique ;
ainsi, $\kappa(u)$ est séparablement clos.
Comme $U' \to U$ est \'etale, nous concluons que le morphisme $U'_{p'} \to U_p$
induit des isomorphismes de corps résiduels.
Soit $u' \in U'_{p'}$ d'image $u \in U_p$.
Par l'hypothèse (3), les morphismes de fibres schématiques
$(s')^{-1}(u') \to s^{-1}(u)$,
$(t')^{-1}(u') \to t^{-1}(u)$ et
$G'_{u'} \to G_u$ sont des isomorphismes. Comme $U_p = t(s^{-1}(u))$
(ensemblistement), nous concluons que l'application de l'ensemble des points de $U'_{p'}$
dans l'ensemble des points de $U_p$ est surjective.
Supposons que $u'_1$ et $u'_2$ soient des points de $U'_{p'}$ ayant
la même image $u$ dans $U_p$. Il existe alors un point
$r' \in R'_{p'}$ tel que $s'(r') = u'_1$ et $t'(r') = u'_2$.
Considérons les deux tours de corps
$$
\kappa(r')/\kappa(u'_1)/\kappa(u)/\kappa(p) \quad
\kappa(r')/\kappa(u'_2)/\kappa(u)/\kappa(p)
$$
dont les ``extrémités'' sont les mêmes que les deux ``extrémités'' des deux tours
$$
\kappa(r')/\kappa(u'_1)/\kappa(p')/\kappa(p) \quad
\kappa(r')/\kappa(u'_2)/\kappa(p')/\kappa(p)
$$
{\emergencystretch=3em Ces deux tours induisent les mêmes applications $\kappa(p') \to \kappa(r')$, puisque
$(U'_{p'}, R'_{p'}, s', t', c')$ est un schéma en groupoïdes sur $p'$.
Puisque $\kappa(u)/\kappa(p)$ est purement inséparable,
nous concluons que les deux applications induites
$\kappa(u) \to \kappa(r')$ sont les mêmes.
Par conséquent, $r'$ s'envoie sur un point de la fibre $G_u$.
Par l'hypothèse (3), nous concluons que $r' \in (G')_{u'_1}$.
En effet, nous pouvons considérer $G$ comme un sous-schéma fermé de $R$,
vu comme un schéma sur $U$ au moyen de $s$, et utiliser le fait que
son changement de base à $U'$ donne $G' \subset R'$.
En particulier, nous avons $u'_1 = u'_2$.
Nous concluons que $U'_{p'} \to U_p$ est une application bijective
sur les points et qu'elle induit des isomorphismes sur les corps résiduels.
Il s'ensuit que $U'_{p'}$ est un ensemble fini de points fermés
(Algèbre, Lemme \ref{algebra-lemma-finite-residue-extension-closed})
et donc que $U'_{p'}$ est fermé dans $U'$.
Soit $J' \subset A'$ l'idéal radical définissant $U'_{p'}$
ensemblistement.\par}

\medskip\noindent
Seconde partie de la preuve de l'affirmation.
Soit $\mathfrak m \subset C$ l'idéal maximal.
Observons que $(A, \mathfrak m A)$ est un couple hensélien d'après
Compléments d'algèbre, Lemme \ref{more-algebra-lemma-integral-over-henselian-pair}.
Soit $J = \sqrt{\mathfrak m A}$.
Alors $(A, J)$ est un couple hensélien
(Compléments d'algèbre, Lemme \ref{more-algebra-lemma-change-ideal-henselian-pair})
et l'homomorphisme d'anneaux \'etale
$A \to A'$ induit un isomorphisme $A/J \to A'/J'$
d'après nos considérations ci-dessus.
Nous concluons que $A' = A \times A''$ d'après
Compléments d'algèbre, Lemme \ref{more-algebra-lemma-characterize-henselian-pair}.
Considérons la décomposition en réunion disjointe
correspondante $U' = U \amalg U''$. L'ouvert $(s')^{-1}(U)$ est l'ensemble
des points de $R'$ qui se spécialisent en un point de $R'_{p'}$.
Il en va de même pour $(t')^{-1}(U)$. De même, nous avons
$(s')^{-1}(U'') = (t')^{-1}(U'')$, car il s'agit de l'ensemble des
points qui ne se spécialisent pas dans $R'_{p'}$.
Nous obtenons donc une décomposition en réunion disjointe
$$
(U', R', s', t', c') =
(U, R, s, t, c) \amalg
(U'', R'', s'', t'', c'')
$$
Celle-ci donne immédiatement $M' = M \amalg M''$, et la preuve de l'affirmation
est achevée.

\medskip\noindent
Il nous reste à montrer que l'application canonique $U' \to M' \times_M U$
est un isomorphisme. C'est un morphisme \'etale
(Morphismes, Lemme \ref{morphisms-lemma-etale-permanence}).
D'autre part, en effectuant des changements de base vers des anneaux locaux strictement henséliens
(comme dans le troisième paragraphe de la preuve) et en utilisant la bijectivité
$U'_{p'} \to U_p$ établie au cours de la preuve de l'affirmation,
nous voyons que $U' \to M' \times_M U$ est universellement bijectif
(certains détails sont omis). Or un morphisme universellement bijectif
et \'etale est un isomorphisme
(Descente, Lemme \ref{descent-lemma-universally-injective-etale-open-immersion}),
ce qui achève la preuve.
\end{proof}

\begin{lemma}
\label{lemma-basis}
Soit $S$ un schéma.
Soit $(U, R, s, t, c)$ un schéma en groupoïdes sur $S$.
Supposons que
\begin{enumerate}
\item $U = \Spec(A)$ et $R = \Spec(B)$ soient affines, et qu'
\item il existe des éléments $x_i \in A$, $i \in I$, tels que
$B = \bigoplus_{i \in I} s^\sharp(A)t^\sharp(x_i)$.
\end{enumerate}
Alors $A = \bigoplus_{i\in I} Cx_i$ et $B \cong A \otimes_C A$,
où $C \subset A$ est l'anneau des fonctions $R$-invariantes
sur $U$, comme dans (\ref{equation-invariants}).
\end{lemma}

\begin{proof}
Dans cette preuve, nous écrirons $s, t : A \to B$ au lieu de
$s^\sharp, t^\sharp$, et de même $c : B \to B \otimes_{s, A, t} B$.
Nous écrivons $p_0 : B \to B \otimes_{s, A, t} B$, $b \mapsto b \otimes 1$ et
$p_1 : B \to B \otimes_{s, A, t} B$, $b \mapsto 1 \otimes b$. D'après le
Lemme \ref{lemma-diagram-pull}
et la définition de $C$, nous avons le diagramme
commutatif suivant :
$$
\xymatrix{
B \otimes_{s, A, t} B &
B \ar@<-1ex>[l]_-c \ar@<1ex>[l]^-{p_0} &
A \ar[l]^t \\
B \ar[u]^{p_1} &
A \ar@<-1ex>[l]_s \ar@<1ex>[l]^t \ar[u]_s &
C \ar[u] \ar[l]
}
$$
De plus, les deux carrés de gauche sont cocartésiens dans la catégorie des anneaux, et
la ligne supérieure est isomorphe au diagramme
$$
\xymatrix{
B \otimes_{t, A, t} B &
B \ar@<-1ex>[l]_-{p_1} \ar@<1ex>[l]^-{p_0} &
A \ar[l]^t
}
$$
qui est un diagramme égalisateur d'après
Descente, Lemme \ref{descent-lemma-ff-exact}, car la condition (2) implique
en particulier que $s$ (et donc aussi la flèche $t$ qui lui est isomorphe)
est fidèlement plat.
La ligne inférieure est un diagramme égalisateur par définition de $C$.
Nous pouvons utiliser les $x_i$ et obtenir un diagramme commutatif
$$
\xymatrix{
B \otimes_{s, A, t} B &
B \ar@<-1ex>[l]_-c \ar@<1ex>[l]^-{p_0} &
A \ar[l]^t \\
\bigoplus_{i \in I} B x_i \ar[u]^{p_1} &
\bigoplus_{i \in I} A x_i \ar@<-1ex>[l]_s \ar@<1ex>[l]^t \ar[u]_s &
\bigoplus_{i \in I} C x_i \ar[u] \ar[l]
}
$$
où la flèche verticale de droite envoie $x_i$ sur $x_i$,
la flèche verticale du milieu envoie $x_i$ sur $t(x_i)$ et
la flèche verticale de gauche envoie $x_i$ sur
$c(t(x_i)) = t(x_i) \otimes 1 = p_0(t(x_i))$ (égalité résultant de la commutativité
de la partie supérieure du diagramme du lemme \ref{lemma-diagram}). Le diagramme
est alors commutatif. De plus, la flèche verticale du milieu est un isomorphisme
par hypothèse. Les deux carrés de gauche étant cocartésiens, on
en conclut que la flèche verticale de gauche est elle aussi un isomorphisme.
D'autre part, les lignes horizontales sont exactes (c'est-à-dire que ce sont des
égalisateurs). On en conclut donc que la flèche verticale de droite
est elle aussi un isomorphisme.
\end{proof}

\begin{proposition}
\label{proposition-finite-flat-equivalence}
Soit $S$ un schéma.
Soit $(U, R, s, t, c)$ un schéma en groupoïdes sur $S$.
Supposons que
\begin{enumerate}
\item $U = \Spec(A)$ et $R = \Spec(B)$ soient affines,
\item $s, t : R \to U$ soient finis localement libres, et
\item $j = (t, s)$ soit une relation d'équivalence.
\end{enumerate}
Dans ce cas, soit $C \subset A$ comme dans
(\ref{equation-invariants}). Alors $U \to M = \Spec(C)$
est fini localement libre et $R = U \times_M U$.
De plus, $M$ représente le faisceau quotient $U/R$
pour la topologie fppf (voir la définition \ref{definition-quotient-sheaf}).
\end{proposition}

\begin{proof}
Dans cette démonstration, on utilise la notation $s, t : A \to B$
au lieu de la notation $s^\sharp, t^\sharp$.
D'après le lemme \ref{lemma-criterion-quotient-representable},
il suffit de montrer que $C \to A$ est fini localement libre
et que l'application
$$
t \otimes s : A \otimes_C A \longrightarrow B
$$
est un isomorphisme. Notons d'abord que $j$ est un monomorphisme, et
qu'il est aussi fini (puisque $s$ et $t$ sont déjà finis). On voit donc
que $j$ est une immersion fermée d'après
Morphismes, lemme \ref{morphisms-lemma-finite-monomorphism-closed}.
Ainsi, $A \otimes_C A \to B$ est surjective.

\medskip\noindent
On effectuera des changements de base par des homomorphismes plats d'anneaux $C \to C'$, comme dans
le lemme \ref{lemma-invariants-base-change}, et l'on utilisera le fait que
la formation des invariants commute au changement de base plat, voir
la partie (3) du lemme cité.
On montrera ci-dessous que, pour tout idéal premier $\mathfrak p \subset C$, il existe
un homomorphisme local plat d'anneaux $C_{\mathfrak p} \to C_{\mathfrak p}'$
tel que le résultat soit vrai après changement de base à $C_{\mathfrak p}'$.
Cela implique immédiatement
que $A \otimes_C A \to B$ est injective (utiliser
Algèbre, lemme \ref{algebra-lemma-characterize-zero-local}).
Cela implique aussi que $C \to A$ est plat, en combinant
Algèbre, lemmes \ref{algebra-lemma-local-flat-ff},
\ref{algebra-lemma-flat-localization} et
\ref{algebra-lemma-flatness-descends}. Puisque $U \to \Spec(C)$
est également surjectif (lemme \ref{lemma-points}), on en conclut que $C \to A$
est fidèlement plat. L'isomorphisme $B \cong A \otimes_C A$
implique alors que $A$ est un $C$-module de présentation finie, voir
Algèbre, lemme \ref{algebra-lemma-descend-properties-modules}.
Ainsi, $A$ est fini localement libre sur $C$, voir
Algèbre, lemme \ref{algebra-lemma-finite-projective}.

\medskip\noindent
D'après le lemme \ref{lemma-finite-locally-free-disjoint-free},
on sait que $A$ est un produit fini
d'anneaux $A_r$ et que $B$ est un produit fini d'anneaux $B_r$,
de sorte que le schéma en groupoïdes se décompose de façon correspondante (voir la démonstration
du lemme \ref{lemma-integral-over-invariants}).
Alors $C$ est lui aussi un produit d'anneaux $C_r$, et
$C'$ se décompose de façon correspondante en un produit. On peut donc, et l'on va,
supposer que les homomorphismes d'anneaux $s, t : A \to B$ sont finis
localement libres d'un rang fixé $r$.

\medskip\noindent
Les homomorphismes locaux d'anneaux $C_{\mathfrak p} \to C_{\mathfrak p}'$ que l'on va
utiliser sont des morphismes locaux plats quelconques tels que le corps résiduel de
$C_{\mathfrak p}'$ soit infini.
D'après Algèbre, lemme \ref{algebra-lemma-flat-local-given-residue-field},
de tels morphismes locaux existent.

\medskip\noindent
Supposons que $C$ soit un anneau local d'idéal maximal $\mathfrak m$ et
de corps résiduel infini, et que $s, t : A \to B$ soient
finis localement libres de rang constant $r > 0$.
Comme l'extension $C \subset A$ est entière (lemme \ref{lemma-integral-over-invariants}),
tous les idéaux premiers au-dessus de $\mathfrak m$ sont maximaux, et tous les idéaux maximaux
de $A$ sont au-dessus de $\mathfrak m$. Il en va de même pour $C \subset B$.
Choisissons un idéal maximal $\mathfrak m'$
de $A$ au-dessus de $\mathfrak m$ (il en existe d'après le lemme \ref{lemma-points}).
Comme $t : A \to B$ est fini localement libre, il n'existe qu'un nombre fini
d'idéaux maximaux de $B$ au-dessus de $\mathfrak m'$. On en conclut
(de nouveau par le lemme \ref{lemma-points})
que $A$ a un nombre fini d'idéaux maximaux, c'est-à-dire que
$A$ est semi-local. Cela implique à son tour que $B$ est également
semi-local. Or, puisque $t \otimes s : A \otimes_C A \to B$ est surjective,
on peut appliquer
Algèbre, lemme \ref{algebra-lemma-semi-local-module-basis-in-submodule}
à l'homomorphisme d'anneaux $C \to A$, au $A$-module $M = B$ (considéré comme un $A$-module
via $t$) et au $C$-sous-module $s(A) \subset B$. Ce lemme implique qu'il
existe $x_1, \ldots, x_r \in A$ tels que $M$ soit libre sur $A$
de base $s(x_1), \ldots, s(x_r)$. On en conclut que $C \to A$ est fini libre
et que $B \cong A \otimes_C A$, en appliquant
le lemme \ref{lemma-basis}.
\end{proof}



\section{Groupoïdes finis plats}
\label{section-finite-flat-general}

\noindent
Dans cette section, on démontre un lemme qui aidera à montrer que le quotient
d'un schéma par une relation d'équivalence finie plate est un schéma, pourvu que
chaque classe d'équivalence soit contenue dans un ouvert affine. Voir
Propriétés des espaces,
proposition \ref{spaces-properties-proposition-finite-flat-equivalence-global}.

\begin{lemma}
\label{lemma-find-invariant-affine}
Soit $S$ un schéma.
Soit $(U, R, s, t, c)$ un schéma en groupoïdes sur $S$.
Supposons que $s$, $t$ soient finis localement libres.
Soit $u \in U$ un point tel que $t(s^{-1}(\{u\}))$
soit contenu dans un ouvert affine de $U$.
Alors il existe un voisinage ouvert affine $R$-invariant
de $u$ dans $U$.
\end{lemma}

\begin{proof}
Puisque $s$ est fini localement libre, ses fibres sont finies. Ainsi,
$t(s^{-1}(\{u\})) = \{u_1, \ldots, u_n\}$ est un ensemble fini.
Notons que $u \in \{u_1, \ldots, u_n\}$.
Soit $W \subset U$ un ouvert affine contenant $\{u_1, \ldots, u_n\}$ ;
en particulier, $u \in W$. Considérons
$Z = R \setminus s^{-1}(W) \cap t^{-1}(W)$. C'est un fermé
de $R$. L'image $t(Z)$ est un fermé de $U$ que l'on peut décrire
informellement comme l'ensemble des points de $U$ qui sont $R$-équivalents à un point
de $U \setminus W$. Ainsi, $W' = U \setminus t(Z)$ est un sous-schéma ouvert $R$-invariant
de $U$ contenu dans $W$, et $\{u_1, \ldots, u_n\} \subset W'$.
La situation est la suivante :
$$
\{u_1, \ldots, u_n\} \subset W' \subset W \subset U.
$$
Soit $f \in \Gamma(W, \mathcal{O}_W)$ un élément tel que
$\{u_1, \ldots, u_n\} \subset D(f) \subset W'$. Un tel $f$ existe d'après
Algèbre, lemme \ref{algebra-lemma-silly}. Par notre choix de $W'$,
on a $s^{-1}(W') \subset t^{-1}(W)$, et l'on obtient donc un diagramme
$$
\xymatrix{
s^{-1}(W') \ar[d]_s \ar[r]_-t & W \\
W'
}
$$
La flèche verticale est finie localement libre par hypothèse. Posons
$$
g = \text{Norm}_s(t^\sharp f) \in \Gamma(W', \mathcal{O}_{W'})
$$
Par construction, $g$ est une fonction sur $W'$ qui est
non nulle en $u$, car $t^\sharp(f)$ est non nulle en chacun des points de
$R$ au-dessus de $u$, puisque $f$ est non nulle en $u_1, \ldots, u_n$.
De même, $D(g) \subset W'$ est égal à
l'ensemble des points $w$ tels que $f$ ne s'annule en aucun des points
équivalents à $w$. Cela signifie que $D(g)$ est un
ouvert affine $R$-invariant de $W'$. La situation finale est
$$
\{u_1, \ldots, u_n\} \subset D(g) \subset D(f) \subset W' \subset W \subset U
$$
et l'on obtient donc le résultat souhaité.
\end{proof}







\section{Descente des schémas quasi-projectifs}
\label{section-quasi-projective}

\noindent
On peut utiliser le lemme \ref{lemma-find-invariant-affine}
pour montrer qu'un certain type de donnée de descente est effectif.

\begin{lemma}
\label{lemma-descend-along-finite-with-bound}
Soit $X \to Y$ un morphisme surjectif fini localement libre.
Soit $d$ un entier strictement positif. Supposons que, pour tout point géométrique
$\bar{y} : \mathrm{Spec}(k) \to Y$, la fibre $X_{\bar{y}}$
ait au plus $d$ points.
Soit $V$ un schéma sur $X$ tel que, pour tout
$(y, v_1, \ldots, v_d)$ où $y \in Y$ et
$v_1, \ldots, v_d \in V_y$, il existe un ouvert affine
$U \subset V$ tel que $v_1, \ldots, v_d \in U$.
Alors toute donnée de descente sur $V/X/Y$ est effective.
\end{lemma}

\begin{proof}
Soit $\varphi$ une donnée de descente au sens de
Descente, définition \ref{descent-definition-descent-datum}.
Rappelons que le foncteur des schémas sur $Y$ vers les données de descente
relatives à $\{X \to Y\}$ est pleinement fidèle, voir
Descente, lemme \ref{descent-lemma-refine-coverings-fully-faithful}.
Ainsi, en utilisant Constructions, lemme \ref{constructions-lemma-relative-glueing},
il suffit de démontrer le lemme dans le cas où $Y$ est affine.
Certains détails sont omis (on peut éviter cet argument si $Y$ est
séparé ou possède une diagonale affine, car tout morphisme d'un
schéma affine vers $X$ est alors affine).

\medskip\noindent
Supposons que $Y$ soit affine. Si $V$ est lui aussi affine, on a l'effectivité d'après
Descente, lemme \ref{descent-lemma-affine}. Ainsi, d'après
Descente, lemme \ref{descent-lemma-effective-for-fpqc-is-local-upstairs},
il suffit de démontrer que tout point $v$ de $V$ possède un voisinage ouvert affine $\varphi$-invariant.
Considérons le groupoïde
$(X, X \times_Y X, \text{pr}_1, \text{pr}_0, \text{pr}_{02})$.
D'après le lemme \ref{lemma-cartesian-equivalent-descent-datum},
la donnée de descente $\varphi$ détermine, et est déterminée par,
un morphisme cartésien de schémas en groupoïdes
$$
(V, R, s, t, c)
\longrightarrow
(X, X \times_Y X, \text{pr}_1, \text{pr}_0, \text{pr}_{02})
$$
sur $\Spec(\mathbf{Z})$.
Puisque $X \to Y$ est fini localement libre, on voit que
$\text{pr}_i : X \times_Y X \to X$, et donc $s$ et $t$,
sont finis localement libres, et leurs fibres géométriques ont au plus $d$
points. En particulier, la $R$-orbite $t(s^{-1}(\{v\}))$ de notre point
$v \in V$ a au plus $d$ points. En utilisant encore une fois l'équivalence de catégories du
lemme \ref{lemma-cartesian-equivalent-descent-datum},
on voit que les ouverts $\varphi$-invariants de $V$
sont exactement les ouverts $R$-invariants de $V$.
Notre hypothèse montre qu'il existe un ouvert affine de $V$
contenant l'orbite $t(s^{-1}(\{v\}))$, car tous les points
de cette orbite ont la même image dans $Y$.
Le lemme \ref{lemma-find-invariant-affine}
fournit donc un ouvert affine $R$-invariant contenant $v$.
\end{proof}

\begin{lemma}
\label{lemma-descend-along-finite}
Soit $X \to Y$ un morphisme surjectif fini localement libre.
Soit $V$ un schéma sur $X$ tel que, pour tout
$(y, v_1, \ldots, v_d)$ où $y \in Y$ et
$v_1, \ldots, v_d \in V_y$, il existe un ouvert affine
$U \subset V$ tel que $v_1, \ldots, v_d \in U$.
Alors toute donnée de descente sur $V/X/Y$ est effective.
\end{lemma}

\begin{proof}
Comme dans la démonstration du lemme \ref{lemma-descend-along-finite-with-bound}, on
peut supposer que $Y$ est affine, donc en particulier quasi-compact. Il
existe alors un entier $d$ tel que toutes les fibres géométriques de $X \to Y$
aient au plus $d$ points, voir la discussion dans Morphismes, section
\ref{morphisms-section-universally-bounded}.
Notre hypothèse sur $V$ implique trivialement que l'on peut appliquer
le lemme \ref{lemma-descend-along-finite-with-bound}, ce qui termine la démonstration.
\end{proof}

\begin{lemma}
\label{lemma-descend-along-finite-quasi-projective}
Soit $X \to Y$ un morphisme surjectif fini localement libre.
Soit $V$ un schéma sur $X$ tel que l'une des conditions suivantes soit satisfaite :
\begin{enumerate}
\item $V \to X$ est projectif,
\item $V \to X$ est quasi-projectif,
\item il existe un faisceau inversible ample sur $V$,
\item il existe un faisceau inversible $X$-ample sur $V$,
\item il existe un faisceau inversible $X$-très ample sur $V$.
\end{enumerate}
Alors toute donnée de descente sur $V/X/Y$ est effective.
\end{lemma}

\begin{proof}
Vérifions la condition du lemme \ref{lemma-descend-along-finite}.
Soient $y \in Y$ et $v_1, \ldots, v_d \in V$ des points au-dessus de $y$.
Le cas (1) est un cas particulier de (2), voir
Morphismes, lemme \ref{morphisms-lemma-projective-quasi-projective}.
Le cas (2) est un cas particulier de (4), voir
Morphismes, définition \ref{morphisms-definition-quasi-projective}.
S'il existe un faisceau inversible ample sur $V$, alors
il existe un ouvert affine contenant $v_1, \ldots, v_d$ d'après
Propriétés, lemme \ref{properties-lemma-ample-finite-set-in-affine}.
Ainsi, la condition (3) est satisfaite.
Dans les cas (4) et (5), on peut sans inconvénient remplacer $Y$ par un
voisinage ouvert affine de $y$.
Alors $X$ est lui aussi affine.
Dans le cas (4), on voit que $V$ possède un faisceau inversible ample
d'après Morphismes, définition \ref{morphisms-definition-relatively-ample},
et le résultat découle du cas (3).
Dans le cas (5), on peut remplacer $V$ par un ouvert quasi-compact contenant
$v_1, \ldots, v_d$, et l'on se ramène au cas (4) par
Morphismes, lemme \ref{morphisms-lemma-ample-very-ample}.
\end{proof}

\begin{lemma}
\label{lemma-descend-along-finite-radicial}
Soit $X \to Y$ un morphisme surjectif fini localement libre qui est radiciel.
Soit $V$ un schéma sur $X$. Alors toute donnée de descente sur $V/X/Y$ est effective.
\end{lemma}

\begin{proof}
Appliquer le lemme \ref{lemma-descend-along-finite-with-bound} avec $d = 1$.
\end{proof}








\begin{multicols}{2}[\section{Autres chapitres}]
\noindent
Préliminaires
\begin{enumerate}
\item \hyperref[introduction-section-phantom]{Introduction}
\item \hyperref[conventions-section-phantom]{Conventions}
\item \hyperref[sets-section-phantom]{Théorie des ensembles}
\item \hyperref[categories-section-phantom]{Catégories}
\item \hyperref[topology-section-phantom]{Topologie}
\item \hyperref[sheaves-section-phantom]{Faisceaux sur les espaces}
\item \hyperref[sites-section-phantom]{Sites et faisceaux}
\item \hyperref[stacks-section-phantom]{Champs}
\item \hyperref[fields-section-phantom]{Corps}
\item \hyperref[algebra-section-phantom]{Algèbre commutative}
\item \hyperref[brauer-section-phantom]{Groupes de Brauer}
\item \hyperref[homology-section-phantom]{Algèbre homologique}
\item \hyperref[derived-section-phantom]{Catégories dérivées}
\item \hyperref[simplicial-section-phantom]{Méthodes simpliciales}
\item \hyperref[more-algebra-section-phantom]{Compléments d'algèbre}
\item \hyperref[smoothing-section-phantom]{Lissage des morphismes d'anneaux}
\item \hyperref[modules-section-phantom]{Faisceaux de modules}
\item \hyperref[sites-modules-section-phantom]{Modules sur les sites}
\item \hyperref[injectives-section-phantom]{Objets injectifs}
\item \hyperref[cohomology-section-phantom]{Cohomologie des faisceaux}
\item \hyperref[sites-cohomology-section-phantom]{Cohomologie sur les sites}
\item \hyperref[dga-section-phantom]{Algèbre différentielle graduée}
\item \hyperref[dpa-section-phantom]{Algèbre à puissances divisées}
\item \hyperref[sdga-section-phantom]{Faisceaux différentiels gradués}
\item \hyperref[hypercovering-section-phantom]{Hyperrecouvrements}
\end{enumerate}
Schémas
\begin{enumerate}
\setcounter{enumi}{25}
\item \hyperref[schemes-section-phantom]{Schémas}
\item \hyperref[constructions-section-phantom]{Constructions de schémas}
\item \hyperref[properties-section-phantom]{Propriétés des schémas}
\item \hyperref[morphisms-section-phantom]{Morphismes de schémas}
\item \hyperref[coherent-section-phantom]{Cohomologie des schémas}
\item \hyperref[divisors-section-phantom]{Diviseurs}
\item \hyperref[limits-section-phantom]{Limites de schémas}
\item \hyperref[varieties-section-phantom]{Variétés}
\item \hyperref[topologies-section-phantom]{Topologies sur les schémas}
\item \hyperref[descent-section-phantom]{Descente}
\item \hyperref[perfect-section-phantom]{Catégories dérivées de schémas}
\item \hyperref[more-morphisms-section-phantom]{Compléments sur les morphismes}
\item \hyperref[flat-section-phantom]{Compléments sur la platitude}
\item \hyperref[groupoids-section-phantom]{Schémas en groupoïdes}
\item \hyperref[more-groupoids-section-phantom]{Compléments sur les schémas en groupoïdes}
\item \hyperref[etale-section-phantom]{Morphismes étales de schémas}
\end{enumerate}
Thèmes de théorie des schémas
\begin{enumerate}
\setcounter{enumi}{41}
\item \hyperref[chow-section-phantom]{Homologie de Chow}
\item \hyperref[intersection-section-phantom]{Théorie de l'intersection}
\item \hyperref[pic-section-phantom]{Schémas de Picard des courbes}
\item \hyperref[weil-section-phantom]{Théories cohomologiques de Weil}
\item \hyperref[adequate-section-phantom]{Modules adéquats}
\item \hyperref[dualizing-section-phantom]{Complexes dualisants}
\item \hyperref[duality-section-phantom]{Dualité pour les schémas}
\item \hyperref[discriminant-section-phantom]{Discriminants et différentes}
\item \hyperref[derham-section-phantom]{Cohomologie de de Rham}
\item \hyperref[local-cohomology-section-phantom]{Cohomologie locale}
\item \hyperref[algebraization-section-phantom]{Géométrie algébrique et formelle}
\item \hyperref[curves-section-phantom]{Courbes algébriques}
\item \hyperref[resolve-section-phantom]{Résolution des surfaces}
\item \hyperref[models-section-phantom]{Réduction semi-stable}
\item \hyperref[functors-section-phantom]{Foncteurs et morphismes}
\item \hyperref[equiv-section-phantom]{Catégories dérivées de variétés}
\item \hyperref[pione-section-phantom]{Groupes fondamentaux de schémas}
\item \hyperref[etale-cohomology-section-phantom]{Cohomologie étale}
\item \hyperref[crystalline-section-phantom]{Cohomologie cristalline}
\item \hyperref[proetale-section-phantom]{Cohomologie pro-étale}
\item \hyperref[relative-cycles-section-phantom]{Cycles relatifs}
\item \hyperref[more-etale-section-phantom]{Compléments de cohomologie étale}
\item \hyperref[trace-section-phantom]{La formule des traces}
\end{enumerate}
Espaces algébriques
\begin{enumerate}
\setcounter{enumi}{64}
\item \hyperref[spaces-section-phantom]{Espaces algébriques}
\item \hyperref[spaces-properties-section-phantom]{Propriétés des espaces algébriques}
\item \hyperref[spaces-morphisms-section-phantom]{Morphismes d'espaces algébriques}
\item \hyperref[decent-spaces-section-phantom]{Espaces algébriques décents}
\item \hyperref[spaces-cohomology-section-phantom]{Cohomologie des espaces algébriques}
\item \hyperref[spaces-limits-section-phantom]{Limites d'espaces algébriques}
\item \hyperref[spaces-divisors-section-phantom]{Diviseurs sur les espaces algébriques}
\item \hyperref[spaces-over-fields-section-phantom]{Espaces algébriques sur des corps}
\item \hyperref[spaces-topologies-section-phantom]{Topologies sur les espaces algébriques}
\item \hyperref[spaces-descent-section-phantom]{Descente et espaces algébriques}
\item \hyperref[spaces-perfect-section-phantom]{Catégories dérivées d'espaces}
\item \hyperref[spaces-more-morphisms-section-phantom]{Compléments sur les morphismes d'espaces}
\item \hyperref[spaces-flat-section-phantom]{Platitude sur les espaces algébriques}
\item \hyperref[spaces-groupoids-section-phantom]{Groupoïdes dans les espaces algébriques}
\item \hyperref[spaces-more-groupoids-section-phantom]{Compléments sur les groupoïdes dans les espaces}
\item \hyperref[bootstrap-section-phantom]{Amorçage}
\item \hyperref[spaces-pushouts-section-phantom]{Sommes amalgamées d'espaces algébriques}
\end{enumerate}
Thèmes de géométrie
\begin{enumerate}
\setcounter{enumi}{81}
\item \hyperref[spaces-chow-section-phantom]{Groupes de Chow des espaces}
\item \hyperref[groupoids-quotients-section-phantom]{Quotients de groupoïdes}
\item \hyperref[spaces-more-cohomology-section-phantom]{Compléments sur la cohomologie des espaces}
\item \hyperref[spaces-simplicial-section-phantom]{Espaces simpliciaux}
\item \hyperref[spaces-duality-section-phantom]{Dualité pour les espaces}
\item \hyperref[formal-spaces-section-phantom]{Espaces algébriques formels}
\item \hyperref[restricted-section-phantom]{Algébrisation des espaces formels}
\item \hyperref[spaces-resolve-section-phantom]{Résolution des surfaces revisitée}
\end{enumerate}
Théorie des déformations
\begin{enumerate}
\setcounter{enumi}{89}
\item \hyperref[formal-defos-section-phantom]{Théorie formelle des déformations}
\item \hyperref[defos-section-phantom]{Théorie des déformations}
\item \hyperref[cotangent-section-phantom]{Le complexe cotangent}
\item \hyperref[examples-defos-section-phantom]{Problèmes de déformation}
\end{enumerate}
Champs algébriques
\begin{enumerate}
\setcounter{enumi}{93}
\item \hyperref[algebraic-section-phantom]{Champs algébriques}
\item \hyperref[examples-stacks-section-phantom]{Exemples de champs}
\item \hyperref[stacks-sheaves-section-phantom]{Faisceaux sur les champs algébriques}
\item \hyperref[criteria-section-phantom]{Critères de représentabilité}
\item \hyperref[artin-section-phantom]{Axiomes d'Artin}
\item \hyperref[quot-section-phantom]{Espaces de Quot et de Hilbert}
\item \hyperref[stacks-properties-section-phantom]{Propriétés des champs algébriques}
\item \hyperref[stacks-morphisms-section-phantom]{Morphismes de champs algébriques}
\item \hyperref[stacks-limits-section-phantom]{Limites de champs algébriques}
\item \hyperref[stacks-cohomology-section-phantom]{Cohomologie des champs algébriques}
\item \hyperref[stacks-perfect-section-phantom]{Catégories dérivées de champs}
\item \hyperref[stacks-introduction-section-phantom]{Introduction aux champs algébriques}
\item \hyperref[stacks-more-morphisms-section-phantom]{Compléments sur les morphismes de champs}
\item \hyperref[stacks-geometry-section-phantom]{La géométrie des champs}
\end{enumerate}
Thèmes de théorie des espaces de modules
\begin{enumerate}
\setcounter{enumi}{107}
\item \hyperref[moduli-section-phantom]{Champs de modules}
\item \hyperref[moduli-curves-section-phantom]{Espaces de modules de courbes}
\end{enumerate}
Divers
\begin{enumerate}
\setcounter{enumi}{109}
\item \hyperref[examples-section-phantom]{Exemples}
\item \hyperref[exercises-section-phantom]{Exercices}
\item \hyperref[guide-section-phantom]{Guide bibliographique}
\item \hyperref[desirables-section-phantom]{Desiderata}
\item \hyperref[coding-section-phantom]{Style de programmation}
\item \hyperref[obsolete-section-phantom]{Obsolète}
\item \hyperref[fdl-section-phantom]{Licence de documentation libre GNU}
\item \hyperref[index-section-phantom]{Index généré automatiquement}
\end{enumerate}
\end{multicols}


\bibliography{my}
\bibliographystyle{amsalpha}

\end{document}
